\PassOptionsToPackage{unicode=true}{hyperref} % options for packages loaded elsewhere
\PassOptionsToPackage{hyphens}{url}
%
\documentclass[]{book}
\usepackage{lmodern}
\usepackage{amssymb,amsmath}
\usepackage{ifxetex,ifluatex}
\usepackage{fixltx2e} % provides \textsubscript
\ifnum 0\ifxetex 1\fi\ifluatex 1\fi=0 % if pdftex
  \usepackage[T1]{fontenc}
  \usepackage[utf8]{inputenc}
  \usepackage{textcomp} % provides euro and other symbols
\else % if luatex or xelatex
  \usepackage{unicode-math}
  \defaultfontfeatures{Ligatures=TeX,Scale=MatchLowercase}
    \setmainfont[]{LiberationSerif}
    \setmonofont[Mapping=tex-ansi,Scale=0.7]{LiberationSerif}
\fi
% use upquote if available, for straight quotes in verbatim environments
\IfFileExists{upquote.sty}{\usepackage{upquote}}{}
% use microtype if available
\IfFileExists{microtype.sty}{%
\usepackage[]{microtype}
\UseMicrotypeSet[protrusion]{basicmath} % disable protrusion for tt fonts
}{}
\IfFileExists{parskip.sty}{%
\usepackage{parskip}
}{% else
\setlength{\parindent}{0pt}
\setlength{\parskip}{6pt plus 2pt minus 1pt}
}
\usepackage{hyperref}
\hypersetup{
            pdftitle={The Sacraments},
            pdfauthor={James Bannerman},
            pdfborder={0 0 0},
            breaklinks=true}
\urlstyle{same}  % don't use monospace font for urls
\usepackage{longtable,booktabs}
% Fix footnotes in tables (requires footnote package)
\IfFileExists{footnote.sty}{\usepackage{footnote}\makesavenoteenv{longtable}}{}
\usepackage{graphicx,grffile}
\makeatletter
\def\maxwidth{\ifdim\Gin@nat@width>\linewidth\linewidth\else\Gin@nat@width\fi}
\def\maxheight{\ifdim\Gin@nat@height>\textheight\textheight\else\Gin@nat@height\fi}
\makeatother
% Scale images if necessary, so that they will not overflow the page
% margins by default, and it is still possible to overwrite the defaults
% using explicit options in \includegraphics[width, height, ...]{}
\setkeys{Gin}{width=\maxwidth,height=\maxheight,keepaspectratio}
\setlength{\emergencystretch}{3em}  % prevent overfull lines
\providecommand{\tightlist}{%
  \setlength{\itemsep}{0pt}\setlength{\parskip}{0pt}}
\setcounter{secnumdepth}{5}
% Redefines (sub)paragraphs to behave more like sections
\ifx\paragraph\undefined\else
\let\oldparagraph\paragraph
\renewcommand{\paragraph}[1]{\oldparagraph{#1}\mbox{}}
\fi
\ifx\subparagraph\undefined\else
\let\oldsubparagraph\subparagraph
\renewcommand{\subparagraph}[1]{\oldsubparagraph{#1}\mbox{}}
\fi

% set default figure placement to htbp
\makeatletter
\def\fps@figure{htbp}
\makeatother

\usepackage{booktabs}
\usepackage{amsthm}
\makeatletter
\def\thm@space@setup{%
  \thm@preskip=8pt plus 2pt minus 4pt
  \thm@postskip=\thm@preskip
}
\makeatother
\usepackage[]{natbib}
\bibliographystyle{apalike}

\title{The Sacraments}
\author{James Bannerman}
\date{1868}

\begin{document}
\maketitle

{
\setcounter{tocdepth}{1}
\tableofcontents
}
\hypertarget{a-warhorn-historicals-book}{%
\chapter*{\texorpdfstring{\href{https://historicals.warhornmedia.com/}{A Warhorn Historicals Book}}{A Warhorn Historicals Book}}\label{a-warhorn-historicals-book}}
\addcontentsline{toc}{chapter}{\href{https://historicals.warhornmedia.com/}{A Warhorn Historicals Book}}

\includegraphics[width=0.5\textwidth,height=\textheight]{images/warhornlogo.png}

Republished by Warhorn Historicals:
Making historical books available for free online in high quality, readable formats.

We hope this book is a blessing to you. If it is, please \href{https://warhornmedia.com/give}{make a one-time or recurring contribution} right now, sponsor a book from our upcoming list, or volunteer your proofreading or technical skills to help produce more content. Contact \href{mailto:lucas@beggarsborn.com}{Lucas Weeks} to get involved.

God bless,

---The Warhorn Team

\hypertarget{cover}{%
\subsubsection*{Cover}\label{cover}}
\addcontentsline{toc}{subsubsection}{Cover}

\hypertarget{introduction}{%
\chapter*{Introduction}\label{introduction}}
\addcontentsline{toc}{chapter}{Introduction}

Nobody understands the Church like Bannerman, and this section on the Sacraments is mandatory reading today.

\hypertarget{the-sacraments-in-general}{%
\chapter{The Sacraments In General}\label{the-sacraments-in-general}}

FOR some time past we have been occupied with the subject of the ordinances of the Christian Church. We have discussed the questions connected with the public worship appointed in the Church, the special time set apart and sanctified for worship, and the ministry by means of which the worship of the Church is conducted. All these are outward ordinances which Christ has established in His Church, as parts of that external provision which He has made for the spiritual benefit and advancement of His people, and which He specially makes effectual to that end by the presence and power of His Spirit. All of these ordinances are in themselves, perhaps, and naturally adapted by their inherent character and influence to promote the edification of Christians; but above and beyond this natural or moral efficacy for that end, there is a spiritual blessing connected with them in consequence of the positive appointment of Christ, and the positive promise of His Spirit fulfilled in the right use of them. There may be a natural or moral efficacy in the ordinances of the Church considered in themselves, so that, apart from any other influence, they would, to a certain extent, be beneficial and advantageous in the case of those who used them. But in addition to this, there is a spiritual efficacy in the ordinances of the Church, distinct from the natural, and which is derived from the blessing of Christ and the working of His Spirit in them who by faith make use of them as He has appointed. What this spiritual and supernatural efficacy of outward ordinance exactly is,---what is the measure or amount of the inward benefit to the believer,---in what way and to what extent grace is connected with the external observance,---how beyond the sphere of this natural or moral influence the positive institutions of the Church have a blessing not natively their own,---these are questions which it is impossible for us distinctly to answer. The only wise and fitting reply to such questions is, that we have now reached the region of the supernatural, and that \emph{there} we have no data to guide us beyond what has been revealed. We know, from revelation, that there is a promise of grace annexed to outward ordinances when rightly used; we know that in the external observances Christ meets with His people to bless them and to do them good;---but beyond this we do not know. The character, the measure, the amount of the blessing promised,---how it stands connected with the outward ordinance, and what is the extent and efficacy of the supernatural grace over and above the natural efficacy of the ordinance,---of all this we know nothing, because we have been told nothing. We can distinctly understand, from the analogy of other cases, how the preaching of the Word, viewed as a system of human teaching of truth, and no more, may have a natural tendency to benefit the understanding and the heart. But we do not understand the supernatural efficacy which, over and above the natural, is imparted to it by the presence and the power of the Spirit in the ordinance.

In passing, as we do at this stage, from the non-sacramental to the sacramental ordinances appointed by Christ in His Church, it is of great importance to carry this general principle along with us. A supernatural grace is not peculiar to the Sacraments, although it may be found in them in larger measure than in other ordinances. It is common to all the ordinances which Christ has appointed in His Church. Whatever mystery there may be in the connection which by the promise of Christ has been established between the outward act and the inward blessing,---between the external observance rightly used and the internal grace divinely bestowed,---it is a mystery not belonging to Sacraments alone, but belonging to them in common with all Church ordinances. \emph{There is the supernatural element in them all.} There is that supernatural element connected in some manner with the outward act of the believer in the use of ordinances. There is a mystery in respect to any ordinance, not less than in respect of sacramental ordinances, which we cannot explain. It is, in short, the mystery of the Spirit of God, promised to dwell in the Church, and making every ordinance of the Church, whether sacramental or not, the channel for the conveyance of supernatural grace. If we would rid ourselves of this mystery, we can only do so by denying that the Spirit is present in ordinances at all. ``As the wind bloweth where it listeth, and thou hearest the sound thereof, but canst not tell whence it cometh, and whither it goeth,''---so is every ordinance, as well as each person, touched and sanctified of the Holy Ghost. There can be no natural explanation of the supernatural.\footnote{{[}Bannerman, \emph{Inspiration: The infallible Truth and Divine Authority of the Holy Scriptures}, Edin. 1865, pp.~217-228, 472 f.{]}}

What, then, is the character of those special ordinances instituted by Christ in His Church, which are usually denominated sacramental ordinances; and in what respect are they to be distinguished from the other ordinances of the Christian Church, not sacramental? In administering Sacraments, what is the peculiar nature or character of the Church's act; and in what manner does the administration differ from that of common ordinances?

The term \emph{Sacrament,} by which these peculiar ordinances are known, is not of scriptural, but of ecclesiastical origin; and there is some doubt as to the manner in which it came to be applied to these special solemnities of the Church, and to be restricted to the peculiar meaning in which it is now almost universally employed. In classical use, the word ``\emph{sacramentum}'' is almost always, if not invariably, employed to signify an oath,---more especially the military oath by which a soldier bound himself to obey the officer placed over him. And it has been conjectured that from its classical use it was transferred into the service of the Church, as significant of the obligation which the Christian comes under, in voluntarily participating in the Sacraments, to serve Christ as the Captain of his salvation,---these Sacraments being the characteristic badges or symbols by which the Christian is distinguished from other men. There is a second explanation, advocated by not a few, of the way in which the Latin term Sacrament came to be appropriated to its present ecclesiastical sense. It is the ordinary translation of the Greek word \emph{μυστηριον} among the ecclesiastical writers of the early ages, and more especially in the Vulgate and other old Latin translations of the Bible. The term Sacrament, according to this supposition, came to be employed to signify the ``mysteries'' of Christianity,---whether ``mystery'' is employed to denote a doctrine unknown until it was revealed, or a type or emblem bearing a hidden and secret meaning.\footnote{Turrettin, \emph{Opera}, loc. xix. qu. i. 1-6. Halley, \emph{The Sacraments}, Lond. 1844, pp.~7-14.} There is some reason to believe that both the Greek term \emph{μυστηριον} and the Latin translation of it---sacramentum---came at an early period to be applied by the primitive Christians to those special solemnities of their faith, which, although made up of outward and sensible signs or actions, bore in them a secret and spiritual meaning. In one or other of these ways, or perhaps in both, the term ``Sacrament'' soon came to be restricted in its meaning and application, by ecclesiastical practice, to those outward ordinances of Christianity which signify and seal its most precious and momentous truths. But as the term itself is of Church origin, and not found in Scripture, we must look not to it, but to the descriptions and intimations given in Scripture in regard to the ordinances themselves, for an explanation of their true nature and import.\footnote{{[}``The Apostle calleth the vocation of the Gentiles a mystery (Eph. iii. 4-6); the conjunction quhilk is begun here betwixt us and Christ is called a mystery (Eph. v. 32), and the Latin Interpreters call it a Sacrament; and, to be short, ye will not find in the Book of God a word mair frequent nor the word mystery. . . . Alwayis, the word Sacrament is very ambiguous in itself, and there raise about the ambiguity of this word many tragedeis quhilk are not yet ceased, nor will cease while the warld lasts; quher otherwise, gif they had keeped the Apostle's words, and called them, as the Apostle calls them, signs and seals, all this digladiatioun, strife, and contention appearandly had not fallen out. But quher men will be wiser than God, and give names to things beside God, upon the wit of man, quhilk is but mere folly, all this cummer falls out. . . . The ancient theologues took the word Sacrament in a fourfold manner. Sometimes they took it for the hail action, that is, for the hail ministrie of the elements. Sometimes they took it, not for the hail action, but for the outward things that are used in the action of Baptism and of the Supper; as they took it for the water and sprinkling of it, for the bread and wine, breaking, distributing, and eating thereof. Again, they took it, not for the hail outward things that are used in the action, but only for the material and earthly things,---the elements; as for bread and wine in the Supper, and water in Baptism. After this sort sayeth Augustine: `The wicked eats the body of our Lord concerning the Sacrament only;' that is, concerning the elements only. (Aug.~in \emph{Joann}. Tract xxvi. 18). Last of all, they took it not only for the elements, but for the things signified by the elements. And after this manner, Irenaeus saith, `that a Sacrament stands of twa things,---the ane earthly, the other heavenly.' (Adv. \emph{Haeres}. lib. iv. cap. 18.) The ancients, then, taking the word after thir sorts, na question all thir ways they took it rightly.''---Robert Bruce, \emph{Sermons on the Sacraments}, p.~6, Wodrow Soc. ed. Edin. 1843.{]}} In what respects, then, do the Scriptures represent the Sacraments of the Church as differing from its other ordinances which are not sacramental? What, according to Scripture, must we regard as the true nature and design of a Sacrament? To this general question we shall direct our attention in the first place, postponing for the present the special consideration of the Sacraments individually. And in endeavouring to ascertain the real nature and design of the Sacraments of the New Testament, we shall be enabled to understand at the same time, and by means of the same inquiry, in what respects they differ from other ordinances not sacramental.

\hypertarget{nature-and-efficacy-of-the-sacraments-of-the-new-testament-and-difference-between-them-and-non-sacramental-ordinances}{%
\section{Nature And Efficacy Of The Sacraments Of The New Testament, And Difference Between Them And Non-Sacramental Ordinances}\label{nature-and-efficacy-of-the-sacraments-of-the-new-testament-and-difference-between-them-and-non-sacramental-ordinances}}

\hypertarget{the-sacraments-of-the-new-testament-are-divine-institutions-appointed-by-christ.}{%
\subsection{The Sacraments of the New Testament are Divine institutions appointed by Christ.}\label{the-sacraments-of-the-new-testament-are-divine-institutions-appointed-by-christ.}}

It is the positive institution by Christ that sets these ordinances apart to the religious purpose for which they are intended, that makes them significant of spiritual things, and connects them with the virtue or blessing which they are made instrumental to impart. An express Divine appointment is necessary to constitute a Sacrament. In this respect they are similar to the other ordinances which form part of Church worship. Like them, they can claim Divine authority for their institution; and without this authority they would not be Sacraments at all. No observance not ordained by God can properly form any part of His service; far less can any observance not instituted by Him become a sign of His spiritual grace, or a pledge of a blessing which it depends upon His pleasure to give or to withhold. Hence, that any outward institution may answer to our idea of a Sacrament, it must be a positive appointment of God, and made both a sign and a pledge of spiritual blessings, in consequence of His promise and command. Without this, it would be a mere human ordinance, not only destitute of all real religious significance and efficacy, but profanely mimicking the form and character of a Divine ordinance in the Church. This is the first element that goes to make up a Sacrament, and which it has in common with all other ordinances, really forming a lawful or proper part of Divine worship,---namely, that it be of positive appointment by Christ.

\hypertarget{the-sacraments-of-the-new-testament-are-sensible-signs-of-spiritual-blessings-teaching-and-representing-by-outward-actions-gospel-truths.}{%
\subsection{The Sacraments of the New Testament are sensible signs of spiritual blessings, teaching and representing by outward actions Gospel truths.}\label{the-sacraments-of-the-new-testament-are-sensible-signs-of-spiritual-blessings-teaching-and-representing-by-outward-actions-gospel-truths.}}

The word or promise of God is an appeal to the understanding only; the Sacraments, embodying the same word or promise in outward and sensible signs, form a twofold appeal, \emph{first,} to the senses, and \emph{secondly,} to the understanding. There is Christ in the Word preached; and in the preaching of the Word, Christ is presented directly to the understanding and heart, and the truth addressed singly to the spiritual nature of man. But Christ is also in the Sacrament administered; and, in the administration of the Sacrament, over and above the same truth taught to the understanding and spiritual nature of man, there is the truth taught to the senses, and impressed by sensible signs upon them. There is a striking similarity between the method God has employed in the Sacraments of the New Testament to embody the Word and promises of Christ, and of a \emph{past} salvation, to the view of His people since His departure, and the method that He employed before Christ's coming to embody the Word and promises of a \emph{future} salvation. Under the Old Testament Church, there were, from the very first, two lines of promise and prediction,---both pointing forward to the coming of the Redeemer, running parallel with each other, and throwing mutual light upon each other's announcements. There was the line of promise embodied in verbal revelation, and there was the line of promise embodied in outward representation or type.

These two revelations ran parallel with each other since the first hour that a revelation was given to man in Paradise concerning the future coming of a Saviour. At that time there was a promise embodied in words, that ``the woman's seed should bruise the serpent's head, while His own heel was to be bruised;'' and side by side with that verbal announcement, there was the same promise embodied in type through means of the ordinance of sacrifice then appointed. There was Christ in the word of promise, and Christ in the sign of promise. When the promise was renewed to Noah, the second father of the human family, we have again the revelation by word, and the revelation by sensible sign; the covenant was repeated in another form, and the bow was set in the cloud as the outward representation of it. Once more: when Abraham was selected by God to be the depositary of a new development of the promise, we have again that promise embodied in words, and also in outward action; we have the special covenant with Abraham revealed in words, and revealed side by side with the word in the external sign of circumcision; and---to mention no further examples of a practice which must be familiar to every reader of the Old Testament---the whole of the Jewish economy was an exemplification of the two parallel lines that run through every economy of God,---the promise in word and the promise in sign revealed together, and throwing mutual light on each other. The typology of the Old Testament shows us God embodying His promises in signs; the revelation of the Old Testament shows us God embodying the same promises in words; and the Sacraments of the New Testament afford, under the Gospel economy, an exemplification of the same great principle.

The connection between the outward action in the Sacraments and the spiritual blessings to which they stand related is not a mere arbitrary one, arising from positive institution: there is a natural analogy or resemblance between the external signs and the things represented; so that, in the Sacraments of the New Testament, as in the types of the Old, our senses are made to minister to our spiritual advantage, and the outward action becomes the image of inward grace. In the Word, Christ is impressed on the understanding; in the Sacraments, Christ is impressed both on the understanding and the senses. They become teaching signs, fitted and designed to address to the believer the very same truths as are addressed to him in the Word; but having this peculiarity, that they speak at the same time and alike to the outward senses and to the inward thought. In this respect the Sacraments differ from other ordinances of the New Testament Church. Prayer and preaching and praise are ordinances that address themselves to the intellectual and spiritual nature of man alone. They are the expressions and utterances of his intellectual and spiritual being in holding intercourse with God; or they are the means fitted to speak to that nature, and that only, in impressing Divine truth upon men. But in those significant and teaching signs, which we call the Sacraments, Christ is embodied in the ordinance in such a manner as to appeal to the twofold being of man, as made up of body and soul, to minister both to the senses and the understanding; and to speak at once to the outward and inward nature of the believer. In addition to Christ in the Word, we have Christ also in the sign, taught as really in the latter way as in the former, and taught with the advantage of being submitted to the eye, and pictured to the outward senses. This, then, is one important difference between the sacramental ordinances of the New Testament Church and those which are not sacramental.

\hypertarget{the-sacraments-of-the-new-testament-are-federal-acts-affording-a-seal-or-confirmation-of-the-covenant-between-god-and-his-people.}{%
\subsection{The Sacraments of the New Testament are federal acts affording a seal or confirmation of the covenant between God and His people.}\label{the-sacraments-of-the-new-testament-are-federal-acts-affording-a-seal-or-confirmation-of-the-covenant-between-god-and-his-people.}}

This is the main and primary characteristic of sacramental ordinances. They constitute a formal testimony to an engagement entered into by two parties through means, not of words, but of speaking and significant actions,---these actions being the visible witnesses to the engagement, and the outward confirmations of its validity. In other words, they become, according to the expression of the apostle in his Epistle to the Romans, when speaking of one of the Sacraments of the Old Testament, visible ``\emph{seals}'' of the covenant, and of the blessings contained in it.\footnote{Rom. iv. 11.}

There are not a few examples to be found in the Old Testament Scriptures of covenants between man and man ratified by some outward monument, framed or chosen to attest and confirm the transaction. When Jacob parted from his father-in-law Laban, they made a covenant together, and raised a heap of stones and a pillar, to be a memorial of the transaction, and to serve as a witness on both sides to attest their fidelity to the terms of the covenant. ``This heap be a witness, and this pillar be a witness, that I will not pass over this heap to thee, and that thou shalt not pass over this heap and this pillar to me, for harm.''\footnote{Gen.~xxxi. 52.} The outward monument or memorial of the covenant entered into between Jacob and Laban was a witness of the engagement, serving to bind the obligation of it more strongly on both parties, and to ratify and confirm, in a formal and significant manner, its validity. And what we find in patriarchal times, we also find, in one shape or other, in every stage of society, some outward sign or significant action being made use of between men to confirm and attest their plighted faith. In addition to the spoken promise or oath, there has been---if not the stone of the times of Jacob---at least the formal signature and solemn deed, and seal attached to the deed, to remain after the verbal engagement, as the witness and ratification of the transaction. Such outward monuments or significant solemnities are intended for the satisfaction of both parties, and to give additional certainty and confirmation to the agreement. And the practice in this respect, which has obtained universally among men, we find to be made use of also by God. There are repeated examples in the Old Testament Scriptures of God ratifying His engagements or covenants with men by means of appropriate signs or solemnities, and making use of these solemnities for the very same purpose that a signed and sealed deed is employed for in the present day, when it attests or confirms a previous engagement, and gives additional security to both parties for the fulfilment of it. That in such a sense the rainbow in the cloud was employed by God, when it became the sign of His covenant with Noah, is very expressly stated by Himself: ``And the bow shall be in the cloud; and I will look upon it, that I may remember the everlasting covenant between God and every living creature of all flesh that is upon the earth. And God said unto Noah, This is the token of the covenant, which I have established between me and all flesh that is upon the earth.''\footnote{Gen.~ix. 16, 17.} In this point of view the bow was a seal, giving validity and additional security to the covenant then made, and serving as a standing witness for the truth of it. In a precisely similar manner, the rite of circumcision was appointed to Abraham for a voucher of the covenant between God and him. The terms of the institution of the rite would themselves lead us to this conclusion, even had they not been interpreted by the inspired commentary of the Apostle Paul in that sense. ``And, God said unto Abraham, Thou shalt keep my covenant therefore, thou, and thy seed after thee in their generations. This is my covenant, which ye shall keep, between me and you, and thy seed after thee. Every man-child among you shall be circumcised. And ye shall circumcise the flesh of your foreskin; and it shall be a \emph{token of the covenant} betwixt me and you.'' And in reference to this transaction, the Apostle Paul expressly says of Abraham: ``And he received the sign of circumcision, \emph{a seal of the righteousness of the faith} which he had yet being uncircumcised.''\footnote{Gen.~xvii. 9-11; Rom. iv. 11.} The outward act of circumcision, then, was a witness or a seal of the covenant transaction between God and the patriarch, and thus became a voucher to ratify and confirm the validity of it.

In exact accordance with the practice, universal in one shape or other among men, and expressly sanctioned by the example of God Himself in the Old Testament Church, we affirm that the Sacraments of the New Testament are parts of a federal transaction between the believer and Christ, and visible and outward attestations or vouchers of the covenant entered into between them. In addition to being signs to represent the blessings of the covenant of grace, they are also seals to vouch and ratify and confirm its validity. That the Sacraments of the Christian Church are thus seals of the covenant, appears to be very explicitly asserted, in so far at least as regards the Lord's Supper, in the words of the institution themselves: ``\emph{This} cup,'' said our Lord, ``\emph{is the new covenant in my blood, which is shed for you,}''\footnote{Luke xxii. 20.}---language which seems undoubtedly intended to convey the idea that the element used in the Supper was to be the witness of the new covenant,---a visible seal or security to ratify and vouch for it. No doubt that covenant in itself is sufficiently secure without any such confirmation, resting as it does on the word of God. That word alone, and without any further guarantee, is enough. But in condescension to the weakness of our faith, and adapting Himself to the feelings and customs of men, God has done more than give a promise. He has also given a guarantee for the promise,---has vouchsafed to bestow an outward confirmation of His word in the shape of a visible sign, appealing to our senses, and witnessing to the certainty and truth of the covenant. In the case of the Sacraments, God has proceeded on the same principle as is announced by the Apostle Paul in reference to His oath: ``God, willing more abundantly to show unto the heirs of promise the immutability of His counsel, confirmed it by an oath; that by two immutable things, in which it was impossible for God to lie, we might have a strong consolation, who have fled for refuge to lay hold upon the hope set before us.''\footnote{Heb. vi. 17, 18.} The word of promise was itself enough to warrant and demand the belief of God's people. But more than enough was granted: He has not only said it, but also sworn it. By two immutable things---His word and His oath---is the faith of the believer confirmed. The oath is the guarantee for His word. And more than this still: In the visible seal of the Sacraments God would add another and a third witness,---that at the mouth, not of two, but of three witnesses, His covenant may be established. He has not only given us the guarantee of His word, and confirmed that word by an oath, but also added to both the seal of visible ordinances. There is the word preached to declare the truth of the covenant to the unbelieving heart. More than that,---there is the oath sworn to guarantee it. More than that still,---there is the sign administered in order to vouch for all. Christ in the word, unseen but heard, is ours, if we will receive that word with the hearing ear and the understanding heart. Over and above this, Christ, both seen and heard in the Sacrament, is ours, if we will see with the eye or hear with the ear.\footnote{{[}``What mister (need) is there that thir Sacraments and seals suld be annexed to the Word? Seeing we get na new thing in the Sacrament but the same thing quhilk we gat in the simple Word, quherefore is the Sacrament appointed to be hung to the Word? It is true certainly, we get na new thing in the Sacrament, nor we get na other thing in the Sacrament nor we gat in the Word; for \emph{quhat mair walde thou crave nor to get the Son of God}, gif thou get Him weil? Thy heart cannot wish nor imagine a greater gift nor to have the Son of God, quha is King of heaven and earth. And therefore I say, quhat new thing walde thou have? For gif thou get Him, thou gettest all things with Him. Quherefore, then, is the Sacrament appointed? Not to get thee a new thing. I say it is appointed to get thee that same thing \emph{better} nor thou gat it in the Word. The Sacrament is appointed that we may get a better grip of Christ nor we gat in the simple Word; that we may possess Christ in our hearts and minds mair fully and largely nor we did of before in the simple Word; that Christ might have a larger space to make residence in our narrow hearts nor we could have by the hearing of the simple Word. And to possess Christ mair fully it is a better thing; for suppose Christ be ae thing in Himself, yet the better grip thou have of Him thou art the surer of His promise.''---Bruce, \emph{Sermons on the Sacraments}, Wodrow Soc. ed. Edin. 1843, p.~28.{]}}

The Sacraments are the outward and sensible testimony and seal of the covenant, added to the word that declares it. This is the grand peculiarity of sacramental ordinances, separating them by a very marked line from ordinances not sacramental. They are federal acts,---seals and vouchers of the covenant between God and the believer. They presuppose and imply a covenant transaction between the man who partakes of them and God; and they are the attestations to and confirmations of that transaction, pledging God by a visible act to fulfil His share of the covenant, and engaging the individual by the same visible act to perform his part in it. Other ordinances, such as the preaching of the Word, presuppose and attest no such personal engagement or federal transaction between the individual and God. Christ in the Word is preached to all, and all are called upon to receive Him; but there is no personal act on the part of the hearer that singles him out as giving or receiving a voucher of his covenant with his Saviour. But when the same individual partakes of the Sacraments, his own personal deed is an act of covenanting with God; and Christ in the ordinance is made his individually, and he is made Christ's by the very action of partaking of the ordinance. He is singled out by his own voluntary act, if he rightly partakes of the ordinance, as giving a voucher for his engagement with Christ; and Christ Himself gives a voucher of His engagement to the individual; and the visible Sacrament is the seal to the personal and mutual engagement. In this respect, as not only signs but seals of the covenant of grace to the individual who in faith partakes of them, the Sacraments are very markedly distinguished from ordinances not sacramental.

\hypertarget{the-sacraments-of-the-new-testament-are-made-means-of-grace-to-the-individual-who-rightly-partakes-of-them.}{%
\subsection{The Sacraments of the New Testament are made means of grace to the individual who rightly partakes of them.}\label{the-sacraments-of-the-new-testament-are-made-means-of-grace-to-the-individual-who-rightly-partakes-of-them.}}

It is carefully to be noted that they presuppose or imply the possession of grace in the case of those who partake of them; but they are also made the means of adding to that grace. They are seals of a covenant already made between the soul and Christ,---attestations of a federal transaction before completed,---confirmations, visible and outward, of engagement between the sinner and his Saviour previously entered into on both sides. They presuppose the existence of grace, else they could not be called seals of it. Just as the signature and seal of some human covenant necessarily presuppose that the covenant exists before they can become vouchers for it, so the seal of God's covenant, affirmed by means of sacramental ordinances, presupposes the existence of that covenant as already subsisting between God and the rightful participator in the ordinance. But although grace exists in the soul before, the Sacraments are made to those who rightly receive them the means of increasing that grace, and communicating yet more of spiritual blessing. They serve to strengthen the faith of those who already believe, and add to the grace of those who previously possessed grace. They become effectual means of imparting saving blessings in addition to those enjoyed before.\footnote{{[}``The Church has always seen in the Sacraments,'' says Mr.~Liddon in his recent very valuable work on the Divinity of our Lord, ``not mere outward signs addressed to the taste or imagination, nor even signs, as Calvinism asserts, which are tokens of grace received independently of them, but signs which, through the power of the promise and Word of Christ, effect what they signify.'' For this very defective statement of the Calvinistic doctrine of the Sacraments the only authority Mr.~Liddon gives is a single secondhand quotation from Cartwright. He then proceeds to contrast with this supposed Calvinistic view the words of the 25th Article: ``The Sacraments are \emph{effectual} signs of grace and God's goodwill toward us, by which He doth work invisibly in us;'' and the definition of the Church Catechism: ``A Sacrament is an outward and visible sign of an inward and spiritual grace given unto us, ordained by Christ Himself as a \emph{means whereby} we receive the same, and a pledge to assure us thereof.'' \emph{Bampton Lectures}, 1866, p.~721. A very slight reference to the symbolical books or the leading theologians of some of the Calvinistic Churches would of course have shown that all these phrases have been constantly used by them with respect to the Sacraments. The isolated sentence from Cartwright adduced by Hooker, from whom Mr.~Liddon takes it, refers to a particular aspect of a particular Sacrament; it was never designed to be a full definition of the efficacy of these ordinances in a typical case. Moreover, the passage in question is just a translation of Calvin, \emph{Inst}. iv. xv. 22. It might as well, therefore, have been brought forward as expressive of \emph{his} whole doctrine on the subject. But Mr.~Liddon must surely be aware that Calvm constantly speaks of the Sacraments both of the Old and New Testaments as ``\emph{effectual} means of grace,'' ``efficacious instruments,'' ``signs in which God \emph{gives} what He holds out to us,'' etc. (``non modo salutaria exercitia, et adjumenta pietatis, sed etiam \emph{eficacia gratiae instrumenta}.'' ``\emph{Praestat} igitur vere Deus quicquid signis promittit ac figurat; nec effectu suo carent signa, ut verax et fidelis probetur eorum Author'').---\emph{Comment. in} Gal. iv. 9, Col. ii. 17, \emph{Inst}. iv. xiv. 17, etc. Cf. i. \emph{Conf. Helv.} c. 21, ii. c. 21. Conf. Gall. Art. 37, Catech. Gen.~v. etc. Some of the expressions in the Church Catechism, indeed, with respect to Baptism seem to Presbyterians to require at least all the explanation which Dean Goode and others have bestowed upon them. And the passage from Martensen about the ``communication of Christ's glorified corporeity'' in the Lord's Supper, which Mr.~Liddon quotes, seemingly as supplementary of the Catechism, would of course be disapproved of by Calvinists generally, although there are statements in the works of Calvin himself which might perhaps be adduced in its favour. Mr.~Liddon concludes by observing that, ``though there have been and are believers in our Lord's Divinity who deny the realities of sacramental grace, experience appears to show that their position is only a transitional one.'' There is ``a law of fatal declension,'' which will ultimately, Mr.~Liddon thinks, bring all who do not hold the High Church doctrine of the Sacraments to the Socinian position. ``Centuries,'' however, ``may intervene between the premisses and the conclusion;'' so that the prediction is a singularly safe one. By a precisely similar process of reasoning, Dr.~Manning and others are prepared to prove that there is an indissoluble connection between the worship of the virgin and a belief in the Divinity of Christ.---\emph{Engl. and Christend.} p.~civ. Faber, \emph{Growth in Holiness}, p.~72.{]}} In this respect they are similar to the other ordinances which Christ has appointed in His Church, and which by His power and Spirit are made instrumental in advancing the interests of His people. But from the very peculiarity that attaches to their distinctive character, as seals of a personal covenant between God and the believer, Sacraments may reasonably be supposed to be more effectual than non-sacramental ordinances in imparting spiritual blessings. The spiritual virtue of Sacraments is more and greater than other ordinances, just because, from their very nature, they imply more of a personal dealing between the sinner and his Saviour than non-sacramental ordinances necessarily involve.

What is the nature and extent of the supernatural grace imparted in Sacraments,---in what manner they work so as to impart spiritual benefit to the soul, it is not possible for us to define. As visible seals of God's promises and covenant, we can understand how they are naturally fitted, in the same way as the vouchers of any human engagement or covenant are naturally fitted, to attest and confirm them. But beyond this, all is unknown. The blessing of Christ and the working of His Spirit in Sacraments we cannot understand, any more than we can understand the operation of the same supernatural causes in respect of other ordinances. They have a virtue in them beyond what reason can discover in them, as naturally fitted to serve the purposes both of signs and seals of spiritual things. They have a blessing to the right receiver of them, not their own to give. ``They are made effectual means of salvation, not from any virtue in them, or in him that doth administer them, but only by the blessing of Christ, and the working of His Spirit in them who by faith receive them.''\footnote{Shorter Catechism, qu. 91.} In this respect their power and virtue are not more and not less mysterious than those of ordinances non-sacramental.

Such are the general conclusions which a consideration of the nature of the Sacraments of the New Testament lead us to acquiesce in. They are Divine institutions appointed by Christ; they are signs and significant representations of spiritual things; they are seals and vouchers of a federal transaction between God and the worthy receiver of Sacraments; they are the means for applying spiritual grace to the soul. To sum up the discussion in the language of the Shorter Catechism: ``A Sacrament is an holy ordinance instituted by Christ, wherein by sensible signs Christ and the benefits of the new covenant are represented, sealed, and applied to believers.''\footnote{Shorter Catechism, qu. 92. Calvin, \emph{Inst}. lib. iv. cap. xiv. \emph{Consensus Tigurinus} in Niemeyer's \emph{Collectio Confess}. Lipsiae 1840, pp.~192-217, translated in Calvin's \emph{Tracts}, Edin. 1849, vol.~ii. pp.~205-244. Turrettin, \emph{Opera}, tom. iii. loc. xix. qu. i.-ix. Cunningham, \emph{Works}, vol.~i. pp.~225-291, vol.~ii. pp.~201-207, vol.~iii. pp.~121-133. Amesius, \emph{Bellarm. Enerv.} tom. iii. lib. i. cap. i. Willison, \emph{Works}, Hetherington's ed. pp.~456 f. Gillespie, \emph{Aaron's Rod Blossoming}, B, iii. chap, xii.-xiv. Mastricht, \emph{Theol. Theoretico-Pract.} tom. ii. lib. vii. cap. 3.}

Sacraments and non-sacramental ordinances are like each other in two respects; and in two respects they differ. In the first place, sacramental and non-sacramental ordinances agree in this: first, that they are both positive institutions of Christ; and second, that they are both means of grace to believers. Without a Divine warrant and institution, neither non-sacramental ordinances nor Sacraments could have any place in the worship of God as part of His service; and both are therefore Divine appointments. They are both likewise means of grace to believers,---there being a positive promise attached to the right use of them, and that promise being fulfilled in the bestowment of spiritual blessing in connection with their use. This spiritual benefit, linked to the proper use of ordinances, whether sacramental or not, is over and above and quite distinct from the natural or moral influence such ordinances may have to benefit those who employ them. There is a benefit, for example, which the ordinance of preaching the Word is naturally fitted to impart, because the truth preached is adapted to man's moral and intellectual nature, and so naturally fitted to be of advantage to the hearers. In like manner there is a benefit which Sacraments are naturally fitted to impart, because they are symbolical ordinances or teaching signs; and the truths represented or taught by them are, upon the very same principle, naturally fitted to be of advantage to the receiver. But in both cases there is a blessing distinct from and additional to the natural or moral effect of the Word preached or the Sacraments administered. There is the work of the Spirit making use of Word and Sacrament to reach the understanding and the heart, and to convey to the worthy hearer or worthy receiver a spiritual blessing. And this work of the Spirit, over and above the natural effect of the truth received, is a mystery, both in the case of the ordinance of preaching and the ordinance of the Sacraments; and not, I think, a greater mystery in the one case than in the other.

We do not plead for the Sacraments as means of grace, viewed merely as natural actions and ceremonies apart from the truths which they represent, any more than we would plead for the preaching of the Word being a means of grace, viewed as the mere letter of the Word apart from the meaning of the truth which is uttered. The case of \emph{infant} Baptism, which is, as we shall afterwards see, in some respects exceptional, and not to be taken as completely bringing out the full and primary idea of the Sacrament,\footnote{Cunningham, \emph{Works,} vol.~iii. pp.~144-154.} we for the present put aside, postponing it for future consideration. But in the case of adult participation in the Sacraments, we do not plead for these generally as means of grace, when viewed simply as outward acts, and apart from the truths represented, any more than the sound of the Word preached would be a means of grace apart from the intelligent apprehension of it. Through the truths, however, in one case impressed on the hearer by significant words, and in the other case impressed on the participator through significant actions, the Spirit of God \emph{does} operate upon the intellectual and moral nature of man, making both the one ordinance and the other a means of grace. How the Spirit thus operates and imparts of His gracious gifts, we cannot tell in the one instance more than in the other. What is the mode or measure of His communications of a spiritual kind, over and above the natural or outward influence of the truth, we cannot tell. It is His own secret and supernatural work, known and recognised by the believer in the effects wrought on His soul, both in the case of the Word preached and the Sacraments administered, but not to be explained or defined in the manner of working. Let it never be forgotten that there is a mystery not to be explained whenever we get beyond the natural effect of the ordinance, whether sacramental or not, necessarily resulting from the fact that it is an effect of the Spirit, and not of any natural cause. All ordinances, as means of grace, must in that character have something in them mysterious and inexplicable. We cannot rid ourselves of the mysterious by simply ridding ourselves of sacramental ordinances,---as very many in the present day seem to imagine. We can only disconnect all mystery from the ordinances of the Church when we limit their efficacy simply to their natural influence, and deny the influence of the Spirit of God as at all connected with them.

In the second place, Sacraments differ from ordinances not sacramental in the New Testament Church, in these two things: first, they are sensible signs of spiritual truths; and second, they are seals or vouchers of a federal transaction. In respect that they are sensible exhibitions and significant actions, having a definite meaning in them, Sacraments stand out distinctly marked from other ordinances. Speaking generally, sacramental ordinances are spiritual acts of the mind or soul embodying themselves in outward and sensible actions, in so far as regards the part of the receiver in the ordinance. They are outward representations, by means of certain actions on the part of the worthy participator, of the great fact that he gives himself to Christ according to the terms of the covenant of grace. In partaking of the ordinance, he embodies in the sensible actions of the ordinance a spiritual surrender of himself to Christ, in the manner and upon the terms which Christ has appointed. This is the receiver's part in the ordinance. On the other side, Christ, through the person of the administrator of the ordinance, embodies in the actions of it a picture or representation of a spiritual communication of Himself and all the blessings of His grace to the worthy receiver. Christ, in the Sacrament, and by means of its sensible signs, gives Himself and the benefits of the new covenant, spiritually, although under an outward representation, to the believing participator. The outward signs of the Sacrament exhibit, then, a twofold action: the believer giving himself to Christ in covenant, and Christ giving Himself to the believer in the same covenant. There is a spiritual act on the part of the believer embodied in outward representation,---the act, namely, of his surrendering of himself to Christ in the way and on the terms which Christ has appointed; and there is a spiritual act on the part of Christ embodied in outward representation also,---the act, namely, of Christ with all His precious and unspeakable blessings communicating Himself to the soul of the worthy receiver. There is thus a double significance comprehended in the administration and in the participation of the sacramental ordinance, each of them having a definite and intelligible meaning of its own. In the administration of the Sacrament, Christ makes over Himself and all the benefits of His atonement to the believer, and accepts in return the believer as His. In the participation of the Sacrament on the part of the worthy receiver, he makes over himself to Christ; and receives, in return for his own soul, Christ and His covenant blessings. The double action of the administration and participation of the Sacrament is the embodiment in outward sign of a double spiritual act. There is a mutual intercommunication spiritually of Christ and the believer embodied and represented in action,---a covenant interchangeably exhibited in sensible signs, whereby Christ becomes the believer's, and the believer becomes Christ's. In their being signs of spiritual truths, Sacraments differ in a marked manner from non-sacramental ordinances.

Sacraments differ also from other ordinances in this, that they are seals or vouchers of a federal or covenant transaction. This, after all, is the grand and essential distinction between sacramental and non-sacramental ordinances. As a kind of types, as speaking and teaching signs, they are fitted to express, by the help of significant actions cognisable by the senses, the twofold spiritual act of Christ making over Himself and all His blessings to the believer, and of the believer making over himself with all his poverty and sins to Christ. But they are more than signs of a covenant thus entered into between the two parties,---they are seals and vouchers for the covenant, serving to give confirmation and validity to the engagement, as one never to be broken. In the Sacraments there is a twofold seal, as well as a twofold action, represented. There is a seal on the part of Christ, and there is a seal on the part of the believer. In marvellous condescension to our infirmity and unbelief, Christ has been pleased to add to the promise of His covenant an outward and visible voucher for it,---thereby, as it were, binding Himself doubly to the fulfilment of it, and pledging Himself, both by word and by sign, to implement all its terms. And in the worthy receiving of the Sacrament, the believer gives also a visible voucher for his part of the engagement,---thereby placing himself under new and additional obligations to give himself to Christ, and adding the outward seal to ratify the inward pledge of his heart. The covenant is mutual, and the seal is mutual. Without either part of the covenant transaction, the Sacrament would be incomplete. Withdraw Christ from the ordinance as both entering into covenant with the believer and giving him a seal of it,---take away Christ sealed to the soul in the Sacrament,---and the ordinance is reduced to a bare sign of spiritual blessing, having, perhaps, a certain natural effect by signifying truth, but empty and destitute of all spiritual grace. Or withdraw the believer from the ordinance in so far as he really by means of it gives himself to Christ,---take away the spiritual act by which the worthy participator surrenders his soul to the Saviour through his outward participation of the Sacrament,---and the Sacrament is made to be a charm, in which Christ and grace are communicated apart from the spiritual act or state of the receiver. Abstract from the ordinance the act of Christ covenanting with the believer and giving to the soul Himself and His blessings, and the remaining portion of the ordinance may continue,---the believer may still be accounted as giving himself to Christ in the Sacrament; but in the absence of Christ's act there is no spiritual blessing given in return, and the believer's act of participating in the Sacrament becomes a mere sign of adherence to Christ on his part, and nothing more than a sign.\footnote{{[}``Quod omnes fere opinantur, hoc ritu, quem Sacramentum appellant, confirmari saltem fidem nostram, ne id quidem verum censeri debet; cum nec ullo sacro testimonio comprobetur, nec ulla ratio sit cur id fieri possit. Quomodo enim potest nos in fide confirmare id \emph{quod nos ipsi facimus, quodque, licet a Domino institutum, opus tamen nostrum est?}''---Faustus Socinus, \emph{De Caena Dom, Tract. Brev.} Racovian Catechism, 1609, p.~144 f.{]}} Again, abstract from the ordinance the act of the believer spiritually covenanting with Christ and giving his soul in faith to the Saviour, and the remaining portion of the ordinance may continue,---Christ may be held as present in the Sacrament giving Himself and His supernatural grace; but in the absence of the receiver's act surrendering his soul in faith to his Saviour, the communication of spiritual grace is degraded to the position of being the result of a charm or talismanic formula,---something effected, \emph{ex opere operato,} apart from the spiritual character or faith of the receiver. It is only when the separate spiritual acts of both parties meet and combine in one transaction, that the covenant is real or complete; or that the ordinance, as a seal of the mutual engagement, is a true and proper Sacrament. As the voucher or seal of a real covenant, spiritually entered into between Christ and the believer through the ordinance, a Sacrament differs, in a very marked and important way, from ordinances not sacramental.

\hypertarget{unscriptural-or-defective-views-of-the-sacraments}{%
\section{Unscriptural Or Defective Views Of The Sacraments}\label{unscriptural-or-defective-views-of-the-sacraments}}

The principles which I have laid down in regard to the nature of Sacraments, and in regard to the difference between them and ordinances not sacramental, stand opposed to the views of two parties holding extreme positions on either side of this question. There is one party who deny the grand and characteristic distinction between sacramental and other ordinances already enunciated, and hold that the Sacraments have no virtue except as badges of a Christian profession, and signs of spiritual truths. There is another party holding opinions on the subject admitting of various modifications, but agreeing in this, that they ascribe a high spiritual efficacy to the Sacraments apart from the faith or spiritual act of the receiver. By the first party the views of the Sacraments already stated by me are held to be erroneous in the way of attributing to them a greater virtue than actually belongs to them. By the second party these views are regarded as defective in the way of ascribing to Sacraments a less virtue than really belongs to them. Let us endeavour briefly and generally to estimate the merits and truth of the principles adopted by these two parties,---reserving until a future stage in our discussions the more particular examination of their theories, in their application to the Sacraments of the New Testament individually.

\hypertarget{signs-and-no-more-than-signs-of-spiritual-things}{%
\subsection{Signs, and no more than signs, of spiritual things}\label{signs-and-no-more-than-signs-of-spiritual-things}}

The Sacraments of the New Testament are regarded by one party as signs, and no more than signs, of spiritual things,---symbolical actions fitted to represent, and impress upon the minds of men, Gospel truths. The Socinian party have made this doctrine peculiarly their own. According to their views, a federal transaction between the believer and Christ founded on His atonement is no part of the Gospel system at all; and hence the Sacraments of the New Testament can be no seals appointed and designed to ratify such a covenant. The Socinian doctrine concerning the nature of the Sacraments allows to them no more than a twofold object and design. They are not essentially distinct from other ordinances, as set apart by themselves to be the seals of the one great covenant between the believer and Christ, at his entrance into the Church at first, and from time to time afterwards, as occasion justifies or demands. But in the first place, they are signs in which something external and material is used to express what is spiritual and invisible,---the only virtue belonging to them being what they are naturally calculated to effect, as memorials, or illustrations, or exhibitions of the important facts and truths of the Gospel; and in the second place, the Sacraments are solemn pledges of discipleship on the part of those who receive them, discriminating them from other men, and forming a public profession of or testimony to their faith as Christians. These are the two grand objects, which, according to the Socinian view, the Sacraments were intended to serve; and such, according to their theory, is the nature of the ordinance.

The same system in substance, making, as it does, Sacraments entirely or essentially teaching and symbolical signs, has been adopted by many who disown the tenets of Socinianism in regard to the Gospel system generally. The theory of the Sacraments now described has been and is held by not a few in the Church of England of somewhat latitudinarian views,---the representative of such, as a class, being Bishop Hoadly. It is avowed and advocated in the present day by a very large proportion of the Independent body, who count the Sacraments to be no more than symbolical institutions, and who are ably represented by Dr.~Halley in his work, entitled, \emph{An Inquiry into the Nature of the Symbolic Institutions of the Christian Religion, usually called the Sacraments.} The single difference between the Socinian doctrine, as maintained by Socinians in the present day, and the Independent doctrine, as maintained by Dr.~Halley and others, is probably this, that Socinians limit the efficacy of the Sacraments to the natural or moral power that belongs to them as signs of Gospel truth, while Independents may admit that beyond the natural and moral power of the ordinance, as symbolical of truth, the Spirit of God makes use of them in representing truth to the mind. Let Dr.~Halley speak his own views as they are generally held by English Independents. ``The opinion we propose is, that the Sacraments are significant rites,---emblems of Divine truth,---sacred signs of the evangelical doctrine,---designed to illustrate, to enforce, or to commemorate the great and most important truths of the Gospel. Baptism, we believe, is the sign of purification, on being admitted into the kingdom of Christ, but neither the cause nor the seal of it; the Lord's Supper the commemoration of the death of Christ, the symbol of its propitiatory character, but not the assurance of our personal interest in its saving benefits. The truth exhibited in the Sacraments, just as when it is propounded in words, may be the means of the communication of Divine grace; but then the evangelical doctrine and not the Sacrament, the truth and not the symbol, the spirit and not the letter, gives life and sanctity to the recipient, as it may even to a spectator.''\footnote{Halley, \emph{The Sacraments: an Inquiry,} etc. Lond. 1844, vol.~i. p.~94 f.} According to this theory, it is the truth signified in the Sacrament---and not, over and above that, the Sacrament itself as a seal---that possesses any spiritual virtue; and that virtue may be, according to Socinians, the natural influence of the truth on the mind,---or, according to Independents, that natural influence, with the addition of the power communicated through the truth by the Spirit.

Now, in reference to this view of the Sacraments, it is necessary to bear in mind that there is no dispute as to the fact that sacramental ordinances are symbolical,---signs fitted to represent and to teach Gospel truths. Further, there is no dispute as to the fact, acknowledged by some of the advocates of this theory, that in so far as they teach or convey truth to the mind, they may be made the means of the communication of Divine grace, in the same manner very much as when the truth is propounded in words.\footnote{{[}``Es geschah in dem Zeitpunkte der Reformation, aber nicht zum ersten Male, dass die hypermystische oder zauberische Vorstellung den entgegengesetzten Fehler, die Behauptung des \emph{signum nudum} oder des blossen Bekenntnisszeichens, hervorrief. Gegen diejenige Kirche, die im Dienste der Verwandlungslehre und des \emph{opus operatum} die symbolische Natur des Sacraments verleugnete und zerstörte, hatte die sogenannte Ketzerei allezeit Recht, zunächst nur wieder das Daseyn des Symbols und die Bedeutung zu behaupten. Diejenige Kirche, die des Sacramentes Wirkung und Wesen vom lebendigen Worte und Glauben, den Sohn vom Geiste losgerissen hatte, durfte einen Gegner nie Lügen strafen oder des Unchristenthums zeihen, der der Gemeinschaft des Erlösers durch die Speise des Wortes als durch die rechte Assimilation mit seinem Leben theilhaft zu werden hoffte, und sich des Sacramentes nur noch als eines Zeichens dieser Gemeinschaft, oder auch dieses Zeichens nicht mehr bediente weil es so sehr vom Wesen abgelenkt und etwa nur habe bei noch nicht ganz befestigter Wirksamkeit des Wortes einem anfänglichen Bedürfnisse dienen sollen.'' (This is still the position of the Quakers as expounded by Barclay in his \emph{Apology.}) ``\emph{Blosse Gebetschristen,} Messalianer und dergleichen, \emph{sind nicht weniger Christen als blosse Sacramentschristen; blosse Symboliker stehen sich nicht schlechter mit der Quelle des Lebens als die Hierurgen die den Leib Christi conficiren.} Diese sind am Ende des verschwindenden Christenthums angelangt, jene stehen am Wiederanfange der Entwickelung.''---Nitzsch, \emph{prot. Beant. der Symb. Möhlers,} Hamburg 1835, p.~162.{]}} But the point in debate is, whether the Sacraments are not more than signs, and more than merely symbolical representations of truth. We hold that they are. We contend that, in addition to being signs, they are also seals,---the visible vouchers of a federal transaction between Christ and the believer who partakes of His Sacraments,---the outward pledges speaking to the eye and the senses of the completed covenant by which Christ becomes the believer's, and the believer becomes Christ's. And further, we contend that, as seals, they are made a means of grace more powerful and efficacious than simply as signs of truth.

The arguments urged by Dr.~Halley against this additional office and virtue attributed to Sacraments as more than signs, and as the seals of a federal engagement between the worthy recipient and Christ, are the two following, as stated in his own words: ``First, The ceremonial institutes of preceding dispensations, the Sacraments of the patriarchal and Jewish Church, correspond only with the view which we take of the Christian Sacraments as sacred signs of Divine truth. Second, The Sacraments considered as the causes or the means, or even the seals of converting or regenerating grace, stand opposed to the great Protestant doctrine of justification by faith without works.''\footnote{Halley, p.~95.} We shall very briefly examine each of these two objections to the view which we have announced. And we do this all the more readily, as it will afford us the better opportunity of bringing out our own principles in contrast with those embodied in the Independent theory of the Sacraments.

1\emph{st,} Dr.~Halley alleges, against the ascription to the New Testament Sacraments of the character of seals, that the ceremonial institutes of preceding dispensations, the Sacraments of the patriarchal and Jewish Church, correspond only with the views which he advocates of the Christian Sacraments as exclusively signs of Divine truth. Perhaps there never was a more unfortunate or unfounded assertion. ``One passage of St.~Paul,'' says Dr.~Halley, ``will establish this proposition.''\footnote{Ibid. p.~96.} And the single passage which is to bear the weight of the whole argument is the following one from the Epistle to the Romans: ``He is not a Jew which is one outwardly; neither is that circumcision which is outward in the flesh. But he is a Jew which is one inwardly; and circumcision is that of the heart, in the spirit, and not in the letter; whose praise is not of men, but of God.''\footnote{Rom. ii. 28, 29.} This is the solitary passage quoted to prove the broad and general assertion, that the Sacraments of the patriarchal and Jewish Church afford no precedent or example of Sacraments as seals, but only of Sacraments as signs. The verses quoted plainly amount to nothing more than a statement of the difference between what the apostle calls circumcision outwardly and circumcision inwardly, the external rite and the internal grace, and a declaration that a man might have the outward rite, and not the inward grace. The apostle does not say, and cannot, except by a violent misapplication of his words, be made to say, that in the case of the man who has both the outward and inward circumcision, the external rite may not be the visible seal of the spiritual grace. The very opposite of this the same apostle in the very same Epistle undeniably asserts. In language as plain as he could possibly select or employ, Paul affirms that in the case of Abraham, who had the inward grace, the outward rite of circumcision was a seal to him of that grace. ``Abraham,'' says the apostle, ``received the sign of circumcision, a seal of the righteousness of the faith which he had yet being uncircumcised.''\footnote{Rom. iv. 11.} And how is it that Dr.~Halley gets rid of this express assertion of the apostle, standing as it does in explicit contradiction to his general averment that the Sacraments of the Jewish Church were signs and not seals? He admits that to Abraham personally and individually circumcision was a seal, and not merely a sign. But by a strange misapprehension of the doctrine of his opponents, he argues that it could not be a seal of faith to others of Abraham's family or countrymen who had not his faith. ``Although,'' says Dr.~Halley, ``to him circumcision was the seal of faith, it could not have been so to his posterity.'' ``Was it,'' he asks, ``was it, in this sense, a seal of the righteousness which they had, an approval of their faith, to the men of his clan, or to Ishmael, or to the infants of his household, or to any of his posterity in subsequent ages?''\footnote{Halley, p.~100.} The answer to such a question is abundantly obvious. If the men of Abraham's clan had not faith, if Ishmael had not faith, circumcision could have been no seal of faith to them. The outward rite could not be a seal of the inward grace, when the latter did not exist. It could not be a seal of a spiritual covenant between them and God which had not been entered into. I do not stop to consider the question of whether or not circumcision is to be accounted, even in such a case, the seal to such individuals of the outward blessings promised to them, as Jews, by God, as the rightful King of Israel as a nation; but, as a seal of a spiritual covenant, it of course could not be a seal at all to those who were not parties to the covenant,---while it was a seal, according to the explicit assertion of the apostle, to those who were. The very express statement of Paul cannot be evaded, but fully bears out the assertion that the Sacraments of the Jewish Church were not signs alone, but seals of a spiritual covenant to those who were really parties to the covenant. ``Abraham received the sign of circumcision, a seal of the righteousness of the faith which he had.''

2d, Dr.~Halley alleges that the Sacraments, if they are considered as the cause or the means, or even the seals of spiritual and saving grace, would be opposed to the great Protestant doctrine of justification by faith without works. Now it is readily admitted, that if Sacraments are regarded as the causes or means of justification, they are utterly inconsistent with the Protestant doctrine of justification by faith alone; and in this point of view the objection is true and unanswerable when directed against some of those theories of the Sacraments which we may be called upon to consider by and by. But it is denied that the objection is true when directed against the theory of the Sacraments which maintains that they are not causes and not means of justification, but seals of it and of other blessings of the new covenant. The Sacraments as seals, not causes of justification, cannot interfere with the doctrine of justification by faith, for this plain reason, that before the seal is added, the justification is completed. The seal implied in the Sacrament presupposes justification, and does not directly or instrumentally cause it; the seal is a voucher given to the believer that he is justified already, and not a means or a cause of procuring justification for him. Justification exists before the seal that attests it is bestowed. The believer has previously been ``justified by faith without the works of the law,'' ere the Sacrament of which he partakes can affix the visible seal to his justification. All this is abundantly obvious; and the objection of Independents, that the doctrine of the Sacraments as personal seals is opposed to the principle of justification by faith, is wholly without foundation. That the Sacraments are a means of grace additional to what the believer possessed before his participation in them, it is not necessary to deny, but rather proper strongly to assert.\footnote{{[}``Das Gläubigste, so zu sagen, am sacramentlichen christlichen Gemeinglauben ist doch wohl dieses: je mehr das Sacrament mit voller Empfänglichkeit genossen wird, desto weniger ist es blosses Zeichen, oder blosses Unterpfand der Lebensmittheilung Christi, desto mehr diese Mittheilung selbst. Das Sacrament ist Leiter, Kanal der Gnade, wie der römische Katechismus sich ausdrückt. Bis auf diesen Punkt wird der Sacramentsbegriff---ich will zugeben unter sehr verschiedenen Bedingungen (from those of the Romanist theory)---durch das in dieser Hinsicht ganz ungetheilte Bekenntniss der Protestanten gesteigert. . . . Der protestantische Begriff des Siegels oder Pfandes ist weit entfernt die \emph{collative} Kraft des Sacraments zu schwächen; er gestattet sogar die mystische Verknüpfung der Elemente des Sacraments mit der \emph{res signata et exhibenda; signa et res significatae sacramentaliter conjunguntur} (\emph{Conf. Helv. post.} xix.). \emph{Bezeichnung, Besiegelung, Darreichung} der Gnade Christi vereinigen sich im Sacramente. (\emph{Decl. Thorun. De Sacr.} 1, 7).''---Nitzsch, \emph{prot. Beant, der Symb. Möhlers,} Hamburg 1835, p.~151 f.{]}} In entering into a personal covenant with Christ through particiption in the Sacraments, or in renewing that covenant from time to time, the faith of the believer is called forth and brought into exercise in the very act of participation, and by the aids to faith which the ordinance affords. And in answer to this faith so exercised and elicited, there is an increase of grace given to the worthy recipient above and beyond what he had before. The faith of the believer, called into exercise in partaking of the ordinance and by means of it, is met by the bestowment of corresponding grace. But it is never to be forgotten that the Sacraments presuppose the existence of grace, however they may give to him that already has it more abundantly. They presuppose, and beforehand require, that a man is justified by faith before they give their seal to his justification.

There is no ground, then, in Scripture, but the very opposite, for asserting that the Sacraments are no more than signs or symbolical actions, as held by Dr.~Halley and those whom on this question he represents. The fundamental error involved in the views now adverted to is, the denial of Christ's part in the federal transaction involved in a Sacrament. Independents overlook His department of the work in the engagement entered into through means of the act of receiving the Sacraments; and in the absence of the act of Christ giving Himself and all His spiritual blessings to the believer in the ordinance, the act of the recipient is not met by the grace that Christ confers, but is reduced to a mere significant dedication of himself to the Saviour unconnected with any grace at all. Take away Christ from the ordinance as present there, to covenant with the believer, actually giving Himself and His blessings spiritually through means of the outward ordinance, in answer to the faith of the believer giving himself to Christ through the same ordinance, and the Sacrament is evacuated of all spiritual grace; the act of the receiver becomes a mere expressive sign of what he is willing to do in the way of dedicating himself to Christ; but not an actual dedication, accomplished through means of a covenant then and there renewed, by which the believer becomes Christ's, and Christ becomes the believer's. The principle of the Independents in regard to the Sacraments cuts the Sacrament, as it were, in twain, and puts asunder what God has joined. It leaves to the believer his part in the transaction, in so far as he employs the Sacrament as a sign of his dedication to Christ; but it takes away Christ's part in the transaction, in so far as He meets with the believer and enters into covenant with him,---accepting the believer as His, and giving Himself to the soul in return. Severed from Christ in the ordinance, and from the covenant with His people into which Christ there enters, the act of the recipient can be no more than an expressive sign, or convenient profession of faith, unconnected with true and proper sacramental grace.

\hypertarget{effectual-to-impart-justifying-and-saving-grace-directly}{%
\subsection{Effectual to impart justifying and saving grace directly}\label{effectual-to-impart-justifying-and-saving-grace-directly}}

The Sacraments of the New Testament are regarded by another party as in themselves, and by reason of the virtue that belongs to them, and not through the instrumentality of the faith or the Spirit in the heart of the recipient, effectual to impart justifying and saving grace directly, in all cases where it is not resisted by an unworthy reception of the ordinance. This general opinion may be held under various modifications; but all of them are opposed to the doctrine I have already laid down, that the Sacraments are seals of a justifying and saving grace already enjoyed by the recipient, and not intended for the conversion of sinners; and that they become means of grace only in so far as the Spirit of God, by the aid of the ordinance, calls forth the faith of the recipient, and no further.

The doctrine of the efficacy of Sacraments, directly and immediately of themselves, and not indirectly and mediately through the faith of the receiver, and through the Spirit in the receiver, is advocated in its extreme and unmodified form by the Church of Rome. According to that Church, these ordinances, as outward and material rites, become, after certain words of institution pronounced by the priest, possessed of a sacramental virtue, which is conveyed infallibly to the soul of the person who receives them, on two conditions, which are necessary to justifying and spiritual grace being really imparted. First, on the side of the priest who pronounces the words of institution, there is required, as a condition of the supernatural grace being imparted, that he have the intention to make the Sacrament and confer it; for without this, the outward matter of the ordinance would remain mere matter, and have no sacramental character or virtue. And second, on the side of the recipient of the ordinance, it is required that he be free from any of those sins which, in the language of Popery, are called ``mortal,'' and which, when contracted and not removed, would resist the operation of the sacramental virtue, and prevent his soul receiving spiritual grace. But when these two conditions are present,---when the priest intends to consecrate and dispense the ordinance, and the recipient is not barred from the reception of its virtue by mortal sin,---such is the efficacy of the Sacrament in itself, and directly, that it infallibly communicates to the partaker of it justifying and saving grace. The doctrine of the Church of Rome is very distinctly brought out in the canons of the Council of Trent, and also in her Catechism. ``If any,'' says the 11th canon concerning the Sacraments in general, ``if any shall say that there is not required in the ministers, when they make and confer the Sacraments, at least the intention of doing what the Church does, let him be accursed.'' ``If any shall say that the Sacraments of the New Law do not contain the grace of which they are the signs, or that they do not confer that grace on those who place no obstacle in the way, as if they were only outward signs of grace or justification already received by faith, and certain badges of the Christian profession, by which believers are distinguished from infidels, let him be accursed.'' ``If any shall say that grace is not conferred by the Sacraments of the New Law, ex opere operato, but that faith in the Divine promise alone avails to secure grace, let him be accursed.''\footnote{\emph{Concil. Trident. Canones et Decreta,} Sess. vii. \emph{De Sacr. in Gen.,} Can. vi. viii. xi. {[}Compare Möhler's statement of the Roman Catholic theory of the Sacraments (\emph{Symb.} 6te Aufl. pp.~253-258). It is worthy of remark, that that acute and dexterous controversialist, in his own exposition of the doctrine, passes over in utter silence the very remarkable and important element of the priest's intention, as defined at Trent,---the sole reference to it being in a quotation from Bellarmine given in a note (p.~256). Nitzsch's remarks on the significance of the point thus ignored, are worth quoting: ``The demand of an \emph{intention} on the part of the priest in order to the Sacrament being savingly effectual, or effectual at all, met opposition among the Romish theologians, both at and before the Council of Trent. At one time the danger was pleaded incurred by the baptized or the absolved, who might now so easily miss obtaining grace, or be left in uncertainty about it; at another, the much greater concession already made, that unbelief, or even mortal sin on the part of the priest, did not destroy the efficacy of the priestly act. On these grounds, the Council saw themselves compelled to restrict the demand as much as possible; it asserted itself, notwithstanding. Sess. vii. Can. 11: `Si quis dixerit in ministris, dum Sacramenta conficiunt et conferunt, non requiri intentionem saltem faciendi quod facit Ecclesia, anathema sit,' with good reason, as is easily to be seen. For if, as Can. x. decided, the private Christian cannot make or confer most of the Sacraments,---if the supernatural qualification of the priest, although a `\emph{gratia gratis data,}' and not `\emph{gratum faciens,}' still was of the essence of the sacramental dispensation,---while, on the other hand, no `\emph{bonus motus}' of the recipient was necessary for the reception of the grace, nay, in the case of private mass, no recipient at all, in the case of infant Baptism, no conscious recipient, was needed,---then, should there be an \emph{entire} absence of harmony between the mental state of the priest and the design of the transaction, there would be absolutely nothing left but the bare, mechanical, accidental, external act; and from this hardly a single believer would have expected any blessing whatsoever. The worth of the words of institution and promise, as appropriated by faith, had already been sacrificed to the worth of the `\emph{opus;}' so had the dignity of the congregation to the dignity of the priest: hence arose great perils and perplexities, if now, after all, the worth of the words should be allowed to stand without moral desert on the part of the priest. The morally indifferent supernatural qualification of the priest must therefore now receive at least a psychological quickening, and the co-operation of the mental state of the priest be thus brought in to give the requisite support to a transaction otherwise bereft of all substance and security. They resigned themselves, accordingly, to the lesser perplexity. The doctrine of the `\emph{intentio ministri}' is a reinforcement to the doctrine of the `\emph{opus operatum,}' which yields at the same time various advantages of another sort; and the latter dogma is again explained and supported by the notion of the `not interposed hindrance.'\,'' \emph{Prot. Beant.} Hamburg 1835, p.~154. Gerhard, \emph{Loci,} xviii. 31-38, ed. Preuss. tom. iv. pp.~151-158.{]}} According to this doctrine, then, Sacraments impart grace, not through the channel of the faith of the receiver, and not in dependence in any way on his spiritual act, but immediately and directly from themselves, ``\emph{ex opere operato.}'' This last expression is to be interpreted in connection with the distinction drawn by the Church of Rome between the Sacraments of the Old and New Testament Churches. The Sacraments of the Gospel Church are superior in efficacy to those under the law, in the Popish theory, because the former, or the New Testament Sacraments, work grace independently of the spiritual disposition or act of the recipient; whereas the latter, or Old Testament Sacraments, were dependent on the spiritual disposition or act of the receiver of them. The ``opus operatum'' of the New Testament Sacraments, or the virtue they have by their own act, apart from the spiritual state of the recipient, is contrasted with the ``opus operantis'' in the Old Testament Sacraments, or the virtue which they had, not in themselves, or in their own operation, but only in connection with the spiritual act of the partaker. According to the proper theory of the Church of Rome, the Sacraments of the New Testament impart grace ex opere operato, or from their own intrinsic virtue and direct act on the soul of him who receives them.\footnote{The statement that Papists hold the Sacraments to be efficacious of themselves, apart from the spiritual condition of the recipient, is often met---especially by English Romanists---with a flat denial. And on this ground. They hold that many elements are, in point of fact, present in every case in which the Sacraments are efficacious; some of these elements are connected with the state of the recipient,---such as a desire to receive the ordinance,---and others with the working of God. Thus Bellarmine objects to Calvin's stating the point in debate to be, not as to grace being conferred in the Sacraments, but only ``whether God works in them by His own proper, and, so to speak, intrinsic virtue, or whether He resigns His place to the external signs.''---\emph{Inst.} lib. iv. c. xiv. 17. (See next note.) Any Romanist, however, who has the slightest regard for the authoritative declarations of his Church, can be fixed down conclusively to this position, that, whatever other elements may, in point of fact, be present, the immediate, efficient, instrumental cause of the grace invariably conveyed to all ``qui non ponunt obicem'' is ``the outward action called a Sacrament,'' and nothing else. See Turrettin's masterly treatment of this point, \emph{Opera,} loc. xix. Qu. viii. 2-6. Cunningham, \emph{Works,} vol.~iii. pp.~124-139. Hodge, \emph{Princeton Essays and Reviews,} New York 1857, pp.~370 f., 388.}

This doctrine of the inherent power of Sacraments in themselves to impart grace, held by the Church of Rome, is also the system maintained, although with some important modifications, by another party beyond the pale of that Church, the representatives of which, at the present day, are to be found in the High Churchmen of the English Establishment. The doctrine of the High Church party in the English Establishment in regard to the Sacraments differs indeed in two important particulars from the full and unmodified development of it found in the Popish system; but in other respects it is substantially the same,---equally implying the inherent power of Sacraments to impart grace, not through the spiritual act of the recipient, but apart from and independently of it. The advocates of High Church principles in the Church of England generally---although there is a numerous and increasing section of them who in this respect approximate more nearly to Rome---generally reject the Popish doctrines,---first, of the opus operatum, and second, of the necessity for the intention of the priest in the Sacrament. They deny that the Sacraments have any immediate physical influence upon the soul, by the very act of outwardly participating in them,---such as is implied in the opus operatum of the Church of Rome; and they deny, further, that the intention of the priest to make and confer the Sacrament is a necessary condition of it, without which it could impart no grace. These two elements in the Popish theory of sacramental ordinances are rejected, generally speaking, by the High Church disciples of the English Establishment, although instances are not awanting---and they seem to be multiplying of late---of both these monstrous pretensions being, in a certain sense, maintained by them. But they agree with the Romish Church in the grand and fundamental principle which belongs to its doctrine of the Sacraments,---namely, that they communicate grace from the sacramental virtue that resides in themselves,---or, as some prefer to put it, that invariably accompanies them by Christ's appointment,\footnote{In, cum, or sub Sacramento. {[}``It is to be observed,'' says Bellarmine, ``that the dispute is not about the \emph{mode} in which Sacraments are causes of justification, \emph{i.e.} whether the effect is produced by physical or moral means; and again, if the influence be of a physical sort, whether it be by some inherent quality, or by the simple will of God; for these points do not belong to the question of faith; but only in general, whether the Sacraments are true and proper causes of justification; so that it truly follows, from a man's being baptized, that he is justified.''---\emph{Disputationes,} tom. iii. lib. ii. cap. i.{]}}---and by their own immediate influence on the soul, and not instrumentally by the operation of the Spirit of God on the worthy recipient and through the medium of his faith. This is the characteristic principle that is common both to the Popish and the High Church theories of Sacraments. Both these parties hold that there is something in or connected with the ordinance which directly and immediately does the work of grace upon the soul; and not merely indirectly and mediately through the Spirit of God working on the soul, and the faith of the soul working in return. The Church of Rome ascribes this efficacy of the ordinances to the opus operatum of the Sacraments, and the act and intention of the priest in consecrating them. The High Churchmen of the English Establishment usually reject both of these doctrines as laid down by the Council of Trent, and ascribe this efficacy of the ordinances to the deposit of spiritual grace which Christ has communicated to the Church, and connected with the Sacraments, and given them the power to impart. But the High Churchmen of Rome and the High Churchmen of England agree equally in this, that there are in the Sacraments an efficacy and power to impart grace of themselves, directly and immediately, to the soul of the recipient; and that they are not merely aids or instruments for bringing the recipient into direct and immediate communication with Christ to receive grace from Him.\footnote{{[}Goode, \emph{Nature of Christ's Presence in the Eucharist,} Lond. 1856, vol.~i. pp.~vi. 11-55. Cunningham, \emph{Works,} vol.~i. pp.~233-237.{]}}

Although both the Canons and Catechism of the Council of Trent lay down, to all appearance, expressly and undeniably the doctrine that there is a physical virtue in Sacraments, whereby they operate upon the recipient, yet there are not awanting doctors of the Romish Church who are anxious to soften down the dogma of the opus operatum, and to explain it in the sense of a moral and spiritual, and not a physical virtue, residing in the ordinance. And in this modified form of it, the Romish doctrine of the Sacraments---apart from the necessity of the priest's intention---approximates very closely to the High Church theory entertained by many in the Church of England. That theory maintains the doctrine of not a physical but a spiritual virtue deposited and residing in the Sacrament, which operates universally, not through the faith or spiritual act of the recipient, but directly and immediately through the act of participation in the outward ordinance. This, in fact, is no more than part of the general doctrine that the Church is the grand storehouse of grace to man, and not Christ Himself; and that it is by communication with the Church, and not by direct communication with Christ, that the soul is made partaker of that grace. The Sacraments, as the chief medium through which the Church communicates of its stores of spiritual blessings, are the efficient instruments for imparting grace directly to the recipient.

Now, there is one preliminary remark which, in proceeding to estimate the value and truth of such principles in regard to the Sacraments, it is necessary to bear in mind. It is not denied, but, on the contrary, strongly maintained and asserted, that the Sacraments are means of grace. To the believer who uses them aright, they are made the means of conveying spiritual blessings. In regard to this, there is no controversy between the opponents and the advocates of High Church views of the Sacraments, whether Popish or Tractarian. But the question in dispute is, whether the Sacraments become effectual, from a virtue in themselves, or in the priest that consecrates them, or only by the work or the Spirit and the faith of the recipient? That the faith of the believer is called forth and exercised in the ordinance, and that through this faith he receives grace additional to what he enjoyed before, we do not dispute, but, on the contrary, strenuously maintain. That the spiritual act of the believer in the ordinance, when in faith he gives himself to his Saviour, is met by the spiritual act of Christ in the ordinance, when in return He gives Himself and His grace to the believer, is a doctrine at all times to be asserted and vindicated. That the faith of the recipient, in the act of committing and engaging himself to Christ, through means of the ordinance, is a faith unto which Christ is given in return, we would constantly affirm; and in this sense, and in this way, the Sacraments become means or channels or instruments whereby grace is given and conveyed. But they are no means of grace except through the faith of the recipient, and in consequence of his own spiritual state and act. There is no inherent power in the ordinance itself to confer blessing, apart from the faith of the participator, and except through the channel of that faith. There is no deposit of power---whether, with the Church of Rome, we deem it physical and ex opere operato, or whether, with Tractarians and High Churchmen, we call it spiritual---in the Sacraments themselves to influence the mind of him who receives them. They have no virtue of themselves, apart from the work of Christ through His Spirit on the one side, and the spiritual act of the recipient through his faith on the other side. In the language of Amesius, in his admirable reply to Bellarmine, Sacraments have no power ``efficere gratiam immediate, sed mediante Spiritu Dei et fide.''\footnote{Amesius, \emph{Bellarm. Enerv.} Amsterdam 1658, tom. iii. lib. i. cap. v. p.~22. ``In the following sentences of the \emph{Declaration of Thorn,}'' observes Nitzsch, ``all Protestants agree: `Sacraments are outward and visible signs, seals, and testimonials of the Divine will, instituted by God Himself, by the combination of Word and element, in order to seal and exhibit, through means of these signs, the invisible grace which is promised in the Word of the covenant. It is obvious that we by no means make them bare signs, empty and inefficacious, or mere badges of outward profession, since, besides their mystical significance, according to the Divine institution, we attribute to the Sacraments a sure sealing of God's promises, and at the same time a true and infallible exhibition of the things promised, in the way suited and proper to them, to be received by a living faith.'\,'' (Niemeyer, p.~680.) \emph{Prot. Beant. der Symb. Möhler's,} Hamburg 1835, p.~175. Bruce, \emph{Serm. on the Sacr.} Wodrow Soc. ed. Edin. 1843, p.~10 f. Calvin, \emph{Antidote to Council of Trent,} Sess. vii. Can. ii. iv.-vi. \emph{Tracts,} vol.~iii. Calvin Transl. Soc. Edin. 1851, pp.~172-175.{]}}

Has the Church, then, ordinances for its administration and use which, either by the original appointment of Christ, or by deposit of grace from Christ, have in themselves virtue to impart spiritual blessing through the administration of them alone? Or has the Church ordinances for its administration and use which have no virtue in themselves to communicate grace, except in connection with the faith of the receiver, and the blessing imparted by the Spirit? Are the Sacraments of the New Testament themselves a quickening power in the soul, apart from the faith or spiritual act of the participator,---the original deposit of grace committed to them being still retained, and still communicable through their administration, and that alone? Or are these Sacraments effectual to impart grace only in connection with the faith and spiritual disposition of the recipient,---there being necessary to their efficacy, both the act of the believer, in the use of them, giving himself to Christ, and the act of Christ, through the same ordinance, giving Himself to the believer. It matters little whether, as with the Popish Church, the Sacraments are invested with a physical virtue, in consequence of which they impart grace; or whether, as with the High Churchmen of other denominations, they are invested with a spiritual virtue in consequence of which they impart grace,---if in both cases the grace is given by the Sacrament itself, and not given through the Spirit and the faith in the heart of the recipient. It matters little whether a physical or a spiritual explanation is given of sacramental efficacy, if it be efficacy exerted apart from Christ in the ordinance giving Himself to the believer, and experienced apart from the believer in the ordinance giving himself to Christ. Whatever be the efficacy and virtue, physical or moral, if it is independent of and separate from the faith of the recipient covenanting in the ordinance with Christ, and the act in answer to that faith of Christ covenanting with the recipient, it is not the sacramental grace which the Scripture recognises. It becomes, when thus separated and drawn apart, a mere charm, a trick of magic, whether physical or spiritual, utterly unknown to the Gospel economy. Let us endeavour to apply to this theory those tests which may serve to try its merits and its truth. There are four different tests by which we may try the merits of this sacramental theory, whether held in its extreme form by Papists, or in its more modified form by High Churchmen of other communions.

1st, Tested by Scripture, which constitutes the rule for the exercise of Church power, there is no warrant for asserting that there is an inherent and independent virtue in Sacraments to impart justifying or saving grace.

The truth of this general proposition may be established by a very wide and ample deduction of evidence from Scripture. It is impossible for us to do more than advert to the leading heads of proof in connection with this question. In the first place, those multiplied and various declarations of Scripture, which state that we are justified by faith alone without works on our part, very distinctly prove that the Sacraments cannot have an independent and inherent power in themselves of conveying justifying and saving grace. Such passages expressly assert that faith is the immediate instrumental cause of justification. They are inconsistent, therefore, with the theory that the Sacraments directly and immediately of themselves impart grace, although they are quite consistent with the doctrine that the Sacraments indirectly, and through the faith of the worthy receiver, may impart grace. In the second place, the doctrine that the Sacraments have an inherent virtue to confer grace, is opposed to the whole tenor of Scripture, which sets forth Christ as the one and the immediate object of faith and hope to the believer, in the matter of his justification and salvation. The Word of God, from its commencement to its close, clearly and constantly and invariably points to Christ, and to nothing but Christ, as the only source to which a sinner must look for forgiveness and acceptance with God. The theory of the Sacraments held by High Churchmen presents another and a different object for his faith, and teaches him to rest in an outward observance as sufficient. It is part of that most destructive system which places the Church and the ordinances of the Church between the sinner and his Saviour. In the third place, the very express testimony of the Apostle Paul, in regard to the insufficiency of the Sacraments under the Old Testament Church to communicate grace of themselves, is an argument equally effectual to show that the New Testament Sacraments are insufficient likewise. Abraham was not justified by circumcision, but by the faith of which his circumcision was the seal.\footnote{Rom. ii. 25-29, iii. 20, 30, iv. 3-11; Heb. ix. 11 f., x. 1-11. {[}Comp. the Apology for the Confession of Augsburg, vii. 18, p.~203, in Hase, \emph{Libri Symbolici Eccles. Evang.} Lipsiae 1827.{]}} \emph{In the fourth place,} the statements of Scripture which at first sight might be construed as if they ascribed a gracious influence to the Sacraments of the New Testament in themselves, and which seem to connect saving benefits with the observance of them, are not stronger or more numerous, but less so, than those which ascribe justifying and saving blessings to the ordinance of the Word, or truth received by the reader or hearer of it. We know that the Word or the truth justifies, not of itself, but through the faith of him that receives it; and that, apart from this faith, it has no virtue or power of a gracious kind at all. In the same manner, Sacraments impart grace, not of themselves, but through the faith of those who receive them; and, apart from that faith, they have no life or blessing whatsoever. In the fifth place, the theory of an inherent virtue or power in the administration of the outward ordinance is utterly opposed to those numerous passages of Scripture which assert that the power of the Gospel is altogether of a spiritual kind, and is in no respect akin to a mere external and material influence, as if such could impart a supernatural grace. It is ``not meat and drink, but righteousness, and peace, and joy in the Holy Ghost.'' And instead of pointing to any outward source of power or efficacy, and exclaiming, ``Lo here, or Lo there!'' the Christian has been taught to think that ``the kingdom of God'' has its source and presence ``within him.''\footnote{Rom. xiv. 17; Luke xvii. 21.} The theory which ascribes to the Sacraments an infallible virtue which, unless counteracted by some obstacle, such as infidelity or open vice, must operate to impart grace, is inconsistent with those numerous statements of Scripture which represent the Gospel as a spiritual power, adapted to the spiritual nature of man.\footnote{Gillespie, \emph{Aaron's Rod Blossoming,} B. iii. chap. xii.-xiv.}

In estimating the bearing of Scripture testimony on this question, there is one consideration of a general kind which it is of great importance to the argument to bear in mind. In every theory of the Sacraments that can be held,---from the lowest to the highest, from the Socinian up to the Popish,---the Sacraments are regarded as at least signs of spiritual things, representing and exhibiting the blessing in outward resemblance. The union thus established, according to any theory that can be held of them, between the sign and the thing signified by it, has introduced into Scripture a kind of phraseology which at first sight appears to give some sanction to the High Church system in regard to sacramental ordinances. There is often an exchange of names between the sign and the thing signified in Scripture, in consequence of which what may be predicated of the one is often asserted of the other, and vice versâ. This usage of language, so frequently exemplified in Scripture in connection with this matter, is a usage found commonly in other writings and in regard to other matters, and gives rise to no sort of misapprehension in our interpretation of it. It is the great foundation indeed of all figurative language.\footnote{{[}``Omnia significantia videntur quodam modo earum rerum quas significant sustinere personas: sicut dictum est ab Apostolo, `Petra erat Christus,' quoniam petra illa, de qua hoc dictum est, significabat utique Christum.''---Aug. \emph{De Civitate Dei,} lib. xviii. cap. 48.{]}} Thus, when Christ is said to be ``the Passover sacrificed for us,'' there is an exchange of this kind, in which the name of the sign is given to the thing signified; and when Christ says of the bread, ``This is my body,'' there is an exchange in the opposite way, and the name of the thing signified is attributed to the sign. And in perfect accordance with this usage of language, there are several passages in Scripture in which the mere outward observance in the case of the New Testament Sacraments, the external sign, has a virtue attributed to it which in reality belongs, not to the sign, but to the grace represented in the observance, or to the thing signified. Thus, for example, ``Baptism'' is said in one passage ``to save us;'' although, from the further explanation contained in the passage itself, it is plain that it is not the outward sign but the thing signified that is spoken of under the name of the sign; for the apostle adds immediately, ``not the putting away of the filth of the flesh, but the answer of a good conscience towards God.''\footnote{1 Pet. iii. 21.} In the same manner the Apostle Paul speaks of ``the cup of blessing'' as ``the communion of the blood of Christ,''\footnote{1 Cor. x. 16.}---language in which that is predicated of the sign which is truly predicated only of the thing signified. In short, the sacramental union between the outward sign and the inward grace gives occasion to not a few examples in Scripture in which what is true of the one only, or the inward grace, is attributed to the other, or the outward sign. Almost the whole plausibility of the argument from Scripture in favour of the High Church theory of the Sacraments comes from this source; and it is completely removed when the familiar canon of criticism, applicable to Scripture in common with other writings, is attended to,---namely, that what truly belongs to the thing signified is often predicated figuratively of the sign, and so ought to be interpreted and understood.\footnote{{[}Westminst. Conf. chap. xxvii. 2, xxix. 5. Goode, \emph{Nature of Christ's Presence in the Eucharist,} Lond. 1856, vol.~i. pp.~241-250, 598.{]}}

2d, The theory of an inherent power, physical or spiritual, in the Sacraments, is inconsistent with the supreme authority of Christ, from whom all Church power is derived.

The doctrine that would deposit in sacramental ordinances a grace communicable to the participator, apart from his communion with Christ, directly and immediately, is inconsistent with the office and right of Christ to hold in His own hand all blessing, and to dispense from His own hand, not mediately through another, but at once from Himself, the grace which His people receive. Such a theory takes the administration of grace out of the hands of Christ, ever present to dispense it, and transfers it to a priest standing in His room. There can be no participation in heavenly blessing except what comes from direct communication with Christ on the part of the soul that receives it; and it is a dishonour to Him, who is the ever-living and ever-present administrator of all grace to His people, to put the mute and conscious ordinance in the place of Christ, and to transfer the dependence of the soul for spiritual blessing from the Divine Head in heaven to the outward ministry of Sacraments on earth. That Christ might by His original appointment have made the Sacraments the receptacle of a physical influence, fitted and able to work a supernatural blessing on the soul, it would perhaps be presumptuous to deny. That Christ might at the first institution of the ordinances have made them a reservoir or storehouse of grace enough for all ages of the Church, and imparted to them a spiritual blessing out of which every subsequent generation of His people might draw their supply, we need not be anxious to dispute. Or that Christ, without communicating at the beginning to Sacraments a store either of physical or spiritual grace sufficient for all generations, might have tied Himself up to the indiscriminate and invariable communication of His Spirit along with the administration of outward Sacraments, and bound Himself down, without any choice or discretion, to link spiritual grace to material rites, apart from the faith of the person observing them,---this, too, is perhaps a possible imagination. But had Christ, as the Head of ordinances in His Church, done either the one or the other of these things, He must to that extent have divested Himself of His office as Mediator, or resigned the exercise of it; He must in so far have abdicated His functions as the sole and living and ever-present administrator of grace to His Church; and been shut out from that exclusive and supreme agency which He maintains as the dispenser as well as author of every blessing by which the soul is to be saved.

3d, The theory of the Sacraments which ascribes to them an independent virtue or power, is inconsistent with the spiritual liberties of Christ's people.

Such a system brings the soul itself into bondage. It keeps the spirit, which Christ has Himself redeemed, waiting upon man for the communication of the blessings of its redemption; it makes the soul which Christ has ransomed dependent for its freedom on the ministry of a fellow-creature. There cannot be a worse or more abject thraldom than that which subordinates the flock of the Saviour's purchase to any one but Himself, and causes them to hang upon the intention entertained or not entertained by a priest for the enjoyment or forfeiture of spiritual blessing. But even apart from the monstrous doctrine of the Romish Church as to the intention of the priest being necessary to the efficacy of the ordinance, the sacramental theory we have been considering, whether Popish or Tractarian, is inconsistent with the spiritual freedom of those whom Christ has redeemed. That freedom consists in subjection to and dependence on Christ, and none but Christ,---in being emancipated from all dependence on any other except their Saviour,---in being kept waiting, not at the footstool of man for saving blessings, but at the footstool of Christ,---and in being taught to look for all the grace they need day by day, not to the ministry of man's hand, but to the hand of Christ. Spiritual freedom for the believer is bound up with a dependence on Christ immediately and directly, and on Him alone, for every blessing that he needs.\footnote{{[}Goode, \emph{Letter to a Lay Friend,} Lond. 1845, pp.~18-24. Litton, \emph{Church of Christ,} Lond. 1851, pp.~11-13, 202-232, 240 f.{]}}

4th, The sacramental theory we have been considering is inconsistent with the spirituality of the Church, and of the power exercised by the Church for the spiritual good of men.

When, according to that theory, the Sacraments become the instruments of justification and the source of faith, instead of the seal of a justification already possessed, and the exercise and aid of a faith already in existence,---when they are made to come between the soul, in its approach to Christ, and Christ Himself, and communion in the external ordinance is substituted for the fellowship of the Spirit, it is a fatal evidence that the Church, which so teaches and so practises her teaching, although she has ``begun in the Spirit,'' has ``sought to be made perfect by the flesh.''\footnote{Gal. iii. 3.} If the external ordinance be made to occupy that place which belongs to the Spirit, and participation in the ordinance be the substitute for faith, the sacramental theory thus reduced to practice will be but the commencement of worse and deeper degradation. It is but the beginning of a course which, consistently followed, must lead to a religion of form and self-righteousness, of sense and sensuous observances, of carnal ordinances and a ceremonial holiness, of outward satisfaction and penances and merit. There will be the priest and the bloodless but efficacious sacrifice, grace conferred by the tricks of a physical or spiritual magic, a religion that manifests itself outwardly and not inwardly, the holiness of houses, and altars, and sacred wood and stone, but not the holiness of the Spirit; the atonement of Sacraments and penances and creature merits, but not the atonement of the Saviour received by faith; a righteousness of bodily discipline and fleshly mortification, but not the righteousness of God imputed to the believer; a justification made out of pains and merits, of sufferings and works, but not a justification freely given by Divine grace and freely accepted by faith; an outward baptism to regenerate the sinner with water at first,---the food of the communion table, made flesh and blood by the consecration of a priest, to sustain the life so begun, and the anointing with oil at last to prepare the soul for the burial. Such are the inevitable fruits of the sacramental theory, consistently carried out in the Church of Christ, making the very temple of God to be the habitation of every carnal and unclean thing.\footnote{Bellarm. \emph{Disputationes,} tom. iii. lib. ii. cap. i.-xxii. Perrone, \emph{Praelectiones Theologicae,} Parisiis 1842, tom. ii. pp.~5-66. Amesius, \emph{Bellarm. Enerv.} tom, iii. lib. i. cap. i.-v. Turrettin, \emph{Opera,} tom. iii. loc. xix. qu. i.-ix. Cunningham, \emph{Works,} vol.~iii. pp, 121-133. {[}Bruce, \emph{Sermons on the Sacraments,} Wodrow Soc. ed. Edin. 1843, pp.~11-33. Newman, \emph{Lectures on Justification,} pp.~316, etc.; Tract No.~90, 2d ed. p.~13. Wilberforce, \emph{Doctrine of the Holy Eucharist,} 3d ed. Lond. 1854, pp.~17-38, 97-130. Goode, \emph{Doctrine of the Church of England as to the Effects of Baptism in the case of Infants,} 2d ed. Lond. 1850, pp.~3-10. \emph{Vind. of the `Defence of the XXXIX.Art.'} etc., 2d ed. p.~38 f. \emph{Unpublished Letter of Martyr to Bullinger,} Lond. 1850, pp.~11-13. Martensen, \emph{Dogmatik,} 4te Aufl. Kiel 1858, pp.~361, 364. Matthes, \emph{Comparative Symbolik,} Leipzig 1854, pp.~492-510. Thomasius, \emph{Dogmatik,} 3ter Theil, 2te Abth. Erlangen 1861, pp.~113-135.{]}}

\hypertarget{the-sacrament-of-baptism}{%
\chapter{The Sacrament Of Baptism}\label{the-sacrament-of-baptism}}

\hypertarget{nature-of-the-ordinance}{%
\section{Nature Of The Ordinance}\label{nature-of-the-ordinance}}

PASSING now from the doctrine of the Sacraments in general, or viewed in respect of what belongs to them in common, I proceed to consider them more in detail and individually; and for this purpose I commence with the Sacrament of Baptism, as the initiatory rite. Upon what grounds are we justified in attributing to Baptism the name and character of a Sacrament? What is the nature of the ordinance, the place which it occupies, and the office it is intended to serve in the Christian Church? The general principles which we have already laid down in regard to Sacraments as such, when applied more particularly to Baptism, will enable us to bring out distinctly the character, authority, and meaning of the ordinance. There were four elements which we found to enter into the idea of a Sacrament. Let us proceed to apply these to the ordinance of Baptism, in order that we may ascertain its true nature and import. And in doing so, we shall have an opportunity, at the same time, of noticing some of the opinions in regard to Baptism which we hold to be unscriptural and erroneous.

\hypertarget{the-first-characteristic-of-a-sacrament-is-that-it-must-be-a-positive-institution-of-christ-in-his-church-and-this-mark-applies-to-baptism.}{%
\subsection{The first characteristic of a Sacrament is, that it must be a positive institution of Christ in His Church; and this mark applies to Baptism.}\label{the-first-characteristic-of-a-sacrament-is-that-it-must-be-a-positive-institution-of-christ-in-his-church-and-this-mark-applies-to-baptism.}}

The doctrine of the Quakers is opposed to this first position. They contend that Baptism, and the Lord's Supper also, were Jewish practices, neither suited to the Gospel economy nor appointed for the Gospel Church, but destined to be done away with under the dispensation of the Spirit. Now, in reference to Baptism, it cannot be doubted that it was a Jewish observance before it became a Christian one, and that it was administered by the Jews to proselytes joining them from among the Gentiles, previously to the time when it was adopted by our Lord as one of the Sacraments of His Church. This is sufficiently attested by the statements of Jewish writers; it may be inferred, indeed, from the narratives of the Evangelists. Baptism, as an initiatory rite and token of discipleship, connected with a sect or school of religion, was familiarly known among the Jews; and it is on the ground of their previous acquaintance with and practice of it amongst themselves, that we can understand the question addressed to John the Baptist: ``Why baptizest thou then, if thou be not that Christ, nor Elias, neither that Prophet?'' Had John been any of those personages come into the country as a teacher or founder of a new school of religion, the Jews would have felt no surprise, and expressed no objection to his practice of baptizing with water; and it was only because he denied that he was either Christ or Elias, that they were led to demand the authority by which he baptized. Although, then, there is no mention of any such ordinance in the law of Moses, yet there seems to be no doubt that it was a ceremony that had found its way into the practice of the Jews.2 But we are not on this account to imagine that Christian Baptism was one of those temporary ordinances destined to be done away with, or that it is not a positive institution of Christ in His Church. During His own personal ministry on earth, we are given to understand that, acting on our Lord's direct authority, His Apostles adopted the rite, and administered it to the Jews who professed their desire to become Christ's disciples. Side by side with the commission to preach the Gospel given to the Apostles, when the Church was set up by our Lord after His own resurrection, we find the command to baptize those whom they taught; and the ordinances of the Word and of Baptism are spoken of in terms significant equally of the authority and standing obligation of both. ``Go ye therefore, and make disciples of all nations, baptizing them in the name of the Father, and of the Son, and of the Holy Ghost; teaching them to observe all things whatsoever I have commanded you: and, lo, I am with you alway, even unto the end of the world.''

The natural and indeed unavoidable interpretation of the apostolic commission seems to establish these two things: first, that a literal Baptism, or washing with water, was to accompany the discipleship brought about by the preaching of the Apostles; and second, that both the ordinance of Baptism and that of preaching were to be continued unto the end of the world. Added to this, we have the evidence for the Divine authority and permanent obligation of Baptism in the Church of Christ, from the unvarying practice of the Apostles in regard to their converts, whether Jewish or Gentile, down to the latest period in the history of the Church to which the inspired narrative refers. Such considerations as these go to prove that Baptism was not a mere Jewish practice, suffered for a time in the Christian Church, and destined to be cast off with other Jewish customs and observances. On the contrary, the positive appointment of our Lord expressed in the commission He addressed to the Apostles as founders of the Christian society,---the apostolic example itself as regards Baptism equally of Gentile and Jewish converts,---and the entire absence of any intimation, either express or implied, that the practice was only temporary and designed to be discontinued, go undeniably to prove that Christian Baptism is a permanent institution of Christ in His Church.

\hypertarget{another-characteristic-of-a-sacrament-is-that-it-be-an-external-and-sensible-sign-of-an-internal-gracea-spiritual-truth-embodied-in-an-outward-action-and-this-mark-is-applicable-to-christian-baptism.}{%
\subsection{Another characteristic of a Sacrament is, that it be an external and sensible sign of an internal grace,---a spiritual truth embodied in an outward action; and this mark is applicable to Christian Baptism.}\label{another-characteristic-of-a-sacrament-is-that-it-be-an-external-and-sensible-sign-of-an-internal-gracea-spiritual-truth-embodied-in-an-outward-action-and-this-mark-is-applicable-to-christian-baptism.}}

That Baptism is symbolical of unseen and spiritual blessings, is admitted by all parties who hold the ordinance itself to be an appointment of Christ, whatever theory they may entertain as to its sacramental character or virtue. Adopted as it was by Christ from Jewish customs and practices, it could hardly fail, indeed, at its original institution in the Christian Church, to appear to those who used it to be of a symbolical character. They had been accustomed to the washings and sprinklings practised under the law as symbolical observances, expressive of the removal of ceremonial uncleanness, and of such a ceremonial purification as secured acceptance with God,---at least outwardly. And when Baptism was appointed by our Lord, the washing with water included in it must have been interpreted, in accordance with the previous use and meaning of the Jewish observances, as a purification, or a putting away of defilement of sin, so that the person baptized was accounted clean, and fitted for acceptance with God. Hence the language of Scripture everywhere in connection with Baptism conveys the idea of its being a symbolical ordinance like the ancient washings and sprinklings customary among the Jews, and indeed among other nations, as expressive of religious purification or cleansing. The body washed with pure water was an emblem of the soul purified and cleansed through the blood and Spirit of Christ. The ``Baptism for the remission of sin'' was expressive of the cleansing by which sin is removed. The action by which water was applied by the administrator to the person, was representative of the application of the blood of Christ to the guilt of the soul. The action by which the washing of Baptism was submitted to by the recipient, was expressive of his passing under the washing of regeneration and the renewing of the Holy Ghost. And the distinguishing practice in Christian Baptism, that the person who received the ordinance was baptized ``into the name of the Father, and the Son, and the Holy Ghost,'' was symbolical of his dedicating himself to the Father, through his justification by the blood of the Son, and his sanctification by the grace of the Spirit.

There was the twofold representation, exhibited in the ordinance of Baptism, of Christ giving Himself to the believer in the two great initial blessings of the covenant,---justification and sanctification,---and of the believer dedicating himself to Christ as one of His justified and sanctified people. Christ united to the believer, and the believer united to Christ, in consequence of the removal both of the guilt and pollution of sin which had separated between them, is the great lesson exhibited in the ordinance of Baptism as a symbol. Hence Baptism, rather than the Lord's Supper, forms the great initiatory rite of the Church. The former ordinance is more especially fitted symbolically to represent the union of the believer to Christ; the latter to set forth the communion of the believer with Christ. Baptism meets us at our entrance into the Church, and by the purification from the guilt and defilement of sin, which it more particularly represents, it exhibits us as entering into union with a Saviour in the only way in which that union can be effected,---in the way, namely, of free justification by the blood of Christ sprinkled upon the soul, and full sanctification by the Spirit of Christ cleansing and renewing our nature. In regard to this office which we assign to Baptism, of being a sign of the spiritual blessings of the covenant by which the believer is united to Christ, all parties who hold Baptism to be an ordinance of Christ at all, agree, whatever additional views they may hold as to its sacramental character or virtue.

\hypertarget{another-characteristic-of-a-sacrament-as-we-have-already-seen-is-that-it-is-a-seal-of-a-federal-transaction-between-two-parties-in-the-ordinance-and-this-third-mark-also-belongs-to-christian-baptism.}{%
\subsection{Another characteristic of a Sacrament, as we have already seen, is, that it is a seal of a federal transaction between two parties in the ordinance; and this third mark also belongs to Christian Baptism.}\label{another-characteristic-of-a-sacrament-as-we-have-already-seen-is-that-it-is-a-seal-of-a-federal-transaction-between-two-parties-in-the-ordinance-and-this-third-mark-also-belongs-to-christian-baptism.}}

It is more than a sign of spiritual blessings; it is a visible seal and voucher of these to those who rightly partake of the ordinance. At this point the theory of Baptism laid down in the standards of our Church differs from the views held in regard to it by Socinians, and by many of the English Independents. They contend that Baptism is a symbol, and nothing more than a symbol, of spiritual blessings. We maintain that the statements of Scripture warrant us in asserting that, in addition to its being a symbol, it is also a seal of the covenant entered into between Christ and the believer through the ordinance. That in the administration and participation of Baptism there is a federal transaction between Christ and the believer who rightly receives it, and that the outward ordinance is a seal of the covenant engagement, may be established by abundant evidence from Scripture.

1st, There are a number of statements of Scripture connected with the ordinance which cannot be understood except upon the supposition that Baptism is not only a sign, but also a seal of a covenant transaction between Christ and the believer. The very words of the institution seem to point to this. Baptism ``into the name (εἰς το ὀνομα) of the Father, and the Son, and the Holy Ghost'' means more than Baptism by their authority, or an expression of our submission to them. It plainly implies, on the part of the baptized person, an act of dedication of himself to the Three Persons of the blessed Godhead, under the separate characters which they bear in the work of redemption,---an act of engagement by the recipient of the ordinance unto the Father, through the Son, and by the Spirit; or, in other words, a dedication of himself to God through the medium of justification and sanctification. In exact accordance with this view, we find in Scripture that Baptism is connected with ``remission of sins,'' obtained through Christ, and with ``the washing of regeneration,'' performed by the Spirit,---expressions which go much farther than merely to represent the ordinance as symbolical of these blessings, and which appear to imply that there is an intimate connection between the right reception of Baptism and the privilege of forgiveness of sins through the blood of Christ, and of sanctification of our nature by the Spirit. What that sort of connection is which is more than a mere sign to represent, and less than an outward charm to impart these blessings, is illustrated by the Apostle Paul in a remarkable passage of his Epistle to the Romans: ``Know ye not,'' says the Apostle, ``that so many of us as were baptized into Jesus Christ were baptized into His death? Therefore we are buried with Him by baptism into death; that like as Christ was raised up from the dead by the glory of the Father, even so we also should walk in newness of life. For if we have been planted together in the likeness of His death, we shall be also in the likeness of His resurrection.'' Of course in this passage the Apostle must be held as referring to the Baptism of a believer, in whose case it was a spiritual act of faith embodying itself in the outward ordinance. There are two things which seem plainly enough to be included in this remarkable statement. In the first place, the immersion in water of the persons of those who are baptized is set forth as their burial with Christ in His grave because of sin; and their being raised again out of the water is their resurrection with Christ in His rising again from the dead because of their justification. Their death with Christ was their bearing the penalty of sin, and their resurrection with Christ was their being freed from it, or justified. And in the second place, their burial in water, when dying with Christ, was the washing away of the corruptness of the old man beneath the water; and their coming forth from the water in the image of His resurrection was their leaving behind them the old man with his sins, and emerging into newness of life. Their immersion beneath the water, and their emerging again, were the putting off the corruption of nature and rising again into holiness, or their sanctification. All this seems to be implied in this statement of the Apostle in regard to a believer's Baptism; and it cannot be doubted that, in accordance with many other passages of Scripture, it makes Baptism in the case of a believer far more than a sign of the initial blessings of justification and regeneration. The Apostle undoubtedly represents the act as a federal one, in which the believer gives himself to God in the way that God has appointed, through faith in Christ for pardon, and through submission to the Spirit for regeneration; and in which these blessings are communicated and confirmed to him. Such statements of Scripture seem to bear out the assertion, that in the Baptism of a believer there is a federal transaction, and that the outward ordinance is the seal of the spiritual covenant.

2d, The same conclusion, that Baptism is not only a sign but also a seal of the covenant, may be supported by the consideration, that Baptism has come in the room of the Old Testament Sacrament of circumcision. That the ordinance of Baptism under the New Testament has taken the place of circumcision in the ancient Church, is apparent from the statements of the Apostle Paul in his Epistle to the Colossians, in which he argues against the necessity of circumcision under the Gospel, on the ground that Baptism was all to believers now that circumcision had been to believers in former times; and where he actually calls Baptism by the name of ``the circumcision of Christ.'' ``In whom also ye are circumcised with the circumcision made without hands, in putting off the body of the sins of the flesh by the circumcision of Christ: buried with Him in Baptism, wherein also ye are risen with Him, through the faith of the operation of God, who hath raised Him from the dead.'' This assertion, that Baptism is now the circumcision of the Christian Church, leads very directly to the inference that we must regard Baptism as being as much a seal of the covenant of grace, as circumcision was a seal of the Abrahamic covenant; and it goes very clearly to establish the position, that Baptism is far more than the simple symbolical institution which many Independents would make it,---that it has more in it than the character of a mere empty sign; that there belongs to it the grand characteristic of a sacramental ordinance, namely, the character of a seal, confirming and attesting a federal transaction between God and the believer.

\hypertarget{another-characteristic-of-a-sacrament-is-that-it-is-a-means-of-grace-and-this-fourth-mark-like-the-former-ones-belongs-to-christian-baptism.}{%
\subsection{Another characteristic of a Sacrament is, that it is a means of grace; and this fourth mark, like the former ones, belongs to Christian Baptism.}\label{another-characteristic-of-a-sacrament-is-that-it-is-a-means-of-grace-and-this-fourth-mark-like-the-former-ones-belongs-to-christian-baptism.}}

Baptism is a means for confirming the faith of the believer, and adding to the grace which he possessed before. It is not intended for the benefit or conversion of unconverted men; it is not designed or fitted to impart justification or spiritual grace to those who were previously strangers to these; but it is made a means of grace by the Spirit to those who are believers already, and fitted and intended to promote their spiritual good. I do not at present speak of the case of infants baptized, or of the benefits which they may be supposed to receive from the administration of the ordinance. Their case, as peculiar and exceptional, I shall reserve for separate and more detailed consideration. But, putting aside the case of infant Baptism for the present, the position that l lay down is, that Baptism is a means of grace fitted and blessed by God for the spiritual good of the believer. And that it is so, the considerations already stated in regard to the nature of the ordinance, if they are correct and scriptural, will sufficiently enable us to understand. If the act of the adult believer in receiving Baptism be an act of making or renewing his covenant with God through the ordinance,---if his part of the transaction be the embodiment in outward sign of the spiritual act whereby he dedicates himself to Christ,---and if Christ's part of the transaction be the giving of Himself and His grace to the believer in return, then it is plain that the ordinance, so understood, must be a divinely instituted means of grace to the parties who rightly partake of it. Christ given to the believer in the Sacrament is not less precious and blessed, but more so, than Christ given to the believer in the Word; and for this reason, that in the Sacrament Christ is not only in the Word, but in the sign also. In both cases, it is, however, only in connection with the faith of the believer that the blessing is received and enjoyed; and apart from that faith, there is no blessing either in Word or Sacrament. Christ in the Word, received into the soul by faith, is the source of saving grace to the soul. Christ in the Sacrament, received into the soul by faith, is not less, but more, a blessing likewise. But in neither case can the grace and blessing be enjoyed except in connection with the exercise of faith on the part of the hearer or receiver. There is no promise connected with Word or Sacrament over and above the promise that ``the just shall live by faith.'' It is only in connection with faith, indeed, that grace can be imparted in a manner consistent with the nature of man as a moral and intelligent being, and without a subversion of its ordinary laws. The case of infants is an exceptional case, to be dealt with apart, and by itself. But in the case of adults, the communication of supernatural grace, whether through Word or Baptism, must be in connection with, and not apart from, the exercise of their own spiritual and intelligent nature, and in connection with that act of the spiritual nature which we call faith. Baptism is no exception to the ordinary principle that represents all the blessings of God's salvation as associated with faith on the part of the receiver of them. It becomes a means of grace in connection with the faith of the believer, which it calls into life and exercise.

The views now stated are of course opposed to the doctrine of what has been called ``baptismal regeneration,'' whether held by Romanists or Romanizing Protestants. The Church of Rome considers Baptism, like the other Sacraments, to be a means of imparting grace ex opere operato, and to carry with it the virtue of so applying to the person baptized, whether infant or adult, the merits of Christ, as that both original and actual transgression are completely removed by the administration of it, in every case, apart altogether from the faith of the recipient. The authorized formularies of the Church of England seem to maintain the doctrine of baptismal regeneration in a sense at least approximating to that of the Church of Rome. The Thirty-nine Articles, indeed, give no countenance to such a theory; but both her Liturgy and her Catechism appear to speak differently on the subject; and the doctrine, under various modifications, is held and asserted by a large number of her ablest divines. It is extremely difficult, in investigating this question, to ascertain the exact sense in which regeneration is understood to be imparted through the ordinance of Baptism, or the precise nature and amount of change which, according to the advocates of this doctrine, actually takes place on the person baptized. In some instances, I believe that the doctrine of baptismal regeneration is held in words, whilst it is not held in reality; the advantage conferred by Baptism on all equally and indiscriminately being nothing more than admission to the outward privileges of the visible Church, in consequence of the reception of it. But although, in the case of a few, the doctrine, as held by them, may be regarded as more nominal than real, yet it cannot be doubted that very many in the Church of England approximate, on this question, more or less closely to the views asserted in the standards of the Church of Rome.

There are at least three different modifications of the doctrine of baptismal regeneration held by divines of the Church of England, which can be readily enough distinguished from each other. First, there is one party who assert that Baptism, by the administration of it, gives the person baptized a place within the covenant of grace, in such a sense that he has a right to all its outward privileges and means of grace, and by a diligent and right use of them, may secure to himself salvation. This is the lowest view of the efficacy of Baptism held by those who assert the doctrine of baptismal regeneration, and amounts apparently to this, that Baptism is necessary in order to the salvability of a man,---all unbaptized persons having no right to the privileges of the covenant, and being left to ``the uncovenanted mercies of God.'' In answer to such a theory, it is enough to assert, with the Word of God, that the Gospel is free to all; that all, without exception of class or character, are invited to avail themselves of it; and that ``the free gift unto justification of life'' is not restricted to any limited number of men, baptized or unbaptized, but is co-extensive in its promises and invitations with ``the judgment that has come upon all unto condemnation.'' Second, there is another party who assert that Baptism conveys to the soul, by the administration of it, regenerating grace---a true spiritual life; which may continue with the baptized person, so as to avail at last to his everlasting salvation, but which may also be forfeited in after years by means of sin. This second form of the doctrine of baptismal regeneration proceeds upon an alleged distinction---held apparently by Augustine, and after him maintained by many Lutheran divines---between those who are predestinated unto life, and those who are regenerated. It is affirmed that the two classes do not coincide, and that regeneration, though once imparted to the soul, may be subsequently lost. Third, there is another party who admit that Baptism imparts saving grace and regeneration to the soul, which under no circumstance can be entirely forfeited, but which entitle the person baptized to everlasting life.

These three different forms of the theory of baptismal regeneration it is not necessary to reply to separately. The only plausible arguments which can be brought in defence of such a doctrine are derived from a few passages of Scripture which apparently, at first sight, connect the inward and spiritual grace with the outward action in Baptism which is its sign. These passages it is not difficult to explain by the help of the canon of interpretation, to which I formerly had occasion to refer, founded on the practice of Scripture, and the practice of every other book, of predicating of the sign figuratively what can only be truly and literally predicated of the thing signified. The sacramental relation between Baptism and regeneration, which it represents, easily explains the application to Baptism, figuratively, of language that belongs literally to regeneration. And while this principle, rightly understood and applied, is sufficient to explain the statements of Scripture that apparently, at first sight, give countenance to the doctrine of baptismal regeneration, the whole tenor of the Word of God clearly and decisively contradicts the theory. It is inconsistent with the fundamental principle which regulates the matter of a sinner's salvation,---the principle that he is saved and lives by faith; and that it is by faith, and not through any other channel, that he receives from God all that is necessary to his present and his everlasting well-being.

\hypertarget{the-subjects-of-baptism-as-regards-adults}{%
\section{The Subjects Of Baptism As Regards Adults}\label{the-subjects-of-baptism-as-regards-adults}}

Having discussed the general nature of Baptism, the question that next awaits our consideration is, as to the subjects of Christian Baptism, or the parties to whom this ordinance ought to be administered. There are three opinions that may be maintained in regard to this matter. There is one party who affirm that Baptism ought to be administered to all, not infants, who are qualified to become members of the Christian Church in virtue of a credible profession of faith in Christ and a corresponding conduct. There is a second party who assert that Baptism rightfully belongs not only to such persons, but also, in virtue of a representative relation between parents and their offspring, to their children. And there is a third party who hold that Baptism ought to be administered without restriction to parents and children, without demanding, as a prerequisite from the applicant, any profession of faith or corresponding conduct. These three classes, holding principles markedly different from each other, probably exhaust the answers to the question: To whom is Baptism to be administered? The first, or the Antipædobaptists, administer the ordinance only to adults, who, by their faith and obedience, appear to be possessed personally of a title to be regarded as members of the Christian Church, and exclude infants, who cannot, by their own faith and profession, make good their claim to be regarded as proper subjects of the ordinance. The second, or the Pædobaptists, administer the ordinance not only to adults, who personally possess a right to be regarded as members of the Christian Church, but also to their infants, who can have no right except what they derive from their parents. And the third class, or the advocates of indiscriminate Baptism, administer the ordinance to all applicants without any restriction, and without demanding, in the case of adults, that they establish their claim to the ordinance by exhibiting a credible profession of faith in their own persons, or, in the case of infants, in the persons of their parents or guardians.

In proceeding to examine these different systems, it will not be necessary for me to discuss over again what occupied our attention at an early period of the course,---the question of what are the qualifications that give a person a title to be regarded as a member of the Christian Church,---or to enter into the controversy between Independents and Presbyterians as to the necessity in order to membership of a true and saving faith, or simply an outward profession and consistent practice. Without entering upon that subject a second time, the three systems of opinion as to the proper subjects of Baptism now mentioned may be conveniently discussed under the head of these two questions. First: Are we warranted by the Word of God to administer the ordinance of Baptism to all applicants for themselves or their children, without any restriction as to religious profession and character in the case of the applicant? And second: Are we warranted by the Word of God to administer the ordinance of Baptism to the children of a parent who would himself be a proper subject for Baptism, and is a member of the Church? The first question, or the point in debate between our Church and the advocates of indiscriminate Baptism, we shall now proceed to deal with, reserving the second, or the question of infant Baptism, for after consideration.

The doctrine of Baptism without restriction, and apart from the religious character and profession of the applicant, has assumed an aspect of more than ordinary importance recently, in consequence of the extent to which it has prevailed and the manner in which it has been advocated among Independents. Dr.~Wardlaw,---who was no friend of such a doctrine, but the reverse,---when speaking in reference to a former statement of opinion, to the effect that all parties were of one mind as to the necessity for a religious profession as a prerequisite to Baptism, says: ``Until of late, I had no idea of the degree or of the extent of this laxity, both as to the requisites in adults to their own baptism, and in parents, to the baptism of their children. It has been a cause of equal surprise and concern to me to find, from the publications of more than one of my brethren which have recently appeared, that in my first statement I have been so very wide of the truth. The lax views to which I now refer have been propounded and argued at length in the Congregational Lecture for 1844, by my esteemed friend, Dr.~Halley of Manchester.'' The surprise expressed by Dr.~Wardlaw at the acceptance which the doctrine of indiscriminate Baptism has received, and the prevalence which the practice has obtained among English Independents, is not without foundation. Dr.~Halley may, I believe, be fairly regarded as the representative of the views of Independents, at least in England, on the subject; and he is perhaps the ablest defender of the practice which prevails, very nearly universally, among them. The doctrine of the class to which he belongs, and whose views he advocates, is expressed by Halley as follows. After stating the principles held by other and opposite parties, he says: ``There are, lastly, those who baptize all applicants whatsoever, provided the application does not appear to be made scoffingly and profanely,---for that would be a manifest desecration of the service,---and all children offered by their parents, guardians, or others who may have the care of them.'' ``The third class maintain that, as no restriction is imposed upon baptism in the New Testament, none ought to be imposed by the ministers of the Gospel.'' ``These views,''---I quote again from Dr.~Wardlaw,---``these views, which he broaches and defends, are characterized by a latitudinarian laxity, which in my eyes is as mischievous as unscriptural,---the former, because the latter.''2 The question, then, of indiscriminate Baptism is one of very great interest and importance,---more especially in the present day,---and amply deserves discussion. In that discussion we must of course appeal for the only arguments which can decide the controversy to the Scriptures themselves. We learn from them that Baptism is a positive institution of Christ in the worship of the Christian Church; and from them also we must learn the terms on which the ordinance is to be dispensed, and the parties entitled to receive it. Is the ordinance, then, to be administered to all applicants indiscriminately without regard to religious profession or character,---to believers and unbelievers alike,---without any restriction, except, according to Dr.~Halley, that they do not apply for it ``scoffingly and profanely?'' Or, on the contrary, does a title to participation in the ordinance of Baptism imply, as a prerequisite, a religious profession and corresponding conduct on the part of the applicant?

Now, in examining into the doctrine and practice of Scripture bearing upon this question, it is important to understand distinctly at the outset the real point in debate. There are two preliminary remarks which may help to place it in its true light.

1st, The question in debate between the advocates and opponents of indiscriminate Baptism is not, as Dr.~Halley has stated it to be: ``Whether the Apostles and their assistants baptized indiscriminately all applicants, leaving their characters to be formed and tested by subsequent events.'' The question rather is: Whether, in such application made to the Apostles for Baptism, there was not included or implied a religious profession of faith in Christ, such as to warrant them to administer the ordinance because of the profession. It is manifest that, in apostolic times, when men were called upon in consequence of a Christian Baptism to forsake all that was dear to them on earth, and to incur the hazard of persecution and death, almost any such application necessarily involved or implied at least a credible profession of faith in Christ; inasmuch as hardly any conceivable motive except a belief in Christ would have induced any one to make the application, except, it may be, in rare and exceptional cases. Generally speaking, the fact of a man's applying for Baptism in apostolic times was itself the evidence of a credible profession, and enough to warrant the administration of the ordinance, not on the principle of baptizing all, believers and unbelievers alike, with a profession or without it; but rather on the principle that the applicant, by the very act of application, in the circumstances of the early Church, professed his faith in Christ. Upon this principle we can easily explain why, in the Scripture narrative of the practice of baptizing in the early Church, we find no example of the applicant being kept for a length of time in the position of candidate for Baptism, so as thereby to test his character and profession.

2d, The question in debate between the advocates and opponents of indiscriminate Baptism is not, whether the Apostles, in their administration of the ordinance, baptized, as Dr.~Halley asserts, ``bad men as well as good.'' That the Apostles did so in particular instances, the case of Simon Magus plainly attests. But that case no less plainly attests that the Baptism was administered, not in the absence of any religious profession, but in consequence of such a profession. Nothing can be more undeniable than that it was upon the ground of his professed belief in the Gospel preached by Philip that Simon Magus was baptized. ``Then Simon,'' says the inspired account of the transaction, ``then Simon himself believed also: and when he was baptized,'' etc. Like the other hearers who were baptized in consequence of their profession of faith in Philip's doctrine, Simon professed to believe, and, on the credit of that profession, was baptized as they were. But although among the number of those who received apostolic Baptism there were good men and bad men, as there must be among the members of the Church in all ages, this is not the real question at issue between the friends and opponents of indiscriminate Baptism. The real question in controversy between them is, whether Baptism was generally, or was ever, administered without a religious profession at all on the part of the applicant; or whether such a profession was invariably present as a prerequisite to Baptism. ``Baptism,'' says the Shorter Catechism, ``is not to be administered to any that are out of the visible Church, till they profess their faith in Christ and obedience to Him.''2

Bearing in mind these preliminary remarks, it is not difficult, I think, from an examination of Scripture doctrine and practice in regard to Baptism, to establish the conclusion, that it is a sacramental ordinance not to be administered indiscriminately and without restriction to all applying for it, but, on the contrary, limited to those maintaining an outward character and profession of Christianity.

\hypertarget{the-nature-and-import-of-the-ordinance-of-baptism-are-inconsistent-with-the-idea-of-an-indiscriminate-administration-of-it-to-all-without-respect-to-religious-character-and-profession.}{%
\subsection{The nature and import of the ordinance of Baptism are inconsistent with the idea of an indiscriminate administration of it to all, without respect to religious character and profession.}\label{the-nature-and-import-of-the-ordinance-of-baptism-are-inconsistent-with-the-idea-of-an-indiscriminate-administration-of-it-to-all-without-respect-to-religious-character-and-profession.}}

The doctrine and practice of the advocates of indiscriminate Baptism very naturally arise out of the system maintained by them as to the nature of the ordinance. With Dr.~Halley and the Independents, whom he represents, Baptism is not, in the proper and peculiar sense of the term, a Sacrament, but only a sign; and a sign, too, of a very restricted meaning indeed. It is a sign that the person holds certain Christian truths, or is willing to learn them; which truths may be held in the way of a mere intellectual apprehension, without the man who so holds them being a Christian, or even seriously professing to be one. Upon this theory,---that Baptism is no more than a sign, expressive of certain truths of Christianity,---it is quite possible to engraft the doctrine of an indiscriminate administration of the ordinance in every instance where those truths, as is usually the case in a Christian country, are not openly renounced or publicly denied. To affix the sign of allegiance to those truths in the case of every man who merely does not deny them, and must be held by the very act of applying for the sign, as at least in some tolerable degree acquainted with them, is consistent enough. To affix the sign to all infants proposed for Baptism, is also consistent; for they are capable of being instructed in the truths represented, and the act of their parents in bringing them to receive the ordinance may be regarded as an acknowledgment that they are willing that their children be so instructed. Restrict the import of Baptism to that of a mere sign of certain Gospel truths, and it is quite in accordance with the theory of indiscriminate administration. ``Practically,'' says Dr.~Halley, ``those who baptize indiscriminately all applicants and all children proposed for baptism, and those who reckon upon the prospect of teaching the baptized, will be found seldom at variance; for scarcely ever is any one proposed whose religious instruction might not be secured by proper care.'' As a sign expressive of acquaintance with certain Christian truths, or of a capacity and willingness to receive them, Baptism may consistently enough be administered without restriction to all applicants, whether adults or infants.

But the very opposite doctrine and practice must be maintained, on the supposition that the Sacrament of Baptism is not a sign merely, and that in a very restricted sense, of Christian truth, but a seal of a federal transaction between two parties in the ordinance, whereby the recipient gives himself in Baptism to Christ, and Christ in Baptism gives Himself and His grace to the recipient. A seal of a covenant which the party baptized does not even profess to make, and has avowedly no intention of entering into,---a voucher to a federal transaction, in which there is no person in the least professing to be a party,---an attestation to a mutual engagement never pretended by the individual who is supposed to give the attestation,---this is a contradiction and inconsistency not to be got over. There is a manifest incongruity in administering equally to those who avow that they are believers, and to unbelievers with no such avowal, the same Christian ordinance,---in dispensing a Gospel Sacrament indiscriminately to those who profess to have received the Gospel, and to those who do not,---in giving a religious privilege to those who make no religious profession, not less than to those who do. If Baptism be no more than a sign of certain religious truths known, or at least that may be learned, by the party baptized, then indeed there is no such incongruity between the nature of the rite and its unrestricted administration. But if Baptism be the outward seal of a federal engagement, distinctively marking the true Christian, then the very nature of the ordinance forbids it to be administered to men with no profession of Christianity. If it be the Sacrament of union to the Saviour and admission into the Christian Church, the ordinance itself points out the necessity of its restriction to those who ``name the name of Christ,'' and whose life and conduct are not outwardly inconsistent with their claim to be numbered among His people.

\hypertarget{the-administration-of-baptism-by-john-the-forerunner-of-our-lord-has-been-very-generally-appealed-to-in-favour-of-an-indiscriminate-dispensation-of-the-ordinance-but-in-point-of-fact-may-be-regarded-as-affording-evidence-of-a-contrary-practice.}{%
\subsection{The administration of Baptism by John, the forerunner of our Lord, has been very generally appealed to in favour of an indiscriminate dispensation of the ordinance, but in point of fact may be regarded as affording evidence of a contrary practice.}\label{the-administration-of-baptism-by-john-the-forerunner-of-our-lord-has-been-very-generally-appealed-to-in-favour-of-an-indiscriminate-dispensation-of-the-ordinance-but-in-point-of-fact-may-be-regarded-as-affording-evidence-of-a-contrary-practice.}}

The Baptism of John, when we are told that multitudes of the Jews flocked to him in the wilderness to be baptized, has been quoted in favour of the doctrine and practice of English Independents. There are two things which it is necessary to establish before any argument for indiscriminate Baptism in the Christian Church could be drawn from the preaching of John; and both these things, so far from being proved, may with good reason be denied. In the first place, it were necessary to prove that the Baptism of John was identical with Christian Baptism, before any countenance could be derived from his practice,---even if it were, as is alleged, that of indiscriminate Baptism,---in favour of the same custom in the Christian Church. And in the second place, it were necessary to establish the assumption that John really baptized all equally who came to him, without regard to their religious profession. I believe that neither the one nor the other of these positions can be established from Scripture, but the reverse.

With regard to the first position, there seems to be warrant from Scripture to say that John's Baptism was not identical with that of Christ. His doctrine and his office occupied an intermediate place between those of the Old Testament teachers and those of the Gospel Church; and his Baptism corresponded with his doctrine. He taught the doctrine of repentance and of preparation for Him that should come after him; he pointed to the future Saviour, rather than preached a present one; and his Baptism was the same in character. We have no reason to believe that he baptized in the name of Christ; and we have ground for asserting that the Baptism of John, in the case of those who received it, was afterwards replaced by Christian Baptism, when they were received into the Christian Church. That such was the case, the instance of the disciples at Ephesus proves; whom Paul rebaptized, as is recorded in the nineteenth chapter of the Acts of the Apostles: ``And he said unto them, Unto what then were ye baptized? And they said, Unto John's baptism. Then said Paul, John verily baptized with the baptism of repentance, saying unto the people, that they should believe on Him which should come after him, that is, on Christ Jesus. When they heard this, they were baptized in the name of the Lord Jesus.''

With regard to the second point, or the assumption that the Baptism of John was really given to all applicants indiscriminately, without respect to religious character, there seems to be no evidence for it in Scripture, but the reverse. We seem to have as good evidence, that John demanded a profession of a religious kind from those whom he baptized, as the character of the very brief and scanty narrative which has come down to us of the transaction would naturally lead us to expect. That vast multitudes of the Jews enrolled themselves by Baptism in the number of John's disciples, would appear to admit of no doubt; for we are expressly told that ``there went out unto him into the wilderness all the land of Judea, and they of Jerusalem, and were all baptized of him in the river Jordan.'' That of this great multitude all were truly brought to repentance, and turned from sin, and savingly taught to look forward to the Messiah who was to come, may, from many circumstances, appear improbable. But that they were all admitted to the ordinance of John's Baptism, without any regard to the religious profession that they made, is undeniably contradicted by the express language of the sacred historian; for it is added: ``They were all baptized of him in the river Jordan, confessing their sins.'' The Baptism and the confession of sins went together,---the one being the accompanying condition of the other. So far is it from being true that the practice of John gives countenance to the theory of indiscriminate Baptism, that the very opposite may be proved from the inspired narrative, brief though it be.

\hypertarget{the-terms-of-the-commission-given-by-our-lord-after-his-resurrection-to-his-apostles-in-regard-to-founding-and-establishing-the-christian-church-seem-very-clearly-to-forbid-the-practice-of-indiscriminate-baptism-and-to-require-a-profession-of-faith-in-christ-as-a-prerequisite-to-baptism-in-his-name.}{%
\subsection{The terms of the commission given by our Lord after His resurrection to His Apostles in regard to founding and establishing the Christian Church, seem very clearly to forbid the practice of indiscriminate Baptism, and to require a profession of faith in Christ as a prerequisite to Baptism in His name.}\label{the-terms-of-the-commission-given-by-our-lord-after-his-resurrection-to-his-apostles-in-regard-to-founding-and-establishing-the-christian-church-seem-very-clearly-to-forbid-the-practice-of-indiscriminate-baptism-and-to-require-a-profession-of-faith-in-christ-as-a-prerequisite-to-baptism-in-his-name.}}

The terms of the commission, as recorded in the Gospel by Matthew, are these: ``Go ye therefore, and disciple---μαθητευσατε---all nations, baptizing them into the name of the Father, and of the Son, and of the Holy Ghost; teaching them to observe all things whatsoever I have commanded you.'' Such is the language employed by our Lord in what must be regarded, I think, as the original institution of Christian Baptism. The commentary of Dr.~Halley on these words brings out his argument in favour of indiscriminate Baptism. ``The question,'' says he, ``respecting the subject of Baptism is here resolved into one of grammar and criticism. It is simply what is the antecedent to the word them, or for what noun is that pronoun substituted. `Going forth, disciple all the nations---παντα τα ἐθνη---baptizing them---αὐτους---all the nations, into the name of the Father, and the Son, and the Holy Ghost; teaching them---all the nations---to observe all things whatsoever I have commanded you.' So far as the grammatical construction is concerned, the meaning of the terms is precisely the same, as it would be if the words of the commission were, `baptize all the nations.' Adhering, therefore, to the grammar of the words, we say, the commission, which no man has a right to alter, is, `baptize all the nations.'\,'' Now, this somewhat summary and confident mode of reasoning may be satisfactorily set aside in two ways.

\begin{enumerate}
\def\labelenumi{\arabic{enumi}.}
\item
  There is some weight due to the order in which the terms of the commission run, as indicating the order in which the discipling, the baptizing, and the teaching of all the nations were to take place, and were to be accounted necessary parts of the Apostles' or the Church's obedience to the commission of Christ. There are three particulars embraced in the authoritative commission addressed to the Apostles, and, through them, binding upon the Church in every age. First, the command is to make disciples of all nations, turning them to the profession and belief of the faith of Christ. Second, there is the command to baptize all nations, granting them the formal and public rite by which their admission into the Church was to be attested and ratified. And third, there is the command to teach all nations to observe all things whatsoever Christ had appointed for His Church collectively, or His people individually. This is the order in which, according to the nature of the various particulars embraced in the commission, they were to be accomplished. That the order of procedure here indicated is in harmony with the nature of the work to be done by the Church in reference to the world, is abundantly plain from the scriptural account given of it in many other passages of the Bible. First of all is the preaching of the Gospel, as the grand instrument employed by the Church to gather in the disciples of Christ within its pale. Next there is the affixing to the disciples thus gathered the characteristic badge of discipleship, and granting them, by the initiatory rite of Baptism, formal admission into the Christian Church. And lastly, there is the instructing those thus admitted in the observance of all their appointed duties as disciples of Christ and members of His Church. This is plainly, I think, the order of procedure indicated in the apostolic commission; and it is an order which implies that a knowledge and profession of the faith as disciples preceded the administration of Baptism to them. The expression, ``all nations''---παντα τα ἐθνη---upon which Dr.~Halley builds his argument for universal and indiscriminate Baptism, is not to be regarded so much as declaring the duty of the Apostles to teach and baptize every individual of the world, or as denoting the absolute extent of the commission, as asserting that individuals of every nation were to be discipled and baptized, and marking out that no nation or class were excluded from the range of the commission. The terms, ``disciple,'' ``baptizing,'' must be taken together, and not separately; and in the order of the inspired declaration, and not in the reverse of that order.
\item
  The words of institution in the baptismal service seem to imply that a knowledge and profession of the faith of Christ are necessary as a prerequisite to Baptism. The recipients of the ordinance are to be baptized ``into the name, εἰς το ὀνομα, of the Father, the Son, and the Holy Ghost,''---language which obviously refers to the peculiar character the Three Persons of the Godhead sustain, and the offices they discharge in the work of man's redemption. Unless Christian Baptism, then, be a mere heathen mystery, to suffice as a sign or to work as a charm, it necessarily implies previous knowledge and instruction in the fundamental truths of the Gospel system; and this, again, implies that the Church, in administering the ordinance, has a right to require some evidence, such as an intelligent profession of the faith, that such knowledge has been obtained. All this points very distinctly to a profession of faith in Christ as a necessary prerequisite to the administration of the ordinance in the case of candidates for Baptism.
\end{enumerate}

\hypertarget{an-examination-in-detail-of-scripture-practice-as-bearing-upon-the-doctrine-of-indiscriminate-baptism-as-contradistinguished-from-baptism-restricted-to-professing-christians-will-sufficiently-bear-out-the-conclusion-to-be-drawn-from-the-previous-considerations-that-at-least-a-profession-of-faith-is-necessary-as-a-prerequisite-to-the-scriptural-administration-of-the-ordinance.}{%
\subsection{An examination in detail of Scripture practice, as bearing upon the doctrine of indiscriminate Baptism as contradistinguished from Baptism restricted to professing Christians, will sufficiently bear out the conclusion to be drawn from the previous considerations, that at least a profession of faith is necessary as a prerequisite to the scriptural administration of the ordinance.}\label{an-examination-in-detail-of-scripture-practice-as-bearing-upon-the-doctrine-of-indiscriminate-baptism-as-contradistinguished-from-baptism-restricted-to-professing-christians-will-sufficiently-bear-out-the-conclusion-to-be-drawn-from-the-previous-considerations-that-at-least-a-profession-of-faith-is-necessary-as-a-prerequisite-to-the-scriptural-administration-of-the-ordinance.}}

It is impossible, and indeed unnecessary, for us to enter at length into this field of argument. Nothing but the most violent injustice done to the language of Scripture by a bold and unscrupulous system of interpretation can suffice to get rid of the evidence which, in the case of the Baptism of converts mentioned in Scripture, connects the administration of the rite with a profession of faith in Christ on the part of the person who was the recipient of it. The association of the person's profession, faith, repentance, or believing, with Baptism, appears in a multitude of passages; while not one passage or example can be quoted in favour of the connection of Baptism with an absence of profession. ``He that believeth, and is baptized, shall be saved;'' ``repent every one of you, and be baptized;'' ``many having believed, and been baptized,''---these and many other passages of a like import connect together, as inseparable in the process by which under the eye of the Apostles many in their days were added to the Christian Church, the two facts of the religious profession of the candidate, and the administration of the religious ordinance by which formally he became a member of the Church of Christ. In the history, although brief and incomplete, of the Baptism of the early converts to the Christian faith, there is almost invariably some statement by which is attested the distinctive Christian profession that stands connected with the administration of the outward rite; while in no instances are there any statements from which it could be proved that Baptism ever stood connected with the absence of such a profession. Connected with the Baptism of the three thousand on the day of Pentecost, there stands the statement, ``Then they that gladly received the Word were baptized.'' Connected with the Baptism of the people of Samaria in consequence of the preaching of Philip, there stands the assertion, ``When they believed Philip preaching the things concerning the kingdom of God, and the name of Jesus Christ, they were baptized, both men and women.'' In regard to the Baptism of the Ethiopian treasurer, we are told that, after the Gospel was preached to him by the same evangelist, ``the eunuch said, See, here is water; what doth hinder me to be baptized? And Philip said, If thou believest with all thine heart, thou mayest. And he answered and said, I believe that Jesus Christ is the Son of God. And he commanded the chariot to stand still: and they went down both into the water, both Philip and the eunuch, and he baptized him.'' In connection with the Baptism of Lydia, and as preceding the administration of the rite, we have the statement: ``whose heart the Lord opened, that she attended unto the things that were spoken of Paul.'' Connected with the Baptism of the Philippian jailer, there stands the statement: ``And he rejoiced, believing in God, with all his house.'' In short, in almost every example of Baptism which the New Testament records, there is enough in the narrative, however scanty and compressed it be, to bring out the fact, that in close association with the administration of the rite appears the religious profession of the recipient. And, on the other hand, it may be safely asserted, that in no example of Baptism recorded in the New Testament can it be distinctly proved that no such profession was made.

What, then, is the answer given to this abundant and apparently satisfactory evidence for a Baptism restricted to and connected with a religious profession by the advocates of its indiscriminate administration? The answer given by them is twofold: first, that there are examples of bad men as well as good baptized by the Apostles; and second, that many or most of these Baptisms were administered so immediately in point of time after the profession made, that there was no opportunity to test by any satisfactory process the sincerity of it. Neither of these replies to the Scripture evidence is satisfactory. With regard to the first, or the fact that unbelievers and hypocrites were baptized, it is enough to say that we do not hold the Independent doctrine that a saving belief is necessary to entitle a man to Church membership; but, on the contrary, maintain that a profession of faith is enough, and that we have no security beyond the mere circumstance of an outwardly decent life against such profession being insincere. With regard to the second, or the fact that the profession on which the apostolic Baptisms in many instances proceeded could have been of no more than a few hours' standing, and therefore not proved by the lapse of time to be true, it is enough to say that there may be, and in apostolic times were, circumstances apart altogether from its duration sufficient to give credibility to the profession.2

\hypertarget{infant-baptism}{%
\section{Infant Baptism}\label{infant-baptism}}

We have now considered the question, To whom ought Baptism to be administered, in so far as it regards adults? The conclusion to which we were conducted was, that the ordinance ought to be dispensed to those alone who ``profess their faith in Christ, and their obedience to Him.'' The theory of indiscriminate Baptism we set aside as inconsistent with the nature and meaning of the Sacrament, as destitute of any countenance from the practice of John the Baptist, as contrary to the terms of the apostolical commission, and opposed to the practice of the apostles and the New Testament Church. There still remains for our consideration the question as to the connection of infants with Baptism, and as to the lawfulness or duty of administering the ordinance to them. The subject is a delicate and a difficult one, and demands a more than usually earnest investigation. The practice of baptizing infants may be regarded at first sight as running counter to all those views which we have already asserted in regard to the nature of Sacraments in general, and of Baptism in particular. Add to this, that it seems at first view directly to traverse the principles we have so lately laid down on the question of indiscriminate Baptism. The advocates of the doctrine of baptismal regeneration, who hold that Baptism is a charm with an inherent and independent power to confer grace in all circumstances and on all parties, can readily defend the practice of administering it to infants, as efficacious in the case of unconscious children, not less than in the case of intelligent adults. The advocates of the doctrine that Baptism is no more than a sign, have also an obvious ground on which they may defend the practice of infant Baptism,---the parents' professional badge being, not without reason or precedent in other matters, affixed to the child. And once more, the party who hold the doctrine of indiscriminate Baptism, and regard themselves as authorized to dispense the rite without regard to religious character or profession, can have no sufficient reason for excluding infants from this comprehensive commission. But if Baptism be the seal of a federal transaction between the party baptized and Christ; if this be the main and characteristic feature of the ordinance; and if a religious profession be a prerequisite to its reception; it would appear as if there were no small difficulty in the way of admitting to the participation of it those who, by reason of nonage, can be no parties to the engagement in virtue of their own act or will. The difficulty that stands in the way of infant Baptism lies on the very surface of the question; and Antipædobaptists have the advantage of an argument on their side which is both popular and plausible.

But in this case, as in all others connected with matters of positive institution in the Church of Christ, the primary and ruling consideration in the controversy must be the express Divine appointment on the subject. In those positive, and in a sense arbitrary, institutions, set up by God in the worship of His Church, mere inferential considerations drawn from reason must be of secondary authority and subordinate force to determine their nature and use, as compared with express intimations of the Divine will. Positive observances, from their very nature, must be regulated by positive institution; and it is only as secondary to such positive institution, that we can listen to arguments drawn from our views of the moral character or meaning of the ordinance. Our first appeal in the case of infant Baptism must, therefore, be to the express statements of the Word of God, and to the view of the ordinance as a positive institution which is there presented. We shall consider, then, in the first place, the scriptural principles which bear upon the question of the lawfulness or duty of infant Baptism. Thereafter we shall examine into the objections, from reason or Scripture, that have been brought against the practice; and also discuss the subject of the efficacy of the ordinance in the case of infants; and lastly, the scriptural mode of administering it.

What, then, is the bearing of Scripture doctrine and practice on the question of the lawfulness or unlawfulness of infant Baptism? The following five propositions I shall endeavour to establish in connection with this subject; and the discussion of these will very nearly exhaust the question. First, the covenant of grace, as revealed by God at different periods for the salvation of His people, has been essentially the same in former and in later times, and has always comprehended infants within it. Second, the Church of God, made up of His professing people, has been essentially the same in character in former and in later times, and has always included infants among its members. Third, the ordinance of outward admission into the Church has, in its essential character and meaning, been the same in former and in later times, and has always been administered to infants. Fourth, the principle on which the initiatory ordinance of admission into the Church has been administered has been the same in former and in later times, and has always applied to the case of infants. And fifth, the practice in regard to the administration of the initiatory rite has been the same in former and in later times, and has always included the case of infants. The illustration of these five propositions must, in consequence of the limits prescribed to us, be very brief, and more in the way of giving the heads of the argument than the argument itself. But taken under consideration even in the briefest way, they will embrace the prominent points of the controversy in regard to infant Baptism. One or more of them separately, if sufficiently established by an appeal to Scripture, would suffice to demonstrate that ``the infants of such as are members of the visible Church are to be baptized;'' while all taken together afford a very full and cumulative proof of the lawfulness of the practice.

\hypertarget{the-covenant-of-grace-as-revealed-by-god-at-different-periods-for-the-salvation-of-his-people-has-been-essentially-the-same-in-former-and-in-later-times-and-has-always-comprehended-infants-within-it.}{%
\subsection{The covenant of grace, as revealed by God at different periods for the salvation of His people, has been essentially the same in former and in later times, and has always comprehended infants within it.}\label{the-covenant-of-grace-as-revealed-by-god-at-different-periods-for-the-salvation-of-his-people-has-been-essentially-the-same-in-former-and-in-later-times-and-has-always-comprehended-infants-within-it.}}

This proposition is, properly speaking, made up of two: first, that the covenant was essentially the same in all ages; and second, that within the covenant, infants were always included. Neither of these two assertions ought to be very difficult of proof. In regard to the first, it is undeniable that God has had a people on the earth since the fall, chosen from the rest of mankind, who called upon His name, and were themselves called by it. The faith and hope of that chosen people, through every generation, have been sustained by a revelation of a Saviour, who either was to come or had come, expressed in promise and in type, in prediction and in symbol before His coming, and in plainer and ampler narrative of actual fact after His appearance. In whatever outward form it was revealed, this was God's covenant---His free promise of grace---His Gospel of glad tidings for the salvation of His people, identical in character and in substance, one in its announcements and its terms in every age from the first revelation in Paradise down to the last in Patmos. It was one and the same covenant of grace which was revealed to Adam in the first promise given to him, and the first ordinance of sacrifice appointed for him; revealed in other terms and form to Noah; repeated to Abraham in the word of promise and type; embodied in history, and prophecy, and symbolic institutions to the Church under the Mosaic economy; and fully brought to light under the Gospel dispensation. That the covenant of grace established under the Gospel was not then for the first time made known, but had been announced long before,---that although in the latter times it was more fully revealed, it had been revealed all along in substance, and proved to be the same at first as at the last,---the plain statements of Scripture very expressly affirm. The Apostle Paul tells us in the Epistle to the Galatians, that ``the Gospel was preached before unto Abraham.'' And in the same Epistle he tells us that ``the covenant confirmed of God in Christ was given to Abraham four hundred and thirty years before the giving of the law'' of Moses,---language fitted to mark both the identity of the covenant of Abraham with the Gospel covenant, and its independence of the Mosaic ceremonial institutions. If we turn to the book of Genesis, we shall find the account of the revelation of the covenant of grace given to Abraham, and referred to by Paul,---a covenant which, as then revealed, comprehended in it temporal blessings, such as the promise of Canaan to the patriarch and his seed, but was in itself independent of these; which preceded the law by more than four hundred years, and was not disannulled by the giving of the law; which was founded on the free grace and unchangeable promise of God, and thus was not bound up with any temporary institution; and which was the very Gospel afterwards ``confirmed in Christ.''2 So clear and abundant is the evidence for the first part of our proposition, that the covenant of grace, revealed under various forms in former and in latter times, was in substance one and the same.

The proof of the second part of our proposition is not less full and satisfactory, that this covenant has always comprehended infants within it. The infants of the parents with whom God's covenant was made, were not left outside that covenant. The promises of grace were not given to the parents, to the exclusion of the children. Infants were not left to their chance of uncovenanted mercies, while to adults the blessings were insured by covenant. On the contrary, that infants were comprehended within the covenant as well as their parents, is a fact that the plainest statements of Scripture demonstrate. In what sense or to what effect infants were so included, may come to be inquired into when we afterwards consider the efficacy of Baptism in their case, or the seal of the covenant as regards infants. But that the covenant made with the parents did not exclude but included their infant children also, the plain assertions of Scripture leave no room to doubt. In the inspired account of the various announcements made by God of His covenant from time to time, the terms of the announcement are almost invariably ``you and your seed.'' In the case of Abraham, as referred to by the Apostle Paul, this is very expressly stated: ``And God said unto Abraham, Thou shalt keep my covenant therefore, thou, and thy seed after thee in their generations. This is my covenant, which ye shall keep, between me and you, and thy seed after thee; Every man-child among you shall be circumcised.'' The covenant of grace, as then revealed to Abraham, included infant children of eight days old; and it has at all times been equally comprehensive and the same. The seal of the covenant, as affixed to the child when eight days old, was the standing evidence and memorial for two thousand years, that infants were included in God's federal promises."2

And in what manner is this argument from the example of Abraham, in favour of the fact that infants are comprehended within the covenant, met by the advocates of Antipædobaptist doctrines. The ordinary reply given by the opponents of infant Baptism is this: They affirm that there were two covenants, distinct and separate from each other, made by God with the patriarch at that time; the one a covenant of temporal, and the other of spiritual blessing. They assert that the ``seed'' mentioned in the history of the transaction, were the natural seed of Abraham, including adults and infants, in so far as regards the temporal covenant; and the spiritual seed of Abraham, or adult believers alone, in so far as regards the spiritual; and that the seal of circumcision administered to his children was the token of a temporal, and not a spiritual blessing. And lastly, they argue that under the Gospel the natural relationship of children to their parents, which under a former economy warranted their admission to the sign and seal of a temporal covenant, does not warrant their admission to the sign and seal of a spiritual one.

Now in regard to this attempted reply to the Scripture evidence for infants being included in the covenant of grace as revealed to Abraham, it is unnecessary to do more than make the following observations.

1st, Even although it were capable of being proved that there were two covenants made with Abraham, and not one simply,---a covenant of temporal blessing separated from the covenant of grace,---and that infants were included in the one but not in the other, this would not do away with the whole tenor of Scripture declaration in many other passages which evinces that the covenant of grace, under whatever shape and to whatever parties it was revealed, included not only the parties themselves, but also their infant offspring. The covenant of grace, as revealed to Abraham, and recorded in Genesis, has been very generally appealed to by the advocates of infant Baptism in demonstration of the interest infants had in it; and it has been so appealed to because it contains a more detailed and distinct evidence of the fact than most other passages of Scripture. But even were the record of the Abrahamic covenant expunged from the Bible, the interest of infants jointly with their parents in the covenant of grace could be satisfactorily established without it. The whole tenor of Scripture justifies us in saying, that it was a covenant which, at whatever time or in whatever form it was revealed to men, embraced both them and their infant seed.

2d, There is certainly no countenance in the narrative in Genesis given to the notion of two covenants, separate and distinct from each other; in the one of which the children of Abraham, being infants, were to have an interest, and in the other of which the descendants of Abraham, not being infants, but adult believers, were alone comprehended. The terms employed very expressly refer to one covenant, and not to two. ``Thou shalt keep my covenant. This is my covenant, which ye shall keep, betwixt me and you, and thy seed after thee.'' Such is the language emphatically reiterated in the original narrative of the transaction, marking a single covenant and not many. It is true, indeed, that there was a twofold blessing, the temporal and the spiritual,---the inheritance of Canaan, and the inheritance of the heavenly Canaan,---embodied in that one covenant. But these two orders of blessing were promised by the same covenant, and referred to the same end. There is no mention of one covenant intended for the natural posterity of the patriarch, and a second intended for his spiritual posterity. The temporal blessings might, indeed, be enjoyed by the descendants of Abraham after the flesh, while they had no interest in the spiritual; just in the same manner as a man under the Gospel may enjoy the outward privileges of a Church state without participation in the inward and saving blessings. But there is nothing whatever in the book of Genesis to warrant the distinction which the opponents of infant Baptism draw between a temporal covenant made with Abraham including infants, and a second and a spiritual one made at the same time and excluding them.

3d, The rite of circumcision, appointed for every man-child when eight days old, in the Abrahamic covenant as the token of it, excludes the theory of the Antipædobaptists, that the covenant in which infants were interested was a temporal covenant only. The fact that circumcision was ordained in connection with the covenant proves that it was not a mere temporal covenant, as Antipædobaptists allege, but a spiritual one,---the very covenant of grace which was the same through all times and dispensations of the Church. It does so in two ways. First, circumcision, as the token of the Abrahamic covenant, was a sign not of temporal, but of spiritual blessings. That this is the case is very expressly asserted by the Apostle Paul in his Epistle to the Romans. ``He is not a Jew,'' says Paul, ``which is one outwardly; neither is that circumcision which is outward in the flesh: but he is a Jew which is one inwardly: and circumcision is that of the heart, in the spirit, and not in the letter; whose praise is not of men, but of God.'' The ordinance of circumcision, then, had a spiritual import; it was expressive of Gospel blessings. And when it was appointed by God as the token of His covenant with Abraham, and administered in that capacity to children, it very plainly declared that the covenant, of which it was the token, and into which it introduced infants, was spiritual too. Circumcision, as the sign of the Gospel blessings, when it was appended to the covenant, demonstrated that the covenant itself was the covenant of grace. Second, circumcision is declared by the Apostle Paul to be more than a sign of grace; it is asserted to be a seal of grace. It is declared to be so, when he tells us, in reference to this very matter of the covenant established with Abraham, that ``he received the sign of circumcision, a seal of the righteousness of the faith which he had being yet uncircumcised.'' As the seal, then, of the covenant according to which Abraham was justified, the ordinance plainly testified that it was the covenant of grace; and, when administered to infants eight days old, it no less plainly indicated that they were interested in that covenant.3

The objections, then, brought by Antipædobaptists against the evidence from Scripture,---more especially derived from the covenant of grace as revealed to Abraham, but by no means confined to that source,---to the fact that infants are interested in that covenant, are of no great force. Our first position seems to be fairly established by Scripture evidence, namely, that the covenant of grace has been, under all the different forms in which from time to time it has been revealed, identical in substance and essentially unchanged; and that it has ever included infants within its provisions. The denial of infant Baptism cannot very well be maintained in the face of this proposition. If included in the provisions of the covenant of grace under the Gospel, infants must have a right to Baptism as one of them. They cannot be excluded from the initiatory ordinance which signifies and seals its blessings, unless the covenant of grace under the New Testament is different essentially both in its extent and in its terms from what it was before. The covenant of grace under former dispensations comprehended within its limits the infants of parties interested in it, as well as the parties themselves. This is undeniable. And the covenant must be altered essentially as to its extent,---it must be a different covenant as to the parties with whom it is made,---if so large a portion of the members included in it formerly, as infants were, should appear under the New Testament Church to be excluded. Further, it must be altered essentially as to the terms of it, and as to its free and gracious character,---it must be a different covenant as to the conditions of it,---if by these conditions one important class, made up of irresponsible parties such as infants, are now cast out when they were formerly comprehended. Unless the covenant of grace, in short, under the New Testament Church is another covenant from what it was under the Old Testament, infants must have a place in it now as much as then. But it is not so altered or restricted. Neither its extent nor its terms are altered. It is God's covenant of grace still; and as it was gracious enough and wide enough to comprehend within its limits infants under a former economy, it does so still.

There are manifold intimations in the New Testament that the covenant of grace is not less comprehensive in latter times than in former. At the first planting of the Christian Church the Apostle Peter assured the Jews that there was no change in this respect of the covenant under the Gospel economy as compared with its comprehensiveness under the Old Testament: ``For,'' said he, ``the promise is unto you, and to your children, and to all that are afar off, even to as many as the Lord our God shall call.'' To the Philippian jailer Paul declared in the very form of the Old Testament promises: ``Believe on the Lord Jesus Christ, and thou shalt be saved, and thy house.'' In these, and a multitude of other expressions of similar force and import, we recognise the great and important truth, that the covenant of grace was the same under the Gospel as under the law; that it was not limited or straitened in latter times in comparison with former; but that in its grace and comprehensiveness it embraces infants under the New Testament dispensation as well as under previous economies. We conclude, then, that the covenant of grace, revealed by God at different periods for the salvation of His people, has been essentially the same in former as in latter times; and has always comprehended infants within it.3

\hypertarget{my-next-proposition-is-that-the-church-of-god-made-up-of-his-professing-people-has-been-essentially-the-same-in-character-in-former-and-in-latter-times-and-has-always-included-infants-among-its-members.-this-second-proposition-like-the-first-consists-of-two-parts-each-of-which-admits-of-being-established-separately-the-first-part-of-the-statement-being-that-the-church-of-god-under-whatever-outward-form-it-has-appeared-has-been-identical-in-substance-throughout-every-dispensation-and-the-second-part-of-it-being-that-it-has-always-included-infants-among-its-members.}{%
\subsection{My next proposition is, that the Church of God, made up of His professing people, has been essentially the same in character in former and in latter times, and has always included infants among its members. This second proposition, like the first, consists of two parts, each of which admits of being established separately; the first part of the statement being, that the Church of God, under whatever outward form it has appeared, has been identical in substance throughout every dispensation; and the second part of it being, that it has always included infants among its members.}\label{my-next-proposition-is-that-the-church-of-god-made-up-of-his-professing-people-has-been-essentially-the-same-in-character-in-former-and-in-latter-times-and-has-always-included-infants-among-its-members.-this-second-proposition-like-the-first-consists-of-two-parts-each-of-which-admits-of-being-established-separately-the-first-part-of-the-statement-being-that-the-church-of-god-under-whatever-outward-form-it-has-appeared-has-been-identical-in-substance-throughout-every-dispensation-and-the-second-part-of-it-being-that-it-has-always-included-infants-among-its-members.}}

The first part of the proposition, which affirms the identity of the Church of God under all its outward forms, in Old Testament times and in New, may be readily demonstrated from two general considerations, independent of other arguments.

\begin{enumerate}
\def\labelenumi{\arabic{enumi}.}
\item
  The oneness of the covenant of grace in every age necessarily implies the oneness of the Church of God in every age. It was on the foundation of that covenant that the Church of God was built at first, and has ever since been maintained. It is that covenant that gives to its members every privilege which, as belonging to the Church of God, they possess; it defines the nature and limits the extent of their rights; it is the title by which they hold their standing and place as members of the Church; it constitutes the badge that distinguishes between a Church state and character, and the absence of them. The covenant is the charter of the Church of God in every age; and that charter remaining unchanged and identical from age to age, the Church that is built upon it must, in all its essential features, be one and the same also,---whatever may be the outward form it may bear, or the circumstantial and accidental changes that may be superinduced upon it. The Church of God in the days of Abraham,---the Church in the days of Moses,---the Church under the Gospel,---are in all vital respects the same; one Church, founded on the same covenant of grace, having the same essential character, and the same chartered rights, although different in outward things, according to the different stages and periods in the development of the Divine dispensations. The reason of this is obvious. The charter that constituted the society was the same in the earlier as in the later times. The covenant that called into existence and defined the character of the Church was essentially identical in the age of Abraham, and in the present age. We are not to confound with the unchanged and unchanging covenant of grace, on which the Church of God was and is built, the covenant made with Israel at Sinai, and destined to be a mere local and temporary ordinance. That subsequent and secondary covenant could neither disannul nor alter the former. It superinduced, indeed, upon the former certain local and temporary ordinances; but nowise enlarged, or contracted, or changed the original charter of the Church's existence and rights. The Church of Israel under the former economy, and the Church of Christ now under the Gospel, are constituted and defined as to their character, their extent, and their membership, by the same covenant of grace. They form the same society in their nature, their essential privileges, and their real members.
\item
  The identity of the Church of God in every age and under every dispensation, might be evinced by the relation which the Church ever bears to Christ as Mediator, and the relation which Christ as Mediator ever bears to the Church. Since the beginning He has been the Prophet, Priest, and King of the Church, immediately discharging all His offices as Mediator towards it, and sustaining it in existence by His continual presence in the midst of it. At different periods, indeed, He has been differently related to the Church, in so far as regards the extent of His manifestations of Himself, and the extent of His communications of spiritual gifts and blessings. But at no period has the Church existed, except through the same presence and power of Christ, as Mediator, that the Christian Church now enjoys,---the same in nature, although different in amount. The Church has ever been the Church of Christ; and this spiritual relationship, the same and unaltered from age to age, has caused the Church itself to be identical as a society throughout all times in its essential character, and privileges, and membership. Such considerations as these very clearly and abundantly attest the truth of the first part of our proposition, namely, that the Church of God, made up of His professing people, has been essentially the same in character in former and in later times.
\end{enumerate}

As regards the second part of the proposition, namely, that the Church has always included infants among its members, the proof, after what has already been said, need not demand a lengthened illustration. If the Church of God, made up of His professing people, be one and the same society at all times, and under all its different dispensations, then the proof that infants were members of it at one period must be a proof that they are competent to be members of it at any subsequent period; unless, indeed, some express and positive enactment can be produced, altering the charter of the society, and excluding, as incompetent to be admitted by the new and altered terms of the deed, those formerly comprehended within it. If no such proof of alteration in the charter or constitution of the society can be produced,---if the society itself remains the same in character and terms of admission as before,---then the proof that infants were once its members may suffice for proof that they are still competent to be so. We know that under the Abrahamic Church infants, as well as their parents, were admitted to the place of members. We have already proved that they were interested and comprehended in the covenant that constituted the Church in those days. The sign and seal of the covenant marked them out at eight days old, as embraced within it. The initiatory ordinance of the Church, which was the formal evidence of admission to its membership, was administered to the infants of such as were themselves members of the Church; and with that token in their flesh they grew up within the pale of the Church in Old Testament times. Circumcision was not part and parcel of the Sinaitic covenant, revealed afterwards through Moses. Our Lord Himself testifies that the ordinance was ``not of Moses, but of the fathers.'' It constituted the door of admission, not into the Sinaitic Church as distinct from the Abrahamic, but into that Church of which Abraham was a member, and of which all in every age are members who have like faith with Abraham. It constituted the door of admission, in the days of Abraham, into that very Church of which Christians are members now. And turning to Gospel times, we have a right to say that infants are competent to be members of the Christian Church now, unless it can be demonstrated that the Church of God is not the same now as in former times; that it is different in character and extent; and that those capable of admission before are, through an express alteration in the fundamental principles of the society, excluded now. Falling back upon our general proposition, already demonstrated, that the Church of God, as the society of His professing people, is one and the same in its essential nature in every age, we are entitled to affirm that infants once competent members of it are competent members of it still.

This proof is sufficient in the absence of any statute of limitation alleged to have been enacted in New Testament times, altering the character of the Church of God, and restricting it to the reception into its membership of adults, and adults alone. But there are very plain intimations in the New Testament, not only that no statute of limitation has been passed excluding infants, but that the privilege they once undeniably enjoyed under the Old Testament economy has been continued to them under the New. I do not dwell again upon the very express declaration of Peter to the Jews, when explaining to them the Gospel privilege: ``the promise is unto you and to your children,''---language which, in the case of a Jewish parent, could have only one meaning. I would refer to the language of our Lord Himself, when the Jewish parents brought their little ones to Christ, and He took them up in His arms and blessed them, accompanying the blessing with the words: ``Suffer little children to come unto me, and forbid them not; for of such is the kingdom of heaven.'' There can be no plausible interpretation of this passage given which proceeds upon the idea that those very infants blessed of Christ, and said by Him to belong to His kingdom, were actually excluded from it as its members. That they were not persons grown up, as one party of Antipædobaptists allege, but infants, who could by no act of their own profess their faith in Christ, is clear from the act of Christ taking them up in His arms when He blessed them. That the expression, ``of such is the kingdom of heaven,'' means no more than that persons of the like dispositions with children belonged to the kingdom of heaven, and that those very children were actually excluded from it, as another class of opponents of infant Baptism affirm, may be safely denied; inasmuch as the act of Christ in blessing them, in connection with the words He used, cannot be explained on the supposition that they were shut out beyond the pale of His covenant, and actually cut off from His Church. In short, the words of our Lord, taken in conjunction with His action, very distinctly demonstrate that the right of infants to be members of His Church, formerly recognised under the Old Testament, was not cancelled, but rather confirmed and continued under the New.2 We are entitled thus far to hold as proved our second grand proposition in all its parts, namely, that the Church of God, made up of all His professing people, has been essentially one in character in former and in latter times; and has always included infants among its members.

The two propositions, which we have already had under consideration, established as we believe them to be by Scripture evidence, go very far indeed, taken by themselves, to decide the question as to the lawfulness of infant Baptism. If infants as well as their parents have an interest in God's covenant,---if infants as well as their parents have a place in the Church as members,---it were difficult to affirm that they have no right to share in the privilege of Baptism, as the seal of the covenant, and the ordinance appointed for the formal admission into the Church of its members. An express prohibition forbidding the administration of the ordinance to them, or an incompatibility no less distinct between the nature of the Sacrament and their condition as infants, might, indeed, force upon us the conclusion that they are excepted. But in the absence of any such exception forced upon us by explicit prohibition or explicit incompatibility, we seem to be warranted in saying that the covenant state of infants and the Church state of infants, fairly demonstrated, unavoidably carry with them the inference that infants are entitled to the administration of Baptism as the seal of the one, and the door of formal admission into the other. The opponents of infant Baptism feel considerable difficulty in giving any explicit or consistent explanation of the relation sustained by infants either to the covenant or to the Church. Some of them deny absolutely that infants have any place either in the covenant or in the visible Church as members; while others of them hesitate about such a sweeping denial in the face of the strong Scripture evidence available to establish the fact, and rather consider infants as possessed of an inchoate and undeveloped right to be members, and as put under the care of the Church in order to be prepared for claiming and exercising the full right afterwards. But the covenant state and the Church state of infants, once fairly established, as they can readily be from Scripture, and the absence of any express bar interposed by Divine authority to the contrary, seem unquestionably to lead to a conclusion in favour of infant Baptism, even were there no further evidence that could be adduced in support of it. But there is much additional evidence at hand. The three propositions which still remain to be discussed and illustrated afford strong additional confirmation of the same conclusion; and, taken along with the positions already established, furnish a complete proof of the lawfulness and duty of baptizing infants.

\hypertarget{the-ordinance-of-outward-admission-into-the-church-has-in-its-essential-character-and-meaning-been-the-same-in-former-and-in-later-times-and-has-always-been-administered-to-infants.}{%
\subsection{The ordinance of outward admission into the Church has, in its essential character and meaning, been the same in former and in later times; and has always been administered to infants.}\label{the-ordinance-of-outward-admission-into-the-church-has-in-its-essential-character-and-meaning-been-the-same-in-former-and-in-later-times-and-has-always-been-administered-to-infants.}}

The main object of this third general proposition, as forming part of the argument for infant Baptism, is to identify, as essentially one and the same in their use and import and character, the Old Testament rite of circumcision with the New Testament rite of Baptism. If we can prove that they meant the same thing, and held the same place, and performed the same office in the Church of God in former and in later times, it were difficult to object to the conclusion that the one ought to be administered to the same infant members of the Church as was the other. To establish this general proposition we may make use of these three steps. First, circumcision and Baptism are both to be regarded as the appointed ordinance for the formal and public admission of its members into the Church. Second, both circumcision and Baptism have essentially the same meaning as the signs and seals of the same Divine truths and the same spiritual grace. Third, Baptism has been appointed to occupy the place and come in the room of circumcision, which has been done away.

In the first place, then, circumcision and Baptism are both to be regarded as the authorized ordinances for the formal admission of members into the Church.

That circumcision was the initiatory ordinance for the Old Testament Church, an appeal to the history of its institution and administration in ancient times will sufficiently evince. Without it no Israelite was accounted a member of the Old Testament Church; with it he could establish a right of membership, and a title to its ordinances. From the days of Abraham down to the date of the discontinuance of the ordinance in Gospel times, circumcision was the only thing that gave a right of admission to the privileges of the Old Testament Church; and apart from circumcision no one had a right to these. There was no access to the membership or ordinances of the ancient Church, except through the door of circumcision. That this was the case, is proved both from the case of infants and the case of adults. In the case of infants, the ordinance was universally administered; and in virtue of it alone, the circumcised infant, as it grew to manhood, was regarded as a member of the visible Church, and ceremonially qualified to receive its privileges without any other initiation or admission. In the case of adults, the administration of the rite to those who had not received it before,---as, for example, in the instance of Gentile proselytes,---was necessary as the door of admission into the fellowship of the Church. Without circumcision they were not admitted. By Divine appointment, circumcision bestowed on ``the stranger, who joined himself to the Lord,'' a right, the same as that of the Israelites themselves, to Church privileges and to partake of the passover. ``When a stranger shall sojourn with thee, and will keep the passover to the Lord,''---such were the terms of the enactment,---``let all his males be circumcised, and then let him come near and keep it; for no uncircumcised person shall eat thereof.'' Both in the case of infants, then, and of adults, circumcision constituted the initiatory ordinance of admission into the ancient Church from the days of Abraham downwards.

Against this fact, so very plainly attested in Scripture, it has been objected on the part of the opponents of infant Baptism, that it was not circumcision, but birth and natural descent, that gave admission into the ancient Church; and that every one born an Israelite became a member of the Israelitish Church. And in confirmation of this view, the fact of the circumcision of the descendants of Ishmael and Esau, without the observance giving them a title to admission to Church membership among the Israelites, is appealed to. The objection has not the least force in it. The tribes that sprang from Ishmael and Esau were divinely separated from the descendants of Abraham in the line of the covenant; and had not, like the other children of the patriarch, any interest in the federal promise. With these, therefore, circumcision could avail nothing to give them admission into the Church. Although practised by them, it was not with them a Church ordinance in connection with the covenant Church; and could not, therefore, admit them among its members. And on the other hand, mere birth did not give to the Israelite a right of admission into the Church, unless when connected with circumcision administered and submitted to. No Israelite was born a Church member. Unless, in addition to his birth as an Israelite, he was also circumcised, he had no right to the privileges of the ancient Church. So very far is it from being true, as some Antipædobaptists affirm, that his birth as an Israelite gave him a right to be considered a member of the Church, without circumcision, that it only placed him under the certainty of a heavy judicial sentence. To be born an Israelite, without circumcision being added to birth, only brought upon his head the sentence of God: ``He shall be cut off from his people.''3

There is quite as little foundation for another objection brought by other opponents of infant Baptism against our position, when they allege that circumcision was no more than a door of entrance to the Mosaic Church, and a token of admission to its outward and ceremonial privileges; and not the initiatory ordinance of the spiritual Church of God in Old Testament times. In answer to this objection, it is enough to say, that circumcision was instituted more than four hundred years before the legal economy was set up; and although it afterwards came to be associated with the law of Moses, yet it never lost its original meaning and use as the initiatory ordinance through which members entered into the Old Testament Church. It was in that character that we are to regard it when first instituted and administered in Abraham's family; and although four hundred years later there was superinduced upon the Church, to which circumcision was the door, a number of outward and ceremonial observances, yet it never ceased to be the initiatory rite of that Church of which Abraham was a member, and of which believers in every age, who have Abraham's faith, are members also. Under the Mosaic law, circumcision used and owned as an outward badge or privilege, admitted a man to an interest in an outward ceremonial institute; but not the less under the Mosaic law circumcision used and owned as a spiritual ordinance, and connected with the faith of the recipient, admitted also to an interest in that inner and spiritual Church, which was one and the same in the days of Abraham, in the time of Moses, and at the present time. Circumcision, although when associated with the Mosaic economy it was an outward badge of an outward Church, never ceased to be what it was at the first hour of its administration to Abraham himself, the ordinance of admission into the true Gospel Church.

The argument from Scripture, then, to prove that circumcision was the authorized ordinance for the admission of members into the Old Testament Church, is clear and satisfactory. It is hardly necessary to prove that Baptism is the authorized ordinance for the admission of members into the New Testament Church. That it is so, is admitted well nigh on all hands. The terms of the apostolic commission prove it to be so. The practice of Apostles and apostolic men in admitting converts to the Christian Church by Baptism, proves it to be so. The meaning of the ordinance as the Sacrament of union to Christ, proves it to be so. In this respect, the two ordinances occupy the same ground, and stand at the entrance of the Church publicly to mark and define its members; being the rites respectively belonging to the Old Testament Church and the New, for accomplishing the same object. To this extent, as the ordinance of admission into the Church of God, circumcision and Baptism are identical.

In the second place, circumcision and Baptism are expressive of the same spiritual truths, and are to be identified as signs and seals of the same covenant blessings.

With reference to circumcision, it is important to bear in mind that it was the sign and seal of a spiritual covenant, and not merely, as has been alleged, of the Sinai covenant, with its outward and ceremonial privileges. It was the covenant of grace as revealed to Abraham of which circumcision was primarily the token; and hence we have distinct evidence in Scripture that the spiritual blessings conveyed in that covenant to the believer were precisely the blessings which the ordinance of circumcision represents. The two cardinal blessings given by the covenant of grace are justification from guilt by faith in the righteousness of Christ, and sanctification from sin by the renewal of the heart through the work of the Holy Spirit; and these two blessings, we have express Scripture warrant to say, circumcision was intended to signify and seal. That circumcision was expressive of justification by faith in the righteousness of Christ, we are distinctly taught by the Apostle Paul to believe, in that passage of the Epistle to the Romans already more than once referred to: ``And Abraham,'' says the Apostle, ``received the sign of circumcision, a seal of the righteousness of the faith which he had, being yet uncircumcised.'' And again, that circumcision was a token of the sanctification of the heart and renewal from sin by the Spirit, is proved by several passages of Scripture which speak of ``the circumcision of the heart'' as the true meaning of the ordinance. ``He is not a Jew,'' says the same Apostle, ``which is one outwardly; neither is that circumcision which is outward in the flesh: but he is a Jew which is one inwardly; and circumcision is that of the heart, in the spirit, and not in the letter; whose praise is not of men, but of God.''2 These passages, and others which might easily be adduced, abundantly demonstrate that circumcision, as a sign and seal, represented and attested those two spiritual blessings of the covenant of grace, which are introductory to all the rest,---the blessings of justification and sanctification. And it is hardly necessary to add, that these are the two very blessings mainly and emphatically represented in the ordinance of Baptism under the New Testament Church. The very words of the Baptismal service tell us, that the member formally admitted into the Church is baptized ``into the name of the Father'' through means of justification by the Son, and sanctification through the Spirit. That is to say, the very same spiritual blessings represented and attested of old time by circumcision, are now represented and attested by Baptism. In this respect, as the signs and seals of the very same covenant blessings, circumcision and Baptism are one and the same.

In the third place, the oneness of circumcision and Baptism is yet further established by the fact that Baptism has come in the room of circumcision.

They are not only both initiatory ordinances for the admission of members into the Church, the one under the Old, and the other under the New Testament. They are not only appointed to be expressions of exactly the same spiritual truths, which stand permanently connected with the admission of a sinner into an interest in the covenant of grace. There is distinct enough evidence to show, that when circumcision was done away with at the establishment of the Gospel Church, Baptism was appointed to stand in its stead and fulfil its office. This appears, among other proofs, from the statement of the Apostle Paul in the Epistle to the Colossians. ``And ye are complete in Him,'' says the Apostle, referring to the unspeakable fulness of blessing laid up in Christ,---``and ye are complete in Him, who is the head of all principality and power; in whom also ye are circumcised with the circumcision made without hands, in putting off the body of the sins of the flesh by the circumcision of Christ: buried with Him in Baptism, wherein also ye are risen with Him.'' Such language seems plainly enough to imply that Baptism comes to Christians now in the room of circumcision to believers under the former dispensation; and that it is both fitted and intended to supply its place as a sign and seal of the blessings of the covenant. The reasoning of the Apostle appears very distinctly to intimate, that all which circumcision could do under the former dispensation, Baptism does now.

Upon these grounds, then, we are warranted to say that our third proposition is established,---namely, that the ordinance of admission into the Church has, in its essential character and meaning, been the same in former and in latter times, and has always been administered to infants.

\hypertarget{the-next-general-proposition-which-i-laid-down-at-the-outset-of-the-discussion-was-this-that-the-principle-on-which-the-initiatory-ordinance-of-admission-into-the-church-of-god-has-been-administered-has-been-the-same-in-former-and-in-latter-times-and-has-always-applied-to-the-case-of-infants.}{%
\subsection{The next general proposition which I laid down at the outset of the discussion was this, that the principle on which the initiatory ordinance of admission into the Church of God has been administered, has been the same in former and in latter times, and has always applied to the case of infants.}\label{the-next-general-proposition-which-i-laid-down-at-the-outset-of-the-discussion-was-this-that-the-principle-on-which-the-initiatory-ordinance-of-admission-into-the-church-of-god-has-been-administered-has-been-the-same-in-former-and-in-latter-times-and-has-always-applied-to-the-case-of-infants.}}

This is a proposition of much interest and importance as forming part of the argument for infant Baptism. What was the principle on which circumcision, recognising a title to membership in the Church under the Old Testament, was administered, and in accordance with which parties had a right to participate in the ordinance? This is the first question. What is the principle on which Baptism, recognising a title to membership in the Church under the Gospel, is administered, and in accordance with which parties have a right to participation in the ordinance? This is the second question. These questions in our present discussion must, of course, be restricted to the case of infants under both economies. The case of adults does not so directly concern our argument; and indeed in itself admits of little dispute. The personal act of the adult professing his religious faith is the ground on which, under the Old Testament in the case of proselytes, and under the Gospel in the case of converts, their right to be admitted as members of the Church, and to receive its initiatory ordinance, as the formal recognition of their admission, is obviously founded. But setting aside the case of adult proselytes or converts, upon what principle were infants entitled to circumcision in ancient times, and are infants entitled to Baptism in these latter days? Can it be established that the principle on which the ordinance is administered is one and the same in both cases?

1st, Upon what principle was the right of infants to circumcision founded under the Old Testament Church?

The analogy of the proceedings of God in providence and in grace not indistinctly points to the principle on which infants in the ancient Church were admitted to the same ordinance and to the membership of the same Church as their parents. By no personal act of theirs could infants become entitled, in the same manner as adults become entitled, to the privileges of the Church. But there is a familiar principle of representation, illustrated in the case of civil society, of providence, and of God's spiritual dispensations, in consequence of which infants, in certain cases and to certain effects, are held to be one with their parents, and through this relationship become entitled to the privileges of their parents. We see this representative principle in civil society, when, in consequence of no personal act of theirs, but simply in consequence of being accounted one with their father, infants become members of the civil society in which their father is a member, and their civil character and standing are the same as his. We see the representative principle, again, in the constitution of God's providence, when, in virtue of no deed of their own, but because of their relationship to their father, his place in society, his moral and intellectual character, his very bodily constitution for good or evil, to a certain extent become theirs. We see the representative principle, once more, in God's spiritual dispensation, where infants, in consequence of no personal act of theirs, but in accordance with that prevailing and universal constitution of things which is found in this world, become, in consequence of their filial relationship and the inheritance of the same flesh and blood as their father, concluded under his sin, and made one with him in original transgression and liability to punishment. In all these cases the representative principle is familiar to us, and infants are seen to partake for good or evil of the relations of their father. In most cases,---perhaps, if we were capable of understanding it, in all cases,---in which God deals with infants so as to show His method or law of dealing, He does so on the representative principle when He cannot deal with them on the principle of personal action and responsibility; and He acts with respect to them as if to a certain extent they were one with their parents.

That God may act towards infants in a way of sovereignty, without regard to their connection with their parents, may be true. But when He deals with them, and desires at the same time to manifest to us His rule or method of dealing, He does so on the principle of representation; a principle revealed to us both in His providential and spiritual economies. And such is unquestionably the principle according to which, in the constitution of the Old Testament Church, infants were dealt with. God made His covenant with infants as well as with adults; and the way in which He did so was never in connection with any personal act of theirs, which was impossible, but in connection with their filial relationship. God made His Church to include infants among its members as well as adult believers; and this too He did not in connection with their personal act, which was impossible, but in connection with the act of their parents. The membership of the father was counted to the infant; and the circumcision of the father gave a right to the infant to be circumcised also.

There are two views somewhat different from each other, that may be held on this point, which it is of considerable importance to discriminate between. The right of the child to circumcision and to the privileges of the Jewish Church, may be viewed as depending on his immediate father; or it may be regarded as depending on his remote progenitor, Abraham. In the one case, his title to be circumcised is counted good because of his relationship to his immediate parent, who was a member of the Jewish Church, and interested in the covenant. In the other case, his title to be circumcised is counted good because of his relationship to Abraham, his remote progenitor, with whom the covenant was made, and independently of his connection with his immediate parent, and without regard to the circumstance of his parent being or not being a member of the Jewish Church. The evidence of Scripture seems not indistinctly to point to the first view as the correct one, or to the view that connects the right of the infant directly with his immediate father's interest in the Church and covenant, and not the view that connects it indirectly with Abraham's. Dr.~Halley advocates the view that connects the infant's right not with the parent's, but with Abraham's interest in the covenant, making that right independent of the parent's connection or non-connection with the Church; and he does so apparently with the view of founding upon it the doctrine of indiscriminate Baptism to all infants alike, whatever be the father's Church state, and whether he be a member of the Church or not. The two following considerations, however, seem very decisively to prove that the right of the infant to circumcision in the Jewish Church was valid in consequence of the Church membership of the father, and not in consequence of his remoter connection with Abraham. First, mere connection with Abraham did not in all cases give a right to the privileges of the Jewish Church, as we see exemplified in the instance of the descendants of Abraham in the lines of Ishmael and Esau. They were directly connected with Abraham as their ancestor, and yet were separated from the communion of the Jewish Church. Second, the case of the infants of Gentile proselytes demonstrates that not remote connection with Abraham, but immediate connection with the parent, is the ground of the infant's right to circumcision. The infants of such Gentile proselytes as were circumcised and members of the Jewish Church, had no connection with Abraham through ordinary descent; and yet in virtue of their father's circumcision they had a right to be circumcised also. These two considerations seem sufficient to prove that the right of the infant to circumcision was not derived remotely from Abraham, passing over his immediate parent, but came directly from the parent. In other words, the case of circumcision under the Old Testament presents to us a complete and perfect illustration of the representative principle, and of the privileges of the child being held to be the same as those of the parent. By no personal act of their own did children become entitled to circumcision; but they were so entitled, in consequence of the right of their father to the ordinance.

2d, Now, what is the principle on which infants under the New Testament Church become entitled to Baptism? Are we warranted by Scripture in identifying the principle on which Baptism is administered now with the principle on which circumcision was administered before? I think that we are. The identity in meaning, and character, and use, already proved between circumcision and Baptism, would afford a strong presumption in favour of the conclusion, even had we no further evidence for it. The strong and close analogy between the two cases would go very far of itself to establish it. But there is one passage of Scripture more especially, which seems of itself explicitly to announce that the very principle of representation found under the Old Testament in the case of parent and child, is not cancelled, but continued under the New, and must be held as a permanent principle in the dealings of God with infants. The passage to which I refer is in the First Epistle to the Corinthians, and is to the following effect. Speaking of the case of husband and wife, when one of the parties is not a Christian but an unbeliever, the Apostle says: ``For the unbelieving husband is sanctified (ἡγιασται) by the wife, and the unbelieving wife is sanctified by the husband: else were your children unclean; but now are they holy (ἁγια).'' The principle of representation found under the Old Testament is the very principle introduced by the Apostle to explain the position and character of children in the case where no more than one parent is a believer and member of the Church.2 That the contrasted terms, ``unclean'' and ``holy,'' are to be understood in the Old Testament sense of not set apart and set apart to the service or fellowship of God, seems to be undoubted. And the assertion of the Apostle is, that one of the parents being a believer, although the other is not, avails, so that the infants are to be accounted clean, or fit for the service of God and the fellowship of His Church. The holiness of the one parent that is a member of the Christian Church, communicates a relative holiness to the infant, so that the child also is fitted to be a member of the Church, and to be baptized. The forced and unnatural interpretation put upon this passage by Antipædobaptists cannot stand a moment's investigation. They interpret the ``cleanness'' of the infant as the legitimacy of the infant,---a construction plainly forbidden by the consideration that marriages are lawful, and the children legitimate, whether the parents be believers or unbelievers. In this passage, then, we have a very express avowal of the principle of representation, proved to obtain in the case of circumcision under the Old Testament. The child is accounted clean because the parent is clean; or, to translate the phrase into ecclesiastical language, the child is entitled to Church membership because the parent is a Church member.2 We recognise at once the identity of the principle under the former economy and the present; and we are entitled to hold as proved the fourth of our general propositions, namely, that the principle on which the initiatory ordinance of admission into the Church of God has been administered, has been the same in former and in latter times, and has always applied to the case of infants.

\hypertarget{the-practice-in-regard-to-the-administration-of-the-initiatory-ordinance-has-been-the-same-in-former-and-in-latter-times-and-has-always-included-the-case-of-infants.}{%
\subsection{The practice in regard to the administration of the initiatory ordinance has been the same in former and in latter times, and has always included the case of infants.}\label{the-practice-in-regard-to-the-administration-of-the-initiatory-ordinance-has-been-the-same-in-former-and-in-latter-times-and-has-always-included-the-case-of-infants.}}

This is the fifth and last of the general propositions which I laid down at the outset; and after what has already been established, it requires no more than the briefest notice. Of course in regard to the practice of the Old Testament Church the proposition may be regarded as proved; the circumcision of the infant eight days old being the standing proof of the practice of the Church in former times. With regard to the practice of the Church under the Gospel, there are two preliminary remarks which it is important to carry along with us. First, the uniform practice of the ancient Church down to the epoch of the Gospel, taken in connection with the total silence of Scripture as to any change of practice when the Jewish passed into the Christian Church, is itself very nearly conclusive as to the practice of the early Christians in regard to infant Baptism. Second, there is not a single instance among all the Baptisms recorded in Scripture in which we find a person, who had grown up a Christian and without Baptism, receiving the ordinance when he became an adult. We have many examples of adult Baptism in Scripture, but none of adults who for years had been Christians before they received the ordinance.

Carrying these remarks along with us, nothing more is necessary, in regard to the practice of the Primitive Church in the matter of infant Baptism, than to refer to the frequent and almost constant mention of the Baptism of ``households'' and ``families,'' in which it is morally certain that there must have been infant members. ``I baptized the household of Stephanas.'' ``He was baptized, and all his, straightway.'' ``She was baptized, and her household,'' etc. Such expressions as these, interpreted in the light of the previous undoubted practice of the Jewish Church, can admit of only one meaning. Infants are not mentioned specifically as baptized along with the parents, because it is taken for granted that everybody understood that they were. Had they been pointedly and separately mentioned in such cases, it would very fairly and reasonably have given rise to the suspicion or inference that infant Baptism was in principle an entire novelty, that it was a new thing for the Church to have infant members. The notices of household and family Baptisms, that occur in the New Testament so repeatedly, cannot be explained on the theory of the Antipædobaptists, that the family or household were adults. In the case of Lydia, for example, it is said: ``She was baptized, and her household.'' If, according to the theory of the opponents of infant Baptism, the household of Lydia consisted of adults, who separately and personally were converted like herself, and on a personal profession of faith like hers were separately baptized, it is very difficult to understand why their conversion and Baptism were not, like hers, separately mentioned, or on what principle they are all merged under her single name. Upon the theory of infant Baptism, on the contrary, it is easy to understand how infants, with no personal profession of faith, and no conversion like her own, were merged under her name as ``her household.'' Under the circumstances of the Apostolic Church, the repeated mention of household or family Baptism is of itself decisive evidence of the practice by which infants were baptized. We are justified in saying that our fifth and last proposition, like the former, is sufficiently established, namely, that the practice in regard to the administration of the initiatory ordinance of the Church has been the same in former and in latter times, and has always included the case of infants.2

\hypertarget{objections-to-infant-baptism}{%
\section{Objections To Infant Baptism}\label{objections-to-infant-baptism}}

We have been occupied of late with the consideration of the general principles laid down in Scripture, upon which the lawfulness and duty of the Baptism of infants may be argued. I have endeavoured to establish and explain five general propositions, from any of which singly, but more especially from all taken together, may be drawn a proof in favour of infant Baptism. In doing this I adopted, as upon the whole the best, the plan of following the natural order of the argument, without caring to turn aside at every step to answer the objections which Antipædobaptists have urged against it, except when these lay directly in the line of my own illustration of it. In the right understanding of the argument itself, there is contained an answer to these objections, so that they may be considered as in a good degree met by anticipation. But still, as the subject is an important one, and as it may better help to develop the principles of the argument, I shall now proceed to consider some of the most common and plausible of the objections brought by Antipædobaptists against the relevancy or conclusiveness of our reasonings.

That in the case of infants baptized, there are difficulties connected with their condition as infants, which it may be hard to solve, it would be useless to deny. But that those difficulties, in one form or other, are peculiar to infant Baptism, and nowhere else to be met with, may reasonably be questioned. Above all, that those difficulties should be permitted to overbear the very strong and cumulative evidence from Scripture in favour of the doctrine and practice, it is not the part of truth or wisdom to assert. And yet I believe that it is mainly those difficulties which have led many to scruple to accept as valid or conclusive the Scripture evidence for infant Baptism. In what sense, or to what effect, infants are interested in the ordinance of Baptism, or benefited by it; what explanation is to be given of the use and efficacy of the Sacrament in their case; in what manner we are to reconcile infant participation in the sign and seal of the covenant of grace with the absence of intelligence and responsibility in infants: these are difficulties which have had more to do in bringing about that state of mind which has led many to declare infant Baptism to be unscriptural, than the force of Scripture argument against it. I believe that these difficulties which have influenced so many against the practice of infant Baptism, and which at first sight appear to be peculiar to it, are not really peculiar to it. In one shape or other, and to a greater or less extent, these difficulties are to be encountered in the case of adult Baptism as much as in the case of infant; and, indeed, are common to the supernatural grace or virtue connected with all Divine ordinances. Such difficulties may appear more palpably and prominently in their association with infant Baptism, and by many have been regarded as connected with it alone; but in reality they will be found in greater or less measure present, wherever we admit that the work of the Spirit of God in His own ordinances is present, making them the means or instruments of supernatural grace.

This matter will come on for consideration at a subsequent stage, when I proceed to deal with the question of the efficacy of Baptism in the case of infants. I advert to it at present for the purpose of indicating my conviction that the source of not a few of the objections to infant Baptism is to be found, not in the Scripture evidence against it, but rather in those difficulties which are thought to embarrass the theory or explanation of its efficacy. It is plain that, in the first instance, our duty is to examine and weigh the Scripture evidence on the subject, and to be guided in our belief and practice by its force and conclusiveness. It is only in the second instance that it is lawful for us to inquire as to what explanation is to be given of the difficulties which stand connected with the Scripture ordinance. Objections drawn from the mere difficulty of framing a theological theory of the Sacrament, in its application to infants, are not for one instant to be allowed to contradict Scripture evidence, where it is clear and conclusive on the subject. That such evidence we have in support of infant Baptism, the heads of argument already given may be enough to evince. Postponing, then, for after consideration, the question of the efficacy of the ordinance in the case of infants, and the difficulties alleged to be connected with that point, because that question ought not to be allowed to interfere with the Scripture evidence to be weighed and examined in the first place, I now go on to consider some of the common and most plausible objections to that evidence as it has been already laid down.

The objections generally urged against the Scripture argument for infant Baptism, may be ranged under two heads: those which deny the relevancy of a large portion of our reasoning; and those which controvert the conclusiveness of it. There are two general objections which I shall examine, as commonly urged against the relevancy of the argument; and there are two objections also which I shall notice, directed against the conclusiveness of our reasoning. Under these heads we shall probably be able to discuss all that is of much weight or plausibility in the objections of Antipædobaptists.

\hypertarget{under-the-head-of-objections-to-the-relevancy-of-our-reasoning-in-favour-of-infant-baptism-i-remark-in-the-first-place-that-not-a-few-object-to-our-argument-as-one-based-upon-as-they-allege-an-outward-and-ceremonial-dispensation-that-was-to-be-done-away-and-which-has-no-place-under-the-gospel.-they-regard-our-reasoning-from-the-abrahamic-covenant-as-irrelevant-to-our-duty-or-practice-under-the-gospel-economy-and-hold-that-in-transplanting-the-custom-of-affixing-to-infants-the-outward-seal-of-the-covenant-from-the-ancient-to-the-present-dispensation-we-are-borrowing-the-carnal-ordinances-of-a-bygone-time-and-giving-them-without-warrant-and-unlawfully-a-place-in-the-spiritual-church-of-christ.}{%
\subsection{Under the head of objections to the relevancy of our reasoning in favour of infant Baptism, I remark in the first place, that not a few object to our argument as one based upon, as they allege, an outward and ceremonial dispensation that was to be done away, and which has no place under the Gospel. They regard our reasoning from the Abrahamic covenant as irrelevant to our duty or practice under the Gospel economy; and hold that, in transplanting the custom of affixing to infants the outward seal of the covenant from the ancient to the present dispensation, we are borrowing the carnal ordinances of a bygone time, and giving them, without warrant and unlawfully, a place in the spiritual Church of Christ.}\label{under-the-head-of-objections-to-the-relevancy-of-our-reasoning-in-favour-of-infant-baptism-i-remark-in-the-first-place-that-not-a-few-object-to-our-argument-as-one-based-upon-as-they-allege-an-outward-and-ceremonial-dispensation-that-was-to-be-done-away-and-which-has-no-place-under-the-gospel.-they-regard-our-reasoning-from-the-abrahamic-covenant-as-irrelevant-to-our-duty-or-practice-under-the-gospel-economy-and-hold-that-in-transplanting-the-custom-of-affixing-to-infants-the-outward-seal-of-the-covenant-from-the-ancient-to-the-present-dispensation-we-are-borrowing-the-carnal-ordinances-of-a-bygone-time-and-giving-them-without-warrant-and-unlawfully-a-place-in-the-spiritual-church-of-christ.}}

Now in reference to this objection, it is at once admitted, that the argument for infant Baptism rests partly, although not by any means exclusively, upon a consideration of the Abrahamic covenant and Church. But it rests upon nothing peculiar to that Church, or that has been done away with. It is not unfrequently demanded of the advocates of infant Baptism, why they so often begin their argument in favour of a New Testament ordinance, such as Baptism, from the days of Abraham and from the nature of the covenant made with him. The answer to such a question is very plain. We not unfrequently begin with the Abrahamic covenant in the argument for infant Baptism, because with Abraham the Gospel Church was first formally established, and endowed with that ordinance which we believe to be in its character and use identical with Baptism. No doubt the Church of God had existed from the days of the first promise made to Adam of a Saviour, and of the first believer in that promise; and downward to the present time, under all its different forms, a Church has existed in this world. But with Abraham, and not before, began that outward provision in the Church for the admission of infants by means of an initiatory rite which was to signify and seal their interest in the covenant of grace; and therefore, in seeking to ascertain the meaning and nature and use of that initiatory rite, whether you view it under the form of circumcision in other days, or of Baptism now, it is both natural and lawful to go back to its origin and first institution the better to understand it. Circumcision was, in short, the Baptism of the Church of God in former days; and in arguing in respect to its use and administration, it is both justifiable and reasonable to inquire into its origin, and into the terms on which it was originally enforced. Nor is there the slightest ground for alleging that in doing this we are guilty of transplanting an Old Testament, carnal, and temporary practice into the New Testament and spiritual Church without warrant, and against the meaning and nature of Gospel ordinances. It is granted, that there is a vast and unspeakable difference between the spirituality of the Gospel dispensation and the outward and ceremonial nature of the Jewish economy. But it is carefully to be remarked,---and if marked, would prevent much confusion in the argument,---that although in popular and common language we are wont to speak of the Jewish and Christian Churches as if they were two separate and contrasted Churches, and not one Church under two dispensations, yet strictly speaking the expression is not correct, and has led to much confusion both of thought and argument on this question as well as on others. There were two dispensations, the Jewish and the Christian; a carnal and outward dispensation, and a spiritual and more inward one. But it was the same Church of God under both, identical in character and essence, and all that is fundamental to a Church; although in the one case, under the Mosaic dispensation, it was the Church encircled by and subsisting in a carnal and outward economy, and in the other case, under the Gospel dispensation, it was the same Church encircled by and subsisting in a less outward and more spiritual economy. What belonged to the mere dispensation within which the Church of God was at any time encircled might be done away; what belonged to the Church itself was not to be done away.

There are two brief considerations that will be sufficient to remove the objection to the relevancy of our argument for infant Baptism, from the alleged fact that it is built upon the practice of a former and temporary dispensation.

\begin{enumerate}
\def\labelenumi{\arabic{enumi}.}
\item
  As already indicated, the objection is founded on the fallacy that the Old Testament Church and the New Testament Church were not one but different Churches; the one being carnal and the other spiritual,---the one being outward and ceremonial, as contrasted with the other, which is not so. It is hardly necessary to repeat what has already been largely established, that the Church of God has been one and the same in all ages, whether it is made up of ``the household of Abraham'' whom the patriarch circumcised, or ``the household of Stephanas'' whom Paul baptized; whether it numbers as its members Jews as in the days of Moses, or Gentiles as in our own. The outward dispensation superinduced upon the Church was changed from time to time; but the Church itself remained the same. Circumcision did not belong to the dispensation; it belonged to the Church. The initiatory ordinance by which infants were admitted as its members, was appointed more than four hundred years before the Jewish dispensation, and was administered before as well as during the period of the ceremonial economy. That economy, with its legal observances and symbolic ritual, might have been removed, as indeed it was removed, at the introduction of the Gospel dispensation; and yet, had God not intended to introduce Baptism in the place of circumcision in these latter times, circumcision might have still remained in force as the initiatory rite of His Church, in virtue of the place which it had in the Abrahamic covenant. Circumcision was independent either of the introduction or abolition of the law of Moses; and would have continued the standing ordinance for admission into the Church of God, as the seal of the covenant of grace, had not Baptism been expressly appointed as a substitute for it.
\item
  The objection to our reasoning, that it is founded on the practice of a bygone and temporary dispensation, arises partly out of a misapprehension in regard to the typical nature of the ordinance. Under the general and comprehensive formula that all types are now merged in their antitypes, and that all that was symbolic in other days is abolished in the New Testament Church, Antipædobaptists have argued in support of the conclusion that circumcision belonged to a temporary economy, which can be no precedent under the Gospel. Now circumcision may, it is frankly admitted, have served the purpose of a type of Christian sanctification under the ancient economy; and as a type, it had place no longer than until the antitype was realized. But it cannot be denied that it served another purpose also. It cannot be denied that it was instituted and used as a sacramental ordinance in the Church of God, altogether apart from its typical character as expressive of Christian regeneration; that it was, in short, a sign and seal of the covenant of grace. And in this character, which it unquestionably sustained, over and above its typical one, we cannot regard it as part and parcel of the Mosaic institute; nor is there any ground for alleging that, in appealing to the authority of circumcision in favour of infant Baptism, we are appealing to a carnal dispensation as a precedent for the practice of the Gospel Church.
\end{enumerate}

\hypertarget{but-under-the-head-of-objections-to-the-relevancy-of-our-reasoning-for-infant-baptism-i-remark-in-the-second-place-that-not-a-few-object-to-our-argument-because-as-they-allege-it-is-applicable-to-an-outward-but-not-applicable-to-a-spiritual-church.-this-second-objection-is-no-more-than-a-modification-of-the-preceding-one.-it-is-allied-to-the-fallacy-that-circumcision-was-the-badge-of-a-temporary-and-typical-dispensation-opposed-to-the-spirit-of-the-gospel-and-not-to-be-represented-under-the-gospel-by-any-parallel-or-identical-ordinance-equally-binding-and-equally-administered-to-infants.}{%
\subsection{But under the head of objections to the relevancy of our reasoning for infant Baptism, I remark, in the second place, that not a few object to our argument, because, as they allege, it is applicable to an outward, but not applicable to a spiritual, Church. This second objection is no more than a modification of the preceding one. It is allied to the fallacy that circumcision was the badge of a temporary and typical dispensation, opposed to the spirit of the Gospel, and not to be represented under the Gospel by any parallel or identical ordinance, equally binding, and equally administered to infants.}\label{but-under-the-head-of-objections-to-the-relevancy-of-our-reasoning-for-infant-baptism-i-remark-in-the-second-place-that-not-a-few-object-to-our-argument-because-as-they-allege-it-is-applicable-to-an-outward-but-not-applicable-to-a-spiritual-church.-this-second-objection-is-no-more-than-a-modification-of-the-preceding-one.-it-is-allied-to-the-fallacy-that-circumcision-was-the-badge-of-a-temporary-and-typical-dispensation-opposed-to-the-spirit-of-the-gospel-and-not-to-be-represented-under-the-gospel-by-any-parallel-or-identical-ordinance-equally-binding-and-equally-administered-to-infants.}}

In many cases, the source of the feeling which regards infant Baptism as akin to an outward but unsuited to the character of a spiritual Church, is to be found in the denial of the Scripture distinction, so important to be kept in mind, between the visible and invisible Church. When the character of the Church as a visible corporate society is ignored or denied,---when the Church on earth is identified with the invisible Church made up of true believers alone,---when the title to membership in the Church here below is restricted to a saving faith in Christ and regeneration by His Spirit, and none but those possessed of saving faith are considered to have a right to entrance,---when such views as to the nature of the Church and its membership are held, it is not unnatural, but the reverse, that infants should be regarded as not members of the Church, and that infant Baptism should be accounted a misapplication of the ordinance. And hence, historically, it is a fact of great significance and interest, that among Independents, who deny the distinction between the visible and invisible Church, mainly, if not entirely, have been found also that religious party who deny infant Baptism; while among Presbyterians, whose principles lead them to mark distinctly and maintain strongly the difference between the visible and invisible Church, few or no deniers of the lawfulness of infant Baptism have been found. I feel myself exempted from the necessity of falling back upon the question of the grounds on which the important distinction between the visible and invisible Church of Christ rests, inasmuch as these have been fully argued at a previous stage in our discussions. It is enough for me to remind you that the Church of Christ, as exhibited in this world, has, as we have already established, a visible and corporate character, and is possessed of certain outward privileges and certain outward ordinances, by which it is known in the eyes of men, as well as an inward and spiritual character, by which it is known in the eyes of God; that the tares grow side by side with the wheat in the enclosure of the Christian Church; and that even the external provision of ordinances and Sacraments, administered, although they may be, in numberless instances, to merely nominal Christians, is not to be undervalued or set aside, but rather esteemed a gift of God to His Church exceedingly great and precious. The ordinance of Baptism, administered to infants as well as to adults, forms part of the outward provision of ordinance which God has made for the visible Church. And it is an unscriptural theory, which, by denying the existence of such a Church, and assuming one purely and exclusively spiritual, would bear with an unfriendly influence on the doctrine and practice of infant Baptism.

But passing from the objections to the relevancy of our argument in favour of infant Baptism, I go on to consider some of the more common objections to the conclusiveness of our reasonings.

1st, Under the head of the objections to the conclusiveness of the reasoning in favour of infant Baptism, I remark, in the first place, that it has been objected against infant Baptism that there is no express or explicit command in the New Testament to administer the ordinance to infants.

It is readily admitted that Baptism is a positive institution; and that in regard to the nature and use of positive institutions in the Church of Christ we must be guided solely by the communications of the Word of God in regard to them. But that the objection to infant Baptism from the absence of a positive and articulate formula, enjoining the administration of the Sacrament to infants, is of no real force, can be readily evinced.

First, the absence in Scripture of an express formula enjoining any duty, is no proof that the duty is not required; and the absence of any express formula imposing the duty of infant Baptism in particular, is no argument against the practice, but the reverse. Looking at the proposition as a general one applicable to all cases, it is evidently both unwarrantable and perilous to lay down as a canon of Scripture interpretation, that whenever there is no express and explicit injunction, in so many words, requiring a duty to be performed, there the deed is unlawful, or at least not commanded. It is unwarrantable; because we have no right to limit God as to the form in which He may be pleased to make known to us His will, if, in one form or other, it is made known. It is perilous as regards ourselves; because there can be no more dangerous position than to assume the attitude of refusing to regard the will of God intimated to us, because it is not intimated in the manner which we may consider the plainest and the best. Whatever is laid upon us in Scripture, whether it be in the way of direct and explicit commandment, or in the way of indirect but necessary inference from what is commanded, is equally binding and of Divine obligation.

But the absence of any express formula enforcing the Baptism of infants in Scripture is more especially and emphatically to be regarded as no argument against the practice, but rather an argument on its side. A positive formula for infant Baptism, parallel to that which was given to the Apostles, to preach the Gospel, and to baptize all nations, would have looked very much as if infant Baptism was a novelty in the Church, unknown in principle and substance before. To preach the Gospel to the Gentiles, to baptize the Gentiles, were duties unknown to the exclusiveness of the Jewish Church; and hence a new and express formula enjoining them was necessary at the outset of the new economy. Had the admission of infants as members been equally unknown to the Church, there would have been a no less urgent necessity for an express and explicit command in regard to it. But infants had been accounted and treated as members of the Church of God for well nigh four thousand years; and at the era of the Gospel dispensation there was no need for the proclamation of any new law in regard to their admission. Any such new law formally enjoining it might well have given rise to the idea that the practice had never been heard of before; that it was as much a new thing in the Church as seeking to proselytize and baptize the Gentile nations was. All that was necessary was a positive intimation that the outward manner of admitting infants into the Church was to be different under the Gospel from what it was before,---that the ordinance of Baptism was to be used instead of circumcision; and such an intimation is very expressly given both in the way of precept and example in the New Testament. Anything beyond this in the shape of an express formula to admit infants into the Church would reasonably have led to the belief that they had been excluded before.

Second, in reply to the objection to infant Baptism taken from the absence of any explicit injunction of the practice, it may be remarked that exactly the same objection may be brought against other Christian duties, which notwithstanding are generally or universally acknowledged to be duties, because, in the absence of an express command, the authority of Scripture imposing them can be certainly learned by ``good and necessary inference.'' For example, the duty of females to commemorate the Lord's death at His table, and the duty of keeping the Sabbath under the Gospel, are not, it has often been remarked, expressly enjoined by any separate formula in the New Testament Scriptures. The duty of females to join in the Lord's Supper is only to be gathered inferentially by a process of reasoning not more direct than that which establishes the lawfulness and duty of infant Baptism. In like manner, the duty of keeping the first day of the week holy unto the Lord can claim no express or separate injunction in the New Testament any more than the practice of infant Baptism can.

There is a marked resemblance, indeed, between the sanctification of the first day of the week and the practice of baptizing infants, in regard both to what is enjoined and what is left to be inferred in respect of each, in the New Testament. The sanctification of one day in seven was not a new appointment in the Christian Church, but rested on the practice and authority of the more ancient dispensation of God; and hence there is no re-enactment in the New Testament of the general Sabbath law. But the change in the circumstance of the time when the Sabbath was to be kept, was a new appointment under the Gospel; and hence, by explicit examples of an authoritative kind, the change of the day is intimated and fixed in the New Testament. Exactly parallel to this, the admission of infants as members of the Church was no new appointment in the Church of God at the introduction of the Gospel dispensation; and hence it was left very much to rest for its authority on the previous law and practice of the Church, without any re-enactment of what was binding before. But the change in the form of admitting infants into the Church,---the change from circumcision to Baptism,---was a new appointment; and hence, by explicit command and example in the New Testament, we have authority for the change.

Third, in reply to the objection against infant Baptism, drawn from the absence of any separate authority for the practice, it might be enough to challenge the Antipædobaptist upon his own principles to prove his own practice to be scriptural; and show an explicit precept or explicit precedent for baptizing the child of a Church member not along with the parent in his infancy, but afterwards when the child has grown to manhood. The inspired history of the Christian Church contained in the Acts of the Apostles embraces a period of more than twice the number of years required to allow the infants of a baptized convert themselves to grow up to the years of discretion, when they might have been accounted able to make a personal profession of their faith, as their parents had done before; and yet there is neither precept nor example in Scripture giving express authority for baptizing the children of Christian parents, after they had grown up to years of maturity, apart from the case of adult converts, which forms common ground to both parties in this controversy. Tried by their own principles, the practice of Antipædobaptists would be found wanting in Scripture authority.

2d, Under the head of objections to the conclusiveness of our reasoning for infant Baptism, I remark further, that it is commonly or universally objected by Antipædobaptists against the practice of infant Baptism, that faith, or at least a profession of faith, in Christ, is positively demanded as a prerequisite to Baptism in all cases; and that as infants cannot have such faith, or make such a profession, they cannot be admitted to the ordinance. Of the fact asserted in this objection, namely, that a profession of faith is required, both by the scriptural commission given to the Apostles to baptize, and by the apostolic examples in this matter, on the part of the person to be baptized in all ordinary cases, there is no room for doubt. We have already had occasion to illustrate and assert the fact against the doctrine and practice of indiscriminate Baptism. But the fact there asserted is too narrow a foundation to build an objection on against infant Baptism.

In the first place, the demand of Scripture for faith or a profession of faith, as a prerequisite for Baptism, is a demand that has respect to adults, and is not addressed to infants; and not being addressed to infants, it cannot be regarded as laying down the conditions or terms on which infants are to be made partakers of the ordinance. It is quite plain that those passages of Scripture in which a profession of faith is connected with Baptism, like the Scriptures at large, are intended for adults and not for infants,---for the common and general case of men in the full possession of their intellectual and moral powers, and not for the exceptional case of infants not in full possession of those powers. That this is the case, the single consideration that the Bible is God's message to men and not to infants, is enough to prove; unless it could be shown, which it cannot, that in those passages, not men but infants are specifically referred to. The passages usually quoted by Antipædobaptists in support of their objection, are the commission to the Apostles, as recorded in Mark, and the saying of Philip to the Ethiopian eunuch, recorded in the Acts of the Apostles. The apostolic commission in Mark is to this effect: ``Go ye into all the world, and preach the Gospel to every creature. He that believeth, and is baptized, shall be saved; he that believeth not shall be condemned.'' It is abundantly obvious that this language applies primarily to the ordinary case of adults, and not to the exceptional case of infants; and while the order---first belief, and then Baptism---refers to adults, it cannot apply to infants, to whom the Gospel cannot be preached, and who cannot be expected to believe it. Are infants, then, in virtue of this passage, to be excluded from Baptism, because in consequence of their infancy they are excluded from believing? Certainly not; for by the very same argument they would be excluded also from salvation. The order of the passage is, first, belief; second, Baptism; third, salvation. And if, on the strength of this passage, infants, as Antipædobaptists assert, are to be excluded from Baptism because they are excluded from believing, they must, in like manner, be excluded from salvation too.

The saying of Philip addressed to the Ethiopian eunuch, is quite as little available for the Antipædobaptist objection. ``If,'' said Philip, addressing the man upon whose understanding and heart there had dawned, through the evangelist's preaching, a saving knowledge of Christ,---``if thou believest with all thine heart, thou mayest be baptized.'' The language was addressed to an adult in the full possession of all his powers of mind, and laid down for him the order of faith as preceding Baptism. But Philip never applied the same language, nor laid down the same order, in the extraordinary case of infants, whose salvation must be according to a different order and a different method. The announcements of Scripture which imply the necessity of faith or a profession of faith in order to Baptism, are framed upon the principle of adult Baptism, not upon the exceptional case of infant Baptism.

In the second place, the objection of Antipædobaptists, grounded on the impossibility of infants complying with the conditions on which Baptism ought to be administered, may be proved to be fallacious by a consideration of the case of circumcised infants. That infants were circumcised, and had a title to be so, will not by any party be denied. And yet circumcision involved in it the very same profession of faith, in all its essential respects, that Baptism now does. Substantially, it is the same ordinance as Baptism. It expressed the same truths. It implied on the part of the worthy recipient essentially the same spiritual qualifications. That this was the case is very expressly asserted by the Apostle Paul in his Epistle to the Galatians. ``Every man,'' says he, ``that is circumcised is a debtor to do the whole law.'' In other words, circumcision in the case of the person circumcised involved a profession of his obligation to keep God's law, very much in the same manner as Baptism involves such a profession now. And yet infants, incapable of making such a profession, were circumcised. And exactly on the same principle, infants incapable now of making such a profession are to be baptized.

In the third place, the objection of Antipædobaptists may be proved to be groundless by a consideration of the case of infants saved. The very same difficulty, if difficulty it can be called, alleged to stand in the way of the doctrine of infant Baptism, applies with undiminished force to the case of infant salvation. ``He that believeth shall be saved; he that believeth not shall be condemned.'' Such is the simple and unchangeable formula that declares in Scripture the order and connection of faith and salvation. It is a formula adopted and intended to apply to the case of adults, responsible for their belief; and it makes the salvation of their souls to be suspended on the existence of their faith. Interpreted in the same manner, and applied in the same unlimited extent to infants, it would close against them the door of the kingdom of heaven, and exclude the possibility of their salvation; for they are incapable, by reason of their infancy, of that faith which stands connected with the justification of the sinner before God. Shall we, in virtue of the Antipædobaptist canon of criticism, proceed to reverse the Saviour's words, and turn His blessing into a curse, and say in regard to infants, that of such is not the kingdom of heaven? Or shall we not, on the contrary, reject a canon of interpretation that would lead to such results, and rather say that infants are subjects both of Baptism and salvation?

\hypertarget{the-efficacy-of-infant-baptism}{%
\section{The Efficacy Of Infant Baptism}\label{the-efficacy-of-infant-baptism}}

The efficacy of Baptism in the case of adults may be understood' from what has been already said of the nature of the Sacraments in general. Baptism, like the Lord's Supper, is a sign and seal of a federal engagement between the receiver and Christ. It presupposes the existence of justifying and saving grace in the person baptized; and it seals or attests that grace to the soul, in this manner becoming the means of further grace.

There is a meaning in the fact that the person receiving the Sacrament has a part to perform in the ordinance,---that in the Lord's Supper he personally takes and partakes of the elements of bread and wine, and that in Baptism he personally submits himself to and receives the sprinkling of water. In both Sacraments there is a personal act on the part of the participator, which has its spiritual meaning, which cannot and ought not to be overlooked in the transaction. That act forms the link that connects the receiver of the ordinance with the ordinance itself; and the spiritual faith embodied in the act forms the link which connects his soul with the covenant blessings which the ordinance represents. The Sacrament is a seal, then, of more than the covenant generally; it is a seal of the covenant in its appropriation by the believer to himself personally in the ordinance.

There are some theologians indeed who in their explanation of the Sacraments make them seals of the covenant in general, and not seals of the believer's own personal interest in the covenant. They make the Sacraments attestations vouching for God's promises of grace at large, but not vouching for those promises as appropriated by the believer and realized in the experience of the worthy receiver of the Sacrament. This explanation of the Sacraments, however, is, I think, much too narrow and limited. It overlooks the personal act of the receiver in the Sacrament, and the spiritual meaning of that act. It disowns or neglects as not essential to the ordinance, the part which the participator has to perform, when in the case of the Lord's Supper he personally takes of the bread and wine, or when in the case of Baptism he personally presents himself to be sprinkled with water in the name of the Trinity. There is a spiritual meaning in these personal acts not to be overlooked in our explanation of the Sacraments, and essential to a right understanding of them. These personal acts constitute the part performed by the believer in the covenant transaction between him and Christ in the ordinance, and are necessary to make up the covenant. And the Sacrament, as a seal, is applicable to that part of the covenant transaction by which the believer appropriated the blessing to himself, not less than to that other part of the covenant transaction by which Christ exhibits or makes offer of the promise of grace to the believer. In other words, the Sacrament is not merely a seal of the covenant offered, or exhibited, or declared in general, but a seal of the covenant appropriated by the believer in particular, and, through means of his own spiritual act in the ordinance as well as Christ's, received in his personal experience.

In the case of Baptism administered to a believing adult, his own personal part in the ordinance, when he presents himself to the sprinkling of water, is the sign of that spiritual act of his through which the blessings of justification and regeneration, represented in the Sacrament, have previously become his; and Baptism is to him a seal not merely of these blessings as exhibited and promised in the covenant generally, but of these blessings realized and enjoyed by himself. Through the channel of his faith, and by means of the Spirit in the ordinance, Baptism becomes a seal in his justification and regeneration, and so a means of grace and spiritual blessing to his soul.

Such is the efficacy of Baptism administered to an adult believer. What is the virtue or efficacy of the ordinance when administered to infants incapable of faith, although not incapable of being made partakers in the grace which the Spirit confers? In entering on the consideration of this delicate and difficult subject, it is necessary, in order to clear our way to it, to lay down one or two preliminary propositions of much importance in the discussion.

First, The proper and true type of Baptism, as a Sacrament in the Church of Christ, is the Baptism of adults, and not the Baptism of infants. In consequence of the altered circumstances of the Christian Church at present, as compared with the era when Baptism was first appointed, we are apt to overlook this truth. The growth and prevalence of the visible Church, and the comparative fewness of the instances of adult conversion to an outward profession of Christianity amongst us, have led to the Baptism of infants being almost the only Baptism with which we are familiar. The very opposite of this was witnessed in the Church of Christ at first. And the true type of Baptism, from examining which we are to gather our notions of its nature and efficacy, is to be found in the adult Baptisms of the early days of Christianity, and not in the only Baptism commonly practised now in the professing Church, the Baptism of infants. It is of very great importance, in dealing with the question of the nature and efficacy of Baptism, to remember this. Both among the enemies and the friends of infant Baptism the neglect of this distinction has been the occasion of numberless errors in regard to the import and effects of the Sacrament. Men have judged of the nature and efficacy of Baptism from the type of the ordinance, as exhibited in the case of baptized adults. They have reversed the legitimate order of the argument, and argued from the case of infants to that of adults, and not from the case of adults to that of infants. It is abundantly obvious that adult Baptism is the rule, and infant Baptism the exceptional case; and we must take our idea of the ordinance in its nature and effects not from the exception, but from the rule. The ordinance of Baptism is no more to be judged of from its ministration to children, than is the ordinance of preaching to be judged of from its ministration to children. The Sacrament in its complete features and perfect character is to be witnessed in the case of those subjects of it whose moral and intellectual nature has been fully developed and is entire, and not in the case of those subjects of it whose moral and intellectual being is no more than rudimental and in embryo. Infants are subjects of Baptism in so far as, and no farther than their spiritual and intellectual nature permits of it. And it is an error, abundant illustration of which could be given from the writings both of the advocates and opponents of infant Baptism, to make Baptism applicable in the same sense and to the same extent to infants and to adults, and to form our notions and frame our theory of the Sacrament from its character as exhibited in the case of infants. It is very plain, and very important to remember, that the only true and complete type of Baptism is found in the instance of those subjects of it who are capable both of faith and repentance, not in the instance of those subjects of it who are not capable of either. The Bible model of Baptism is adult Baptism, and not infant.

Second, The virtue of infant Baptism, whatever that may be, is not more mysterious than the virtue ascribed to adult Baptism, although it may have the appearance of being so. It is a very common idea, that the difficulty in framing an explanation of the efficacy of Baptism in the case of infants, is peculiar to the ordinance in its administration to them, and does not attach to it in its administration to adults. I believe that this is not the case. There may be greater difficulty in gathering from the statements of Scripture what the virtue of Baptism really is in its application to infants, than in ascertaining what it is in its application to adults. But to explain the supernatural virtue itself is just as difficult in the one case as in the other, and simply from this reason, that it is supernatural. Up to a certain point it is easy enough to explain the efficacy of adult Baptism, but beyond that fixed point it is impossible to explain it. That point is where the natural efficacy of the ordinance passes into the supernatural efficacy. There is a certain natural influence which Baptism, as expressive of certain spiritual truths, and through means of these truths, is fitted to exert upon the adult, because he is a moral and intelligent being, with his faculties developed and complete. And this natural influence of Baptism, through means of the truths expressed by it, cannot be exerted upon the infant, because, although he is a moral and intelligent being, his faculties are not developed or complete. As a sign of spiritual truths understood by the adult, and not understood by the infant, Baptism has a certain natural effect on the one and not on the other, which it is not difficult to explain. But this effect is moral or natural, and not, properly speaking, the sacramental efficacy that is peculiar to the ordinance. The sacramental efficacy peculiar to the ordinance is not natural, but supernatural,---an efficacy not belonging to it from its moral character, but belonging to it in consequence of the presence and power of the Spirit of God in the ordinance. This distinctive efficacy of Baptism as a Sacrament, we cannot understand or explain, either in the case of adults or the case of infants. It is a supernatural effect of a gracious kind, wrought by the Spirit of God in connection with the ordinance; and because it is supernatural, it is not more and not less a mystery in the case of infants than in the case of adults.

The supernatural efficacy connected with Baptism, and owing to the presence of the Spirit of God with the ordinance, is an efficacy competent to infants as much as to adults. Even upon their unconscious natures the Spirit is free to work His work of grace, not less than upon the natures of adults whose understandings and hearts are consciously consenting to the work. The work of regeneration by the Holy Ghost is a work which it is as easy for Him to accomplish upon the infant of days as upon the man of mature age,---upon the child who enjoys but the rudiments of his moral and intellectual life, as upon the adult whose moral and intellectual powers are co-operating in and consenting to the gracious change. But broadly marked although the regeneration of the infant and the regeneration of the adult be, by the absence in the one instance, and the presence in the other, of a capacity moral and intellectual for faith and repentance, yet it is never to be lost sight of or forgotten that the work is the work of the Spirit of God, and not to be explained on any natural principle either in the former case or in the latter. The presence of his complete and perfect intellectual and moral powers in the case of the baptized adult, and the exercise of those powers in connection with the truths represented and signified in the Sacrament, afford no adequate explanation of the sacramental grace or efficacy connected with the ordinance in consequence of the power of the Spirit in it. At this point we have got beyond the limits of the natural, and into the region of the supernatural; and it is not more and not less supernatural in the case of infants than in the case of adults. Sacramental grace, properly so called, is a mystery of which there is no explanation, except that it is the grace of the Spirit of God. Admit that this grace is conveyed in any given case through the channel of Baptism to the believing adult, and you admit a mystery, which the presence and active exercise of his moral and intellectual powers do not in the least explain. Admit that this grace is conveyed in any given case through the channel of Baptism to the infant incapable of believing, and you admit a mystery too, but one not more mysterious than the former, and not more difficult to explain, from the absence or incapacity of his moral and intellectual faculties. In one word, the efficacy of infant Baptism, whatever that may be shown from Scripture to be, is not more mysterious than the sacramental virtue ascribed to adult Baptism.

Bearing in mind these preliminary remarks, what, I ask, are the effects of Baptism in so far as regards infants baptized? I do not pause at present in order to examine into the nature and benefit of the ordinance in so far as regards parents, who, in the exercise of a parent's right to represent their unconscious children, claim the administration of the ordinance for their offspring. In acting as the substitute for the infant, who cannot act for itself, in the solemn federal transaction between it and Christ,---in becoming a party in its name to the covenant made between the baptized infant and its Saviour through the ordinance,---the parent comes under a very great and solemn obligation on behalf of the child, thus pledged and given to the Redeemer through the parent's deed and not its own. But passing by this, let us confine our attention to the case of the infant, and proceed to inquire what are the benefits and efficacy of Baptism to the infant participators in the ordinance? In the case of adults, we know that Baptism is fitted and designed not to confer faith, but rather to confirm it,---not to originate grace, but to increase it,---not to effect that inward change of regeneration by which we are numbered with the children of God, or that outward change of justification by which we are accepted of Him, but to seal these blessings before bestowed. With adults, Baptism is not regeneration or justification, but the seal of both to the regenerated and justified man. And in the case of infants, the Sacrament cannot be regarded as accomplishing without their faith, what in the case of adults with their faith, it fails to accomplish. In other words, infant Baptism is not infant regeneration or justification, any more than in the instance of adults. The Baptism with water to a child is not the same thing as the birth by the Spirit. It is not a supernatural charm. It is not a magic spell to confer the washing of regeneration and the renewal of the Holy Ghost. Sacraments in the case of infants, as in the case of adults, have no mysterious and supernatural power of their own to impart, by the bare administration of them, spiritual life. Let us endeavour to understand what are the effects of Baptism in the case of infants.

\hypertarget{baptism-in-the-case-of-all-infants-baptized-gives-to-them-an-interest-in-the-church-of-christ-as-its-members.}{%
\subsection{Baptism, in the case of all infants baptized, gives to them an interest in the Church of Christ, as its members.}\label{baptism-in-the-case-of-all-infants-baptized-gives-to-them-an-interest-in-the-church-of-christ-as-its-members.}}

Circumcision gave to infants in other days a place in the ancient Church as its members; and they grew up within its pale entitled to all its outward privileges and rights, needing no other admission in after life. And what circumcision did during the time when it was in force, that Baptism does now in regard to infants baptized. It constitutes the door of admission into that visible Church of God on earth of which the parent himself is a member; and the baptized one grows up within the pale of its distinctive communion, needing no other admission, marked off at least outwardly from a world that has no interest in God, and having a right to the enjoyment of privileges which, as an outward provision for His own in this earth, God has given to them and not to the world. And this of itself is no small privilege, outward and temporal though it be, and not inward and spiritual. That outward provision of the means of grace, which has been given to the visible Church in this world for its establishment and benefit, is always represented in Scripture as a gift of Christ to His people, not to be undervalued or despised because it comes short, in those who enjoy it, of a saving blessing, but rather to be accounted exceeding great and precious. It is a gift of Christ to His Church which is of such worth and moment that the giving of it is spoken of in the Word of God as one of the great purposes for which the Saviour ascended up on high. ``When He ascended up on high,'' says the Apostle Paul in his Epistle to the Ephesians,---``when He ascended up on high, He led captivity captive, and gave gifts unto men. And He gave some, apostles; and some, prophets; and some, evangelists; and some, pastors and teachers; for the perfecting of the saints, for the work of the ministry, for the edifying of the body of Christ.'' That outward provision of ordinances and means of grace for the visible Church, the bestowment of which is thus represented as one of the grand objects for which Christ left this world and ascended to the Father, must be to that Church of no ordinary importance and value. It is a right to this provision of outward ordinances and means of grace which the baptized infant receives, when by his Baptism he becomes formally a member of the visible Church; and growing up in the use and enjoyment of them, the benefit to him, although short of a saving benefit, is beyond all price. Baptism as the sign of membership and the passport to the infant into the sanctuary of the visible Church, does not bestow the saving blessing, but brings him in after life into contact with the blessing; it does not constitute him a member of the kingdom of heaven, but it brings him to the very door, and bids him there knock and it shall be opened unto him.

\hypertarget{baptism-in-the-case-of-all-infants-baptized-gives-them-a-right-of-property-in-the-covenant-of-grace-which-may-in-after-life-by-means-of-their-personal-faith-be-supplemented-by-a-right-of-possession.}{%
\subsection{Baptism, in the case of all infants baptized, gives them a right of property in the covenant of grace; which may in after life, by means of their personal faith, be supplemented by a right of possession.}\label{baptism-in-the-case-of-all-infants-baptized-gives-them-a-right-of-property-in-the-covenant-of-grace-which-may-in-after-life-by-means-of-their-personal-faith-be-supplemented-by-a-right-of-possession.}}

In regard to this matter, I would have recourse again to a distinction, which in other discussions we have found it necessary to adopt, and which has more than once helped us to clear our way to a right understanding of the question in debate. A man may have a right of property in an estate, and yet a stranger may be in possession of it; and he may require to add to his right of property a right of possession, acquired by making good the former in a court of law, before the stranger is extruded, and he himself introduced into the enjoyment of the inheritance. Now, to apply this distinction to the case in hand, a right of property in the blessings of the covenant of grace is conferred by the gift and promise of God, made over to every man who hears the Gospel message addressed to him. ``And this is the record, that God has given to us eternal life, and this life is in His Son.'' This right of property in the blessings of the covenant of grace, belonging to every man, is written down in these words. The charter which every man has, bearing in it inscribed his right of property to these blessings, is the revealed Word of God. This is the first and superior title. But in itself it is incomplete, and inadequate to put him into the personal possession of his heritage. It requires to be supplemented by another title, before he can actually enjoy the salvation so made over to him by right of property, and certified by God's word and promise. To his right of property there must be added a right of possession; and this latter is obtained by means of his own personal act of faith, appropriating to himself the salvation before made over to him. The Word of God addressed to him, giving him a right of property in the blessings of the covenant, and his faith receiving that Word, giving him a right of possession, complete the full and perfect title to the blessing; and both together admit him to the enjoyment of it. There are many, who have the right of property in the covenant of grace, who never complete their title by seeking for themselves a right of possession in it. The Word of God giving the one, is not supplemented by the faith in that Word which would confer the other; and hence they are never put in actual possession of the salvation of which they are invited to partake.

Now, what the Word of God addressed to the intelligent and responsible adult is, that Baptism is when administered to the unconscious and irresponsible infant. The word of God's promise, giving a right of property in His covenant to all who hear it, cannot penetrate the silent ear, nor reach the unconscious spirit of the little child. That word cannot convey to its mind the glad tidings of its covenant right to God's grace. But is it therefore denied that right, which adults have by the hearing of the ear and the perception of the understanding, in connection with the word of promise addressed to them? Not so. If the outward word that speaks the promise of God cannot pierce to its dormant spirit,---sleeping in the germ of its moral and intellectual being,---the outward sign, that represents the promises of God, can be impressed upon it, giving to the unconscious infant, as the word gives to the intelligent adult, a right of property in the blessing of the covenant. And that is much. The infant, sprinkled with the water of that Baptism which is a sign of the covenant, has---even as the adult addressed with the word of the covenant has---a right of property in the blessings which the covenant contains; and in after life he may, by his own personal act, supplement his right of property by a right of possession obtained through faith. When the period of infancy is passed and he is a child no longer, he bears about with him, in virtue of his Baptism, a right of property in the promise of his God; and laying his hand upon that right, and pleading it with God in faith, he may add to it the right of possession, and so enter into the full enjoyment of the salvation that he requires for his soul. The written or preached Word cannot speak to the mute and insensible infant, as it speaks to the hearing ear and understanding mind of the adult, making over to him in conscious possession a right of property in the blessings of the everlasting covenant. But the little one is not thereby shut out from all interest in the covenant. The outward sign suited to his state of infancy, the outward mark impressed upon his outward person, when the significant Word were in vain addressed to his ear, have been given by God in gracious condescension to supply to him the want of that Word heard and understood. By the act of Baptism, suited and appropriate to his wholly sensitive condition of being and life, his name is put into the covenant with his God. And after years may witness the infant,---then an infant no more,---reading in faith his name there, and with the charter of his right in his hand making good his right, not of property merely, but of personal possession in all the blessings which are written in it.

Baptism, then, in the case of all baptized infants, gives them a right of property in the covenant of grace; which may in after life, by means of their personal faith, be supplemented by a right of possession, so that they shall enter into the full enjoyment of all the blessings of the covenant. The benefits of Baptism in the case of infants are not fully experienced by them until in after years they add to Baptism their personal faith, thereby really making out a complete title, not only to the property, but also to the possession of salvation. In this respect there is an obvious distinction between the Baptism of infants and the Baptism of adults. Infants are not capable of faith and repentance; and Baptism can be to infants no seal of the blessings which these stand connected with, at the time of its administration. But it may become a seal of such blessings afterwards, when the child has grown to years of intelligence, and has superinduced upon his Baptism a personal act of faith, and thereby become possessed of the salvation which he had not before. In such a case, he can look back upon his Baptism with water, administered in the days of his unconscious infancy; and through the faith that he has subsequently received, that Baptism which his own memory cannot recall, and to which his own consciousness at the time was a stranger, becomes to him a seal of his now found salvation. In adults it is otherwise; and the difference is appropriate to their condition as adults. Baptism to the believing adult is a seal at the moment of his interest in the covenant of grace; a sensible attestation of the blessings of justification and regeneration, of which at the time he is in possession, through the exercise of his faith contemporaneously with his Baptism. In the case of the adult, Baptism is a present seal in connection with the faith which he presently has. In the case of the infant, it is a prospective seal in connection with the faith which he has not at the moment, but which he may have afterwards. The full enjoyment of the benefits of the ordinance the adult experiences at the moment of its administration, in virtue of the faith which at the moment makes him a partaker in the blessings of the covenant. The full enjoyment of the benefits of the ordinance the infant cannot experience at the moment of its administration, in virtue of his incapacity of faith; but it may be experienced afterwards, when, in consequence of his newly formed faith in Christ, he too is made partaker of the covenant, and can look back in believing confidence on his former Baptism as a seal. ``The efficacy of Baptism,'' says the Confession of Faith, ``is not tied to that moment of time wherein it is administered; yet notwithstanding, by the right use of this ordinance, the grace promised is not only offered, but really exhibited and conferred by the Holy Ghost, to such (whether of age or infants) as that grace belongeth unto, according to the counsel of God's own will in His appointed time.''

\hypertarget{there-seems-to-be-reason-for-inferring-that-in-the-case-of-infants-regenerated-in-infancy-baptism-is-ordinarily-connected-with-that-regeneration.}{%
\subsection{There seems to be reason for inferring that, in the case of infants regenerated in infancy, Baptism is ordinarily connected with that regeneration.}\label{there-seems-to-be-reason-for-inferring-that-in-the-case-of-infants-regenerated-in-infancy-baptism-is-ordinarily-connected-with-that-regeneration.}}

To all infants without exception, Baptism, as we have already asserted, gives an interest in the Church of Christ as its members. To all infants without exception, Baptism, as we have also already asserted, gives a right of property in the covenant of grace, which may, by their personal faith in after life, be completed by a right of possession, so that they shall enter on the full enjoyment of all the blessings sealed to them by their previous Baptism. And beyond these two positions, in so far as infants are concerned, it is perhaps hazardous to go, in the absence of any very explicit Scripture evidence; and certainly, in going further, it were the reverse of wisdom to dogmatize. But I think that there is some reason to add to these positions the third one, which I have just announced, namely, that in the case of infants regenerated in infancy, Baptism is ordinarily connected with such regeneration. I would limit myself to the case of baptized infants regenerated in infancy,---a class of course to be distinguished broadly from baptized infants who never at any time in their lives experience a saving change; and also to be distinguished from baptized infants who experience that change, not in infancy, but in maturer years. There are these three cases, plainly to be distinguished from one another. There are, first, those infants baptized with an outward Baptism who never at any period come to know a saving change of state or nature. To such Baptism may be an ordinance giving them a place in the visible Church, and giving them also a right of property in the covenant of grace, never completed by a right of possession, and therefore given to them in vain; but it can be nothing more. There are, secondly, those infants baptized with water in infancy, but not regenerated in infancy by the Spirit of God, whose saving change of state and nature is experienced by them in after life. To such Baptism is an ordinance giving them a place in the visible Church, and giving them also a right of property in the covenant, at the moment of its administration; and in after years, when born again by the Spirit through faith, Baptism becomes to them, in addition, the seal, as it had previously been the sign, of the covenant,---their right of property having been completed by the right of possession, and the Sacrament, although long past, having become in consequence a present grace to their souls. But there are, thirdly, those infants baptized with water in infancy and also regenerated in infancy; and with regard to them I think there is reason to believe that this Baptism with water stands connected ordinarily with the Baptism of the Spirit.

That many an infant is sanctified and called by God even from its mother's womb, and undergoes, while yet incapable of faith or repentance, that blessed change of nature which is wrought by the Spirit of God, there can be no reason to doubt. There are multitudes born into this world who die ere their infancy is past,---who open their unconscious eyes upon the light only to shut them again ere they have gazed their fill,---and who, in the brief moment of their earthly being, know nothing of life save the sorrow which marks both its beginning and its close. And with regard to such infants dying in infancy, there is a blessed hope, which the Scriptures give us to entertain, that they are not lost but saved,---that they suffer, and sorrow, and die here from their interest in Adam's sin, but that, not knowing sin by their own personal act or thought, they are redeemed through their interest in Christ's righteousness. But saved though infants dying in infancy may be, yet there is no exemption, even in their case, from the universal law of God's spiritual dispensation towards men, that ``except a man be born again, he cannot see the kingdom of God.'' Within the brief hour of an infant's life, and ere the unconscious babe passes through the avenue of death into the Divine presence, must that mighty change of regeneration be undergone, which none but the Spirit of God can work; and among the rudiments of its intellectual and moral life, sleeping in the germ, there must be planted the seed of that higher life, which in heaven is destined to expand and endure through all eternity. And where, in the brief history of the young life and early death of these baptized little ones, shall we say that this mysterious work is wrought? At what moment, rather than another, is this regeneration by the Spirit accomplished? We dare not limit the free Spirit of God. The beginning of the life that comes from Him may be contemporaneous with the commencement of natural life in the infant, or it may be contemporaneous with its close. The Spirit of God is free to do His own work at His own time. But in the appointment of an ordinance to signify and represent that very work,---in the command to administer that ordinance as a sign to the little infant during the brief hour of its earthly life and ere it passes into eternity, there does seem to me some ground to believe that in such a case, of infants regenerated in infancy, the sign is meant to be connected with the thing signified,---that the moment of its Baptism is the appointed moment of its regeneration too,---and that, ordinarily, its birth by water and its birth by the Spirit of God are bound in one. It is Baptism which gives the baptized infant a right of property in the blessings of the covenant of grace; and when the infant is placed,---not from its own fault,---in such circumstances as to bar the possibility of its completing its title to those blessings by seeking through its personal faith a right of possession in them also, then it is consistent with the analogy of God's appointments in other departments of His Church, that in such extraordinary cases the absence of a right of possession should not exclude from the blessings, but that the right of property alone should avail to secure them; or in other words, that in the case of infants regenerated and dying in infancy, their Baptism should coincide with their regeneration.

I do not wish to speak dogmatically on such a question as this, when Scripture has given us so little light to enable us to read the truth with certainty. But in the particular case of infants regenerated in infancy, there does seem to be some reason to believe, that the washing with water in virtue of God's own appointment stands ordinarily connected with the renewing of nature by God's own Spirit. In the instance of believing adults, regeneration is linked inseparably with the Word believed. In connection with the Word,---although the Spirit of God is free to work without it,---He does His mysterious work of regeneration upon the adult's nature. But that Word cannot profit the little infant who is to die ere his eyes can look upon it. The Spirit of God cannot, therefore, do His gracious work of spiritual renewal and cleansing on the unconscious babe in connection with the Word believed. But there is another ordinance adapted to the infant nature, which needs to be regenerated ere it passes into another state of being. There is another ordinance, not the Word, which we are commanded to administer to the babe, incapable of receiving or profiting by the Word. There is the Baptism with water, expressive of that very regeneration which, before the little one shall pass from us to eternity, its unconscious nature must undergo. And when the infant carries with it to the tomb the sign of the covenant, administered in faith, shall we not say that with the sign, and mysteriously linked to it, there was also the thing that was signified; and that in such a case of a dying babe regenerated in infancy, the laver of Baptism was the laver of regeneration too? In the sign of the covenant thus administered to the child, and linked, as we believe, in such a case to a new and spiritual life, there is a ground of hope and consolation to a bereaved but Christian parent beyond all price. There is a joy at its birth, which none but a mother can feel, when it is said unto her that a man-child is born into the world; and there is a bitter sorrow at its early death, which none but a mother can know, when she is called upon to resign the little one whom she brought forth in sorrow, and to give it to the dust in sorrow deeper still. And when a Christian mother has been called upon thus to weep at the open grave of many of her infants, ere it close in peace upon herself, it is an unspeakable consolation for her to know, that the little one, whom she took from off her bosom to lay in the tomb, was indeed signed with the sign of a Christian Baptism; and that in its case the Baptism with water and the Baptism with the Spirit were bound up in one.

``Oh when a mother meets on high\\
The babe she lost in infancy,\\
Hath she not then for pains and fears,\\
The day of woe, the watchful night,\\
For all her sorrows, all her tears,\\
An over-payment of delight?''

\hypertarget{the-modes-of-baptism}{%
\section{The Modes Of Baptism}\label{the-modes-of-baptism}}

Before passing altogether from the subject of Baptism, it may be desirable briefly to consider the mode or modes in which the ordinance may lawfully be administered. It may seem, indeed, at first sight, a question of no great importance whether we baptize by sprinkling or by immersion,---the former being the method adopted by almost all Protestant Churches and by Western Christendom generally, the latter prevailing to a great extent in the early centuries, and still practised largely in the East. The almost unanimous opinion of orthodox theologians has always been, that Baptism in the name of the Trinity was equally valid in whichever of the two ways referred to it was administered. The position, however, taken up in our own day by many of the advocates of Baptism by immersion has given to the question an importance not properly belonging to it. The Evangelical Baptists in America, for example,---a numerous and energetic denomination,---deny the validity of Baptism by sprinkling, and declare that all persons thus baptized are living in open sin, should not be regarded as members of the Church of Christ, nor be admitted to the Lord's table. Further, they aver that the English authorized version of the Scriptures is false and unfaithful on the subject of Baptism,---purposely so, many of them add. They have issued accordingly a translation of their own with the requisite changes, and consider,---to use the words of a resolution of the Baptist American and Foreign Bible Society,---``That the nations of the earth must now look to the Baptist denomination alone for faithful translations of the Word of God.''

Our translators, in point of fact, seeing that they had to frame their version of the Bible in the very heat of a controversy about Baptism, strove carefully to stand neutral on the subject. They simply gave the Greek word an English dress; instead of βαπτιζω and βαπτισμα, they wrote ``baptize'' and ``baptism,'' thereby deciding nothing either way.

The real question at issue has been very clearly stated by President Beecher, to whose valuable work on the Mode of Baptism I would refer you for an exceedingly able and exhaustive discussion of this whole subject. ``The case,'' he says, ``is this: Christ has enjoined the performance of a duty in the command to baptize. What is the duty enjoined? or, in other words, What does the word `baptize,' in which the command is given, mean? One of two things must be true: Either it is, as to mode, generic, denoting merely the production of an effect (as purity), so that the command may be fulfilled in many ways; or it is so specific, denoting a definite mode, that it can be fulfilled in but one. To illustrate by an analogous case, Christ said: `Go, teach all nations.' Here the word go is so generic as to include all modes of going which any one may choose to adopt. If a man walks, or runs, or rides, or sails, he equally fulfils the command. On the other hand, some king or ruler, for particular reasons, might command motion by a word entirely specific, as, for example, that certain mourners should walk in a funeral procession. Now it is plain that such a command could not be fulfilled by riding or by running, for, though these are modes of going, they are not modes of walking, and the command is not to go in general, but specifically to walk.\ldots{} So likewise, when Christ said, `baptize,' He either used a word which had a generic sense, denoting the production of an effect, in any mode, such as `purify,' `cleanse;' or a specific sense, denoting a particular mode, such as `immerse,' `sprinkle,' `pour.'\,''

Now the scriptural meaning of the term βαπτιζω, I believe there is abundant evidence to show, is generic and not specific; it denotes the production of an effect which can be brought about equally well in more ways than one. The adherents of Baptist views, on the other hand, consider that the word is so specific in its signification as to fix down the lawful performance of the duty enjoined to one method only; they hold that ``in Baptism the mode is the ordinance; and if the mode is altered, the ordinance is abolished.''

The word βαπτω, from which βαπτιζω is derived, was long maintained by Dr.~Gale and other advocates of the Baptist theory to have one meaning, and only one, alike in classic, Hellenistic, and ecclesiastical Greek. It meant, they held, to immerse or dip; and it never meant anything else. This view, however, was with good reason abandoned by Dr.~Carson, probably the ablest defender of the Baptist theory in our own days. It is now very generally admitted by our opponents on this question that βαπτω has at least two meanings; first, to immerse, and second, to dye or colour. The same is true of the Latin ``tingo,'' and various similar words in other languages. It will not therefore be thought improbable that the derivative βαπτιζω should also have a primary and a secondary meaning. In point of fact, we find that, especially in later Greek, while often denoting to immerse or overwhelm, it means also, in many cases, to wash, sprinkle, cleanse. It is natural, however, to suppose that when transferred from common to ecclesiastical use, and applied in Scripture to a religious ordinance which is confessed by all parties to symbolize regeneration or spiritual purification, the meaning of the word might undergo some change. The question therefore comes to be, What is the usus loquendi of the New Testament as regards the term βαπτιζω? Looking, then, to all the passages in which the word occurs, it becomes plain, I think, that the only meaning which will carry us consistently through all of them is that of purification or cleansing. It is perfectly clear that whatever signification of the word we adopt, we must adhere to it throughout. It is quite true that βαπτιζω may have, and has, more meanings than one in ordinary Greek; but that is when it is applied to different things, and used under different circumstances. It can have but one meaning when used with respect to one definite appointment or rite, and under the same circumstances. This test can be easily applied to the various interpretations of the word in question. Take, for example, the first passage in the New Testament in which the term baptize occurs, the third chapter of Matthew, and substitute for it first the rendering which I have adopted, and then that of our Baptist brethren. It is not difficult, I think, to see which of the two best suits the whole scope of the passage: ``Then went out unto John Jerusalem, and all Judæa, and all the region round about Jordan, and were purified (immersed, or plunged) of him in Jordan, confessing their sins. But when he saw many of the Pharisees and Sadducees come to his purification (immersion, or plunging), he said, \ldots{} I indeed purify (immerse or plunge) you with water unto repentance: but He that cometh after me is mightier than I, whose shoes I am not worthy to bear: He shall purify (immerse or plunge) you with the Holy Ghost, and with fire.\ldots{} Then cometh Jesus from Galilee to Jordan unto John, to be purified (immersed or plunged) of him. But John forbade Him, saying, I have need to be purified (immersed or plunged) of Thee, and comest Thou to me? And Jesus answering said unto him, Suffer it to be so now: for thus it becometh us to fulfil all righteousness.''

That such a transition of meaning should have taken place in the case of the word βαπτιζω, appears very natural when we consider the historical circumstances connected with it. It is repeatedly used in the Septuagint, and in the works of Jewish writers who employed the Hellenistic or Alexandrian dialect, to denote the ceremonial immersions, washings, and sprinklings with water, blood, or ashes, common among the Jews. These ``divers baptisms,'' as the Apostle Paul calls them, were all practised for the sake of purification, legal or ceremonial. The two ideas, of ``baptizing'' and of ``purifying,'' were therefore constantly associated in the minds of the Jewish people; and nothing seems more natural than that in the course of time the one should pass into the other, and the words come to be used as synonymous. To recur to the history of the kindred word already alluded to: Men dipped objects in liquid in order to impart colour to them; and βαπτω came to signify ``to dye.'' The Jews immersed, or washed, or sprinkled, in order to attain purity; and so βαπτιζω came to mean ``to purify.'' In Jewish ecclesiastical language, considerably before our Lord's time, βαπτιζω seems to have dropped all reference to mode, and to have become a general term for purifying, practically equivalent to καθαριζω. A remarkable confirmation of this may be found in the third chapter of John. We are there told that a dispute had arisen between the disciples of John the Baptist and a Jew (as the true reading seems to be; not Jews as in the A. V.) ``about purifying'' (περικαθαρισμου). Now this dispute, as is shown by the context, was simply about the respective Baptisms of John and of Christ. The followers of the former were jealous on their master's behalf of the seemingly rival claims of our Lord, which had apparently been urged against them by this Jew. ``They came unto John, and said, Rabbi, He that was with thee beyond Jordan, to whom thou barest witness, behold, the same baptizeth, and all men come unto Him.'' The ``question about purifying'' was just a ``question about baptizing;'' and the Evangelist uses the words interchangeably, just because in the ecclesiastical language of his day the two meant the same thing.

The evidence by which the position which I have laid down on this subject can be still further established and strengthened, is of a cumulative sort, and for the details of it I must refer you to such works as that by Dr.~Beecher, already referred to. With respect to the apostolic practice in this matter, I am disposed to agree with the author last named, that ``it is not possible decisively to prove the mode used by the Apostles; for if going to rivers, going down to the water and up from it, etc., create a presumption in favour of immersion; so does the Baptism of three thousand on the day of Pentecost in a city where water was scarce, and of the jailor (and his household) in a prison, create a presumption in favour of sprinkling. And if a possibility of immersion can be shown in the latter cases, so can a possibility of sprinkling or pouring be shown in the former. The command being to purify, and the facts being as stated, the decided probability is, that either sprinkling, pouring, or immersion, was allowed, and Christian liberty was everywhere enjoyed. A tendency to formalism led to a misinterpretation of Paul in Rom. 6:3, 4, and Col. 2:12; and this gave the ascendency to immersion, which increased (in the postapostolic Church) until it became general, though it was not insisted on as absolutely essential on philological grounds.''

In conclusion, I remark, that many take up what appears to me a wrong ground on this question, in seeking first to prove that the word βαπτιζω, in the whole wide field in which it occurs, sometimes means to immerse, sometimes to wash, sometimes to sprinkle or pour; and then drawing from that the inference that we may lawfully baptize in any of these ways. It may be perfectly true that in profane literature the word has several meanings, but it by no means follows from that fact that, when used ecclesiastically, and applied definitely to one thing, it has more meanings than one. As employed to denote a definite religious rite, the term Baptism must have but one definite signification. And whatever we hold that to be, we must adhere to it throughout, and in all cases in which the word occurs. The true meaning of Baptism in the New Testament I believe to be purification or cleansing. That purification may be effected either by sprinkling or by immersion, according to the dictates of Christian expediency. The command to baptize is a generic command, which may be carried out in either way with equal lawfulness.

\hypertarget{the-sacrament-of-the-lords-supper}{%
\chapter{The Sacrament Of The Lord's Supper}\label{the-sacrament-of-the-lords-supper}}

\hypertarget{nature-of-the-ordinance-1}{%
\section{Nature Of The Ordinance}\label{nature-of-the-ordinance-1}}

CHRIST, as Head of His Church, has dealt out to it with a guarded hand merely outward and visible rites. In the provision which He has made for it there is enough in the way of outward and sensible ordinances for creatures made up of flesh as well as spirit to repose upon for the strengthening and confirmation of their faith; and yet not enough to convert their religion from a spiritual to a bodily service, and to transmute their faith into sight. There are but two ordinances, properly speaking, that link the Spirit with the flesh in the Christian Church; and lend the aid of a seen and sensible confirmation to an unseen and saving faith. There is one ordinance adapted to, and, it may be, specially designed for the case of infants, whose moral and intellectual life, still in the germ, lies hidden in a merely sensitive nature; and Baptism administered to the unconscious babe, whose ear cannot hear the word of salvation, becomes a visible and sensible token and seal impressed upon its flesh, of its interest in the covenant of its God. There is a second ordinance in a similar manner adapted for adults, in which an outward and sensible seal gives witness to their inward and unseen faith; and the Lord's Supper, preaching Christ by sign as well as word, is a fleshly witness, speaking to the flesh as well as to the spirit of the believer, of the blessings of the covenant of grace. There are these two, but no more than these two, outward and visible ordinances in the Church of Christ, like material buttresses, to strengthen and confirm a spiritual and immaterial faith,---the guarded and sparing acknowledgments of the fleshly nature, as well as the spiritual, which in the person of the Christian has shared in the sin, and shared also in the salvation from sin, which he knows.

We cannot doubt that a religion with these two, and neither more nor less than these two, outward rites is divinely proportioned and adapted to the need and benefit of our twofold nature, made up as it is of the fleshly and the spiritual, and both partners in the redemption, as they were formerly partners in the ruin, that belong to us. More than this in the way of the outward and sensible in the religion of Christ would have ministered all too strongly to the carnal and sensuous propensities of our nature, and would have tended towards a system which would have been ``meat and drink,'' and not ``righteousness, and peace, and joy in the Holy Ghost.'' Less than this in the way of outward and sensible ordinance would have left no room in the provision made in the Church for the adequate acknowledgment of our fleshly nature; and denied to our spiritual faith the benefit and support which it derives from some visible witness and confirmation of what it surely believes. Again, Baptism, as commonly administered to entrants into the Church, takes infeftment, so to speak, of our flesh when we enter into covenant with Christ, that not even the lower part of our being may be left without the attestation that He has redeemed it. The Lord's Supper, as administered from time to time to those who have been admitted into the Church before, renews this infeftment at intervals, and attests that the covenant by which we are Christ's still holds good both for the body and spirit which He has ransomed to Himself. The Sacrament of union to and the Sacrament of communion with Christ, tell that our very dust is precious in His sight, and has shared with the spirit in His glorious redemption. Other ordinances address themselves to the intellectual and moral nature exclusively, and speak of the care of Christ and the provision He has made for the growth and advancement of the spirit in all spiritual strength and life. The two ordinances of Baptism and the Lord's Supper, at different periods of our natural existence, and commonly in infancy and age, address themselves to both our outward and inward nature; and speak to us the testimony that both body and soul are cared for and redeemed by Christ, and that both in body and in soul we are His.

In formerly dealing with the case of Baptism as a sacramental ordinance, I endeavoured to ascertain its nature by an appeal to those marks or characteristics, in their application to Baptism, which we have found to define a Sacrament generally. Let us endeavour, by the same process, to make out the true nature and import of the Lord's Supper as a sacramental ordinance.

\hypertarget{the-first-mark-or-characteristic-of-a-sacrament-which-we-laid-down-is-that-it-be-a-divine-institute-appointed-by-christ-for-his-church.-there-is-no-religious-party-whatever-be-their-opinions-in-regard-to-the-meaning-of-the-ordinance-who-do-not-hold-the-divine-appointment-of-the-lords-supper-as-a-permanent-institution-in-the-christian-church-with-the-single-exception-of-the-quakers.-according-to-their-view-the-lords-supper-like-baptism-is-to-be-regarded-as-a-jewish-ordinance-and-the-practice-of-it-in-early-times-as-an-accommodation-to-jewish-prejudices-and-customs-but-an-ordinance-really-opposed-in-its-nature-to-the-spirituality-of-the-gospel-dispensation-and-not-intended-for-continuance-in-the-gospel-church.}{%
\subsection{The first mark or characteristic of a Sacrament which we laid down is, that it be a Divine institute appointed by Christ for His Church. There is no religious party, whatever be their opinions in regard to the meaning of the ordinance, who do not hold the Divine appointment of the Lord's Supper as a permanent institution in the Christian Church, with the single exception of the Quakers. According to their view, the Lord's Supper, like Baptism, is to be regarded as a Jewish ordinance, and the practice of it in early times as an accommodation to Jewish prejudices and customs, but an ordinance really opposed in its nature to the spirituality of the Gospel dispensation, and not intended for continuance in the Gospel Church.}\label{the-first-mark-or-characteristic-of-a-sacrament-which-we-laid-down-is-that-it-be-a-divine-institute-appointed-by-christ-for-his-church.-there-is-no-religious-party-whatever-be-their-opinions-in-regard-to-the-meaning-of-the-ordinance-who-do-not-hold-the-divine-appointment-of-the-lords-supper-as-a-permanent-institution-in-the-christian-church-with-the-single-exception-of-the-quakers.-according-to-their-view-the-lords-supper-like-baptism-is-to-be-regarded-as-a-jewish-ordinance-and-the-practice-of-it-in-early-times-as-an-accommodation-to-jewish-prejudices-and-customs-but-an-ordinance-really-opposed-in-its-nature-to-the-spirituality-of-the-gospel-dispensation-and-not-intended-for-continuance-in-the-gospel-church.}}

Now, in reference to this averment by the Quakers, it cannot be denied that, in the case of the Lord's Supper, as in the case of Baptism formerly noticed, our Lord adopted a Jewish practice or observance, and consecrated it as an ordinance in the Christian Church. The parts and ritual of the Supper are evidently derived from the observances connected with the passover as practised among the Jews. The Christian ordinance seems to be grafted upon the Jewish. We know from the Jewish accounts that we have of the passover service, that the master of the family or priest took unleavened bread, and broke it, and gave thanks to God, in much the same manner as we find it recorded of our Lord at the institution of the Supper. We know also from the same quarter, that there was one particular cup called ``the cup of blessing,'' or of ``thanksgiving,'' used at the paschal feast, of which the guests partook; and this was followed by the singing of psalms. These usages, connected with the Jewish passover, Christ adopted and accommodated to the ritual of that ordinance which we regard as the commemoration of His own death,---very much in the same manner as the washing with water employed in the Jewish baptisms or purifications was adopted and accommodated by Him to the other Sacrament which He established in the Christian Church. All this must be conceded to the Quaker theory in regard to the origin of the Christian Sacrament of the Supper. But all this, so far from making the ordinance a Jewish one, or justifying the explanation given by Quakers of the apos tolic practice of administering it, as a mere accommodation to Jewish customs or feelings, is very evidently calculated to demonstrate the reverse. The adoption of some parts of the paschal feast without the rest,---the eating bread and drinking wine as at the passover by Christians, without the slaying of the paschal lamb,---the observance of the practice at other times than once a year on the return of the anniversary of its first institution,---must, so far from being an accommodation or concession to Jewish feeling or prejudice on the part of the Apostles and first Christians, have been in reality a usage most repugnant to all the habits and prepossessions of the Israelites. The withdrawment of the outward ritual of the paschal service from the object of its original institution, and its destination to the purposes of a feast in commemoration of an event by which that service was abolished, were the very circumstances, above all others, calculated to make the ordinance not acceptable, but revolting, to Jewish feeling.

There is no truth, therefore, but the reverse, in the Quaker assumption, that the temporary continuance of the Lord's Supper in the Christian Church is to be accounted for on the theory of a concession to prejudices on the part of the Jewish converts. Add to this, that both in the statements of Scripture, and in the practice of apostolic men as recorded in Scripture, there is abundant evidence to prove that the Lord's Supper was no temporary ordinance, destined to pass away with the first merging of the Jewish into the Christian Church; but, on the contrary, was intended to be an abiding appointment for the use of its members. The command of our Lord to the disciples at the moment of the institution of the ordinance, spoke of its standing and permanent observance: ``This do in remembrance of me.'' The connection intimated by the Apostle Paul, in his account of the Supper, between the keeping of it and the second coming of Christ, evinces his opinion of the perpetual duration of the ordinance: ``As often as ye eat this bread, and drink this cup, ye do show forth the Lord's death till He come.'' The practice in the primitive Church, while under inspired direction in regard to the Lord's Supper, taken in connection with the absence of the faintest indication that it was meant for no more than a temporary purpose, is decisive evidence of the same conclusion. In short, the nature of the ordinance, as a memorial of Christ until that memorial shall be no more required on earth, in consequence of His second appearing,---the command to Jew and Gentile alike to keep the feast,---the universal practice of the Church under apostolic guidance,---and the absence of any statement express or implied in regard to the temporary character of the ordinance,---very clearly and abundantly demonstrate that the Supper of our Lord was a Divine and permanent appointment for the Church.

\hypertarget{the-next-mark-laid-down-by-us-as-characteristic-of-sacramental-ordinances-was-that-they-be-sensible-and-outward-signs-of-spiritual-truths-and-this-mark-applies-to-the-ordinance-of-the-lords-supper.}{%
\subsection{The next mark laid down by us as characteristic of sacramental ordinances, was, that they be sensible and outward signs of spiritual truths; and this mark applies to the ordinance of the Lord's Supper.}\label{the-next-mark-laid-down-by-us-as-characteristic-of-sacramental-ordinances-was-that-they-be-sensible-and-outward-signs-of-spiritual-truths-and-this-mark-applies-to-the-ordinance-of-the-lords-supper.}}

Simple and obvious although the idea be, that in the Lord's Supper we are commemorating, by appropriate and sensible images and actions, the grand spiritual truths characteristic of the Gospel, yet it is the omission or denial of this that has been the primary cause of numberless errors in regard to the nature of the ordinance. The Lord's Supper is not merely a commemoration; it is much more. But the fundamental idea which must be carried along with us in all our explanations of its nature and meaning is, that it is in the first instance a commemoration of the great truths connected with the death of Christ, as the sacrifice for the sins of His people. Nothing is easier, indeed, than to confound the sign with the thing signified; and nothing is more common in theological argument in reference to this matter. The nature and necessities of language lead us to attribute to the type what is only actually and literally true of the thing imaged or represented by the type; and in the frequent or common identification of the one with the other, we may be led not unnaturally to one or other extreme,---that of sinking the sign in the thing signified, or that of sinking the thing signified in the sign. The result is, either that we make the Sacrament to be nothing more than a sign, with no spiritual reality; or that we make it a mysterious spiritual reality, without being a sign at all. The identifying of the sign with the supernatural grace, and making them one and the same thing, must either lead to the Socinian notion that the Sacraments are nothing but symbols,---thereby evacuating the ordinance of all sacramental grace; or must lead to the Romanist or semi-Romanist notion that they are charms embodying and conveying spiritual grace, without regard to the spiritual meaning realized and appropriated by the believer in the ordinance. Hence the necessity and importance of bringing out distinctly, and laying down broadly, the character which Sacraments possess as signs of spiritual truths.

In regard to the Lord's Supper, nothing can be more distinct or conclusive than the commemorative character which is impressed upon the original institution of the ordinance by our Lord. With regard to the bread, the commandment was: ``Take, eat: this is my body broken for you: this do in remembrance of me.'' With regard to the second element in the ordinance---the cup---the appointment was no less explicit: ``This is the New Testament in my blood: this do ye, as oft as ye drink it, in remembrance of me.'' And in entire accordance with these declarations of our Lord as to the grand object of the Supper as commemorative, we have the further statement by the Apostle Paul, received by immediate revelation, as to the nature of the institution: ``For as often as ye eat this bread and drink this cup, ye do show forth the Lord's death till He come.'' In addition to all this, which very clearly exhibits the Sacrament of the Supper as in its first and most obvious character commemorative, we have the natural significance or pictorial meaning of the elements and actions in the ordinance. A rite may be in its sole or primary character commemorative in consequence of arbitrary appointment, although it may have nothing in itself naturally representative of the event commemorated. But this is not the case with the ordinance of the Communion Table. Over and above its positive institution in remembrance of the death and crucifixion of our Lord, there is a pictorial significance in the actions and elements of the Sacrament, fitted to keep constantly in view the grand and essential idea of the rite, as a rite of commemoration. The broken bread representing the broken and crucified body,---the wine poured out, the shed blood,---the eating and drinking of them, the participation in Christ's blessings to nourish the soul and make it glad,---the ``one bread'' and ``one cup,'' the communion of Christ with His people, and of them with each other,2---all these are no dumb or dark signs, but speaking and expressive of what it is intended to commemorate. This obvious characteristic of a sacramental ordinance, then, is most clearly seen in the Lord's Supper, that it is an outward and sensible sign of an inward and spiritual truth. It is the primary idea of the institution, never to be forgotten without infinite damage done to our understanding of its meaning, that, both naturally and by express Divine appointment, it is a symbolical and commemorative observance.

That the Sacrament of the Lord's Supper is an outward and sensible sign expressive of the grand and central truths connected with His death and sacrifice, is professedly held by all parties who hold that it is a Christian ordinance at all, and consider it to be binding upon Christians. And yet, notwithstanding of this professed and apparent unanimity upon the point, there is one religious denomination whose principles amount to a denial of this simple truth; and who virtually and really make the Lord's Supper to be not a sign, and not a commemorative ordinance at all,---thereby denying to it the proper character of a Sacrament. I allude to the Church of Rome. I do not mean to enter upon a consideration of the doctrine of that Church with regard to the Lord's Supper at present---for I intend to take up that subject afterwards,---but it may be not unsuitable or unimportant, meanwhile, to remark, that many of the errors of the Church of Rome in regard to this Sacrament are to be traced back to the neglect or denial of the simple but fundamental truth, that in its primary and essential character the Lord's Supper is a commemorative ordinance,---a remembrance of a sacrifice, and not a sacrifice itself,---a memorial of the great atonement and offering up of Christ on the Cross, and not a repetition of that atonement. By the doctrine of transubstantiation held by the Church of Rome, the elements of bread and wine are asserted to be changed into the actual body and blood of Christ, the Son of God; so that the use of these elements in the Sacrament is not to represent, but to repeat or continue the offering once made for sinners upon the Cross. The sign is identified with the thing signified; the symbol, instead of remaining a symbol, becomes one and the same with what was symbolized; the image and the reality are not two separate and independent things, but are confounded together. This is the unavoidable consequence of the doctrine of transubstantiation held in regard to the communion elements. The bread in the ordinance ceases to be the sensible sign of the Lord's body, and actually becomes that body; the wine in the cup ceases to be the representation symbolically of the blood of the Lord, and is transmuted into that very blood. There is no separating idea which continues to divide the symbol from the reality represented. The two are lost in one. The grand and fundamental characteristic of a Sacrament---that it is the outward and sensible sign of an inward and spiritual truth---is utterly forgotten or denied; and the consequence is the subversion of every idea essential to a Sacrament. While professedly, in some sort of way not easily understood, the Church of Rome holds that the Lord's Supper is a commemorative Sacrament, it in reality does away with the fundamental characteristic of a Sacrament as a sensible sign of spiritual truth.2

\hypertarget{the-third-mark-laid-down-by-us-as-characteristic-of-sacramental-ordinances-is-that-they-are-the-seals-of-a-federal-transaction-between-the-believer-and-christ-through-means-of-the-ordinance-and-this-mark-is-applicable-to-the-lords-supper.}{%
\subsection{The third mark laid down by us as characteristic of sacramental ordinances, is, that they are the seals of a federal transaction between the believer and Christ through means of the ordinance; and this mark is applicable to the Lord's Supper.}\label{the-third-mark-laid-down-by-us-as-characteristic-of-sacramental-ordinances-is-that-they-are-the-seals-of-a-federal-transaction-between-the-believer-and-christ-through-means-of-the-ordinance-and-this-mark-is-applicable-to-the-lords-supper.}}

There are not a few who rest contented with the position already laid down in regard to the Lord's Supper, and restrict themselves to the view which makes it a sensible sign of spiritual truth. At the date of the Reformation the subject of the Lord's Supper was very keenly canvassed amongst the Protestant Churches; and the Sacramentarian controversy, or the dispute as to the true meaning and nature of the Lord's Supper, went further than any other to divide the opinions of the early Reformers. While Luther held views approximating to those of the Church of Rome on this subject, although denying the doctrine of transubstantiation, there was another party among the first Reformers, especially in Switzerland, headed by Zwingli, who advocated principles differing very widely from those of Luther. Zwingli, the chief founder of the Protestant Churches in Switzerland, and the predecessor of Calvin in the Swiss Reformation, is not uncommonly regarded as the originator of those views of the Lord's Supper which represent it as a symbolical action commemorative of the death of Christ, and as nothing more than this. There seems to be good ground to question this opinion, and to doubt whether Zwingli ever really meant to deny that the Lord's Supper is a seal, as well as a sign of spiritual grace,---the outward voucher as well as representation of a spiritual and federal transaction between the believer and Christ through means of the ordinance. Under the strong reaction then felt from the views of the Lord's Supper entertained by the Church of Rome, which virtually set aside and denied the symbolical character of the ordinance, and superseded the outward sign by the thing signified, Zwingli and others felt that the true source of the doctrine of transubstantiation was the denial of the primary character of the ordinance as a commemorative sign, and the making the symbol give place to the reality symbolized under it. In other words, Zwingli and his associates in Switzerland held that the root of the evil lay in denying that the bread and wine in the Lord's Supper were signs, and constituting them the thing signified,---the very body and blood of the Lord. And in bringing out this principle as against the dogma of transubstantiation, they were led in their argument to speak somewhat unguardedly, as if, while Scripture represented the Sacrament as symbolical, it did not represent it as anything more than symbolical. Notwithstanding the violent controversy which the opinions of Zwingli and his followers excited, and the opposition they encountered from Luther and others of the German section of the Reformation, it is very doubtful indeed whether their opinion really excluded or denied the idea of a seal of a federal transaction, as well as a sign, as really belonging to the character of the Lord's Supper. However this may be, it was reserved for the successor of Zwingli, as the leader in the Swiss Reformation, to bring out from Scripture, and to establish on its true foundation, the proper notion of the Lord's Supper as more generally entertained by Protestant Churches since his time; and it is not the least of the many debts due by the Church to the illustrious Calvin, that we owe to him the first full and accurate development and decided maintenance of the true doctrine of the ordinance, as neither a sign alone, nor yet the thing signified alone,---as neither an empty symbol, nor yet the transubstantiated body and blood of Christ,---but as a sign and, at the same time, a seal of spiritual and covenant blessings, made over in the ordinance to the believer. The doctrine of the Sacrament of the Lord's Supper as a sign or symbol, and nothing more, has become the characteristic system of the Socinian party. More recently still, it has become the theory of not a few of the Independent body in England, as represented by Dr.~Halley.

That the Lord's Supper, in addition to being a sign, is also a seal of a federal transaction, in which the believer through the ordinance makes himself over to Christ, and Christ makes Himself over with His blessings to the believer, may be satisfactorily evinced from a brief review of the statements of Scripture on the subject. There are four different occasions on which the Lord's Supper is more especially referred to in Scripture; and from the statements made in regard to it on these occasions, it may be conclusively proved that much more is attributed to the ordinance than merely the character of a sign.

1st, There is the description given of the nature and meaning of the ordinance in connection with the history of its institution, as given by the different evangelists, and educed from a comparison of them, which seems not indistinctly to intimate that the Lord's Supper is more than a commemorative sign. In the words of the institution, our Lord calls the cup ``the New Testament or covenant in His blood,''---language which can be interpreted, and apparently requires to be interpreted, so as to assert a more intimate connection than any between a symbol and the thing signified, between the cup drunk in the Supper and the covenant of grace which secures the blessings represented. Add to this, that our Lord asserts the bread to be His body, and the wine to be His blood, in such terms as certainly imply that the one was a sign of the other, but apparently imply more than this,---the words seeming to intimate a sacredness in the symbols more than could belong to mere outward signs, and unavoidably suggesting a more intimate relationship between the elements of the ordinance and the spiritual blessings represented,---even such a connection as that which would make the use of the one by the worthy receiver stand connected with the actual enjoyment spiritually of the other.

2d, There is a separate account of the institution of the Lord's Supper given by the Apostle Paul in the 11th chapter of 1st Corinthians, in which the intimacy and sacredness of the connection between the symbols of the ordinance and the blessings represented are still more strongly brought out. The ``eating and drinking unworthily'' is represented as the sin of being ``guilty of the body and blood of the Lord;'' a second time it is spoken of by the apostle as the guilt on the part of the unworthy participator of ``eating and drinking judgment to himself,''---the reason assigned for the heinousness of the offence being, that he ``has not discerned the Lord's body;'' and, as a precaution against the danger of such transgression, a man is commanded to ``examine himself'' before he partake of the Supper. It seems impossible, with any show of reason, to assert that the ``discernment'' (διακρισις) here spoken of is the mere power of interpreting the signs as representative of Christ's death; or that the ``guilt'' incurred is nothing more than the danger of abusing certain outward symbols; or that the ``examination'' enjoined is no more than an inquiry into one's knowledge of the meaning of the commemorative rite. All these expressions evidently point to a spiritual discernment and participation by the believer, not of the sign, but of the blessing signified; and to a spiritual and awful sin, not of misusing and profaning outward symbols, but of misusing and profaning Christ actually present in them.

3d, There is a brief but most emphatic reference to the Lord's Supper in the 10th chapter of 1st Corinthians, which can be interpreted upon no principle which limits the meaning of the ordinance to a mere sign, but which very plainly asserts a federal transaction between the believer and Christ in the ordinance, and the communication through the ordinance of spiritual blessings. ``I speak as to wise men,'' says the apostle; ``judge ye what I say. The cup of blessing which we bless, is it not the communion of the blood of Christ? The bread which we break, is it not the communion of the body of Christ?'' The κοινωνια---the communion, or participation, or interchange, or mutual fellowship of the blood of Christ and the body of Christ---cannot possibly be understood of the mere signs of the body and blood, without a very violent experiment practised on the language of the apostle. And if ``the fellowship'' does not refer to the outward symbol, it can only refer to the spiritual blessings represented in the ordinance,---to Christ Himself present after a spiritual manner in the Sacrament, and giving Himself to the believer, while the believer gives himself to Christ, so as to establish a true κοινωνια, or fellowship, or communion between them. It is hardly possible with any plausibility to interpret the language of the apostle in any other way than as expressive of a federal transaction between the believer and Christ in the ordinance.

4th, There is a lengthened discourse in the 6th chapter of the Gospel by John, in which our Lord indeed makes no express reference to the Supper by name, but which it is hardly possible, I think, to avoid applying in its spiritual meaning to the ordinance. In that discourse our Saviour declares Himself to the Jews to be ``the bread of life which came down from heaven;'' He tells them that ``except they eat the flesh and drink the blood of the Son of man, they have no life in them;'' He asserts that ``His flesh is meat indeed, and His blood drink indeed;'' and He affirms that ``He that eateth my flesh, and drinketh my blood, dwelleth in me, and I in him.'' Whether this discourse refers directly and expressly to the ordinance of the Lord's Supper or not, it is quite plain that it affords, by the parallelism of the language employed to that used in connection with the ordinance, a key to interpret the sacramental phraseology applied to the Supper. It very plainly points to a spiritual eating and drinking of the flesh and blood of the Son of God, and a spiritual participation, far beyond a mere fellowship in an outward and empty symbol.4

On such grounds as these, we hold that the theory which explains the Sacrament of the Supper to be no more than a commemorative sign comes very far short of the Scripture representations of the ordinance; and that nothing but the idea of a seal of a federal transaction between the believer and Christ in the Sacrament will come up to the full import of the observance.

\hypertarget{the-fourth-and-last-mark-laid-down-by-us-as-characteristic-of-a-sacramental-ordinance-is-that-it-is-a-means-of-grace-and-this-mark-also-applies-to-the-ordinance-of-the-lords-supper.}{%
\subsection{The fourth and last mark laid down by us as characteristic of a sacramental ordinance, is, that it is a means of grace; and this mark also applies to the ordinance of the Lord's Supper.}\label{the-fourth-and-last-mark-laid-down-by-us-as-characteristic-of-a-sacramental-ordinance-is-that-it-is-a-means-of-grace-and-this-mark-also-applies-to-the-ordinance-of-the-lords-supper.}}

After what has been said, it is not necessary to do more than lay down this position. As the sign and seal of a federal transaction between the believer and Christ, it is plain that it must be the means of grace to his soul. It presupposes, indeed, the existence of saving grace on the part of the participator in the ordinance; it is a seal to him of the covenant actually and previously realized and appropriated by him; but, as a seal, it is fitted to add to the grace previously enjoyed, and to impart yet higher and further blessing. What is the manner in which this grace is imparted; how the Sacrament of the Supper becomes a living virtue in the heart of the participator; what is the efficacy of the ordinance,---these are questions the consideration of which opens up to us those further discussions to which we have next to address ourselves. While we believe that the Sacrament of the Supper is an eminent and effectual means of grace, as a seal of the covenant transaction represented in the ordinance, and through the faith of the participator, Romanists and semi-Romanists attribute to the ordinance a character and an efficacy which we believe that Scripture does not sanction, but, on the contrary, disowns. To the unscriptural views of the Supper held by the Church of Rome we shall now turn our attention.

\hypertarget{transubstantiation}{%
\section{Transubstantiation}\label{transubstantiation}}

Both the Lord's Supper and Baptism are Divine appointments of perpetual authority in the Christian Church. Both are outward and sensible signs, expressive of spiritual truths; both are seals of a federal transaction between Christ and the believer in the ordinance; and both, while they presuppose the existence of grace on the part of the receiver, are at the same time the means, by the Spirit, and through the believer's faith, of adding to that grace, and imparting a fresh spiritual blessing. And thus, parallel as the Sacraments of the Christian Church are in their nature and efficacy, they are alike also in the misapprehensions to which they have been exposed. Baptism has been misrepresented as an ordinance possessed in itself of an independent and supernatural virtue, apart from the spiritual state or disposition of the participator, so that, ex opere operato, it infallibly communicates saving grace to the soul. And, in like manner, the Sacrament of the Lord's Supper has been misrepresented as an ordinance embodying in itself a spiritual power, and efficacious of itself to impart saving grace. The full-grown and legitimate development of these views in regard to the Lord's Supper is to be found in the principles of the Church of Rome, and in the doctrine which she propounds under the name of transubstantiation.

The Romish system of belief and instruction in regard to the ordinance of the Supper is briefly this. At the original institution of the ordinance, it is believed by the Church of Rome that our Lord, by an exertion of His almighty power, changed miraculously the bread and wine into His body and blood, His human soul and His Divine Godhead; that this supernatural change was effected in connection with the words of institution uttered by Him: ``This is my body; this is my blood;'' that in giving the appearance of ordinary elements into the hands of His Apostles, He actually gave Himself, including both His humanity and His Divinity; and that they really received and ate His flesh, and drank His blood, with all their accompanying blessings to their souls. And what was thus done in a supernatural manner by Christ Himself at the first institution of the ordinance, is repeated in a manner no less supernatural every time the Lord's Supper is administered by a priest of Rome with a good intention. The priest stands in the place of Christ, with an office and power similar to Christ's, in every case in which he dispenses the Supper; the words of institution repeated by the lips of the priest are accompanied or followed by the same supernatural change as took place at first; the substance of the bread and wine used in the ordinance is annihilated, while the properties of bread and wine remain. In place of the substance of the natural elements, the substance of Christ in His human and Divine nature is truly present, although under all the outward attributes of bread and wine; and those who receive what the priest has thus miraculously transubstantiated are actual partakers of whole Christ, under the appearance of the ordinary sacramental elements.

Under this fearful and blasphemous system there are properly two grand and fundamental errors from which the rest flow; and which it is important to mark and deal with separately, although they are intimately connected, and form part of the same revolting theory of the Sacrament. There is, first of all, that supernatural change alleged to be wrought upon the elements by the authority of the priest in uttering the words of institution,---the transubstantiation properly so called,---by which the bread and wine become not a sign or symbol, but the actual substance of the crucified Saviour; and there is, secondly, and in consequence of such transubstantiation, the making of the elements not the signs of Christ's sacrifice, but the reality of it,---the bread and wine having become Christ Himself, and the priest having, in so transubstantiating them, actually made the sacrifice of the Cross once more, and offered it to God. These two doctrines of real transubstantiation, and a real sacrifice in the ordinance of the Supper, are both avowed as fundamental in the theory of the Church of Rome; and from these two doctrines all the others connected with the subject are derived. First, From the doctrine of the transubstantiation of the elements into the actual humanity and Divinity of the Lord Jesus Christ, there very obviously, and perhaps not unnaturally, follows that other doctrine, which declares that the elements are proper objects for the worship of Christians; and hence we have the elevation and adoration of the Host in connection with the Romanist doctrine of the Supper. Second, From the doctrine that the elements, transubstantiated into a crucified Saviour, become a real sacrifice, and a true repetition or continuation of the offering made upon the Cross, there very obviously and naturally follows that other doctrine, which teaches that the ordinance procures for the participator in it atonement and forgiveness of sin; and hence we have the saving grace infallibly communicated by the Sacrament wherever there is a priest to dispense it, or a soul to be saved by the participation of it. We shall consider, then, the doctrine of the Church of Rome in connection with the Supper, under the twofold aspect of the real transubstantiation alleged to pass upon the elements, and the real sacrifice alleged to be offered in the ordinance. These two points form the grand and essential features of the Romanist theory of this Sacrament; and, separately discussed, will enable us to review all that is of chief importance connected with it.

The doctrine of transubstantiation is thus laid down in the Canons of the Council of Trent: ``If any shall deny that in the Sacrament of the most holy Eucharist there is contained truly, really, and substantially the body and blood, together with the soul and Divinity of our Lord Jesus Christ, and so whole Christ, but shall say that He is only in it in sign, or figure, or virtue, let him be accursed.'' ``If any shall say that in the Holy Sacrament of the Eucharist there remains the substance of bread and wine, together with the body and blood of our Lord Jesus Christ, and shall deny that wonderful and singular conversion of the whole substance of the bread into the body, and of the whole substance of the wine into the blood, while only the appearances (species) of bread and wine remain---which conversion the Catholic Church most aptly styles transubstantiation,---let him be accursed.'' ``If any shall say that Christ, as exhibited in the Eucharist, is only spiritually eaten, and not also sacramentally and really, let him be accursed.''

This monstrous and audacious perversion of the doctrine of Scripture by the Church of Rome is founded upon and defended by an appeal to the literal meaning of the words of Scripture in speaking of the ordinance, in contradistinction to the figurative meaning of them. It is on this literal sense of the Scripture language that the only argument of Romanists in support of their system is built; and, over and above an appeal to the bare literalities of the expressions employed, there is not the shadow of a reason that can be alleged in defence of it. ``It is impossible for me,'' says Cardinal Wiseman in his Lectures on the Principal Doctrines and Practices of the Catholic Church,---``it is impossible for me, by any commentary or paraphrase that I can make, to render our Saviour's words more explicit, or reduce them to a form more completely expressing the Catholic doctrine than they do themselves: `This is my body; this is my blood.' The Catholic doctrine teaches that it was Christ's body, that it was Christ's blood. It would consequently appear as though all we had here to do were simply and exclusively to rest at once on these words, and leave to others to show reason why we should depart from the literal interpretation which we give them.'' Since Romanists, then, take up their position in defence of transubstantiation substantiation on the literal construction of the words employed in reference to the ordinance, and on that alone, what is material or essential to the argument is brought within a very narrow compass indeed. That argument may be, and indeed often is, encumbered with much irrelevant matter. But the main and only essential point to be discussed is simply this: Are we bound to interpret the Scripture phraseology employed in connection with the Lord's Supper in a literal sense, as affirming that the true body and blood of Christ are given in the ordinance; or, do the very terms of that phraseology, and the nature of the thing spoken of, compel us to adopt not a literal, but a figurative interpretation? This is evidently the status quœstionis between the Romanists and their adversaries in reference to the debate about transubstantiation. Romanists never pretend to bring any argument in aid of their theory of the Supper, except the argument of the literal meaning of the sacramental words. This disposed of, there is no other in the least available to defend their position. Is it, then, possible to adopt a literal interpretation of the words which Scripture employs to describe the sacramental elements? Is it competent to adopt a figurative interpretation? Is it necessary to adopt a figurative interpretation? These three questions, fairly answered, will embrace the whole controversy necessary to the discussion of the Romanist dogma of transubstantiation.

\hypertarget{it-is-impossible-to-adopt-a-literal-interpretation-of-the-sacramental-phraseology-and-this-is-evinced-by-romanists-themselves-in-their-own-departure-from-it-in-the-very-matter-under-discussion.}{%
\subsection{It is impossible to adopt a literal interpretation of the sacramental phraseology; and this is evinced by Romanists themselves, in their own departure from it in the very matter under discussion.}\label{it-is-impossible-to-adopt-a-literal-interpretation-of-the-sacramental-phraseology-and-this-is-evinced-by-romanists-themselves-in-their-own-departure-from-it-in-the-very-matter-under-discussion.}}

The principle of a strictly literal interpretation of the sacramental language of Scripture is the only principle which furnishes a single plea in favour of the dogma of transubstantiation; and yet the necessities of the language employed compel Romanists to surrender that principle in its application to the very case in which they demand that we shall observe it. The advocate of transubstantiation, by his own practice in the very matter in hand, nullifies his own solitary argument. He demands from us a literal rendering of the Scripture language; and yet in the very same passage of Scripture he is himself forced to adopt a non-literal. Take the words of Luke as he records the first institution of the Supper, and we see at once that in these the Romanist is forced again and again to abandon a literal, and have recourse to a figurative interpretation. ``And He took the cup,'' says the evangelist, describing our Lord's action, ``and gave thanks, and said, Take this, and divide it among yourselves.'' According to the strictly literal method of interpretation advocated and demanded by the Romanist, it was the cup, and not the wine in the cup, that was to be taken and shared by the disciples; and the Romanist is obliged to adopt the non-literal rendering in this case to suit his views of what occurred. Again, we find the inspired historian saying, in reference to what our Lord did, ``Likewise also the cup after supper, saying, This cup is the New Testament in my blood,''---language which once more demands that the Romanist shall surrender his literal, and have recourse to a non-literal interpretation, so that he may not identify the vessel in which the wine was contained with the New Covenant, nor transubstantiate the cup into a covenant, but make the one merely a sign or symbol of the other by a figurative use of the language. Once more, the Romanist departs from his principle of a literal interpretation, when the evangelist tells us that Christ spoke of His blood ``which is shed for you.'' At the moment of the utterance of these words, the shedding of His blood was a future event, to happen some hours afterwards, and not a present one, as the words literally rendered would assert; and, accordingly, the Romanist has no scruple in interpreting it in a non-literal sense, as indeed he is forced to do by the very necessity of the language. Or, take the words of the Apostle Paul in his account of the ordinance of the Supper, which he had, separately from the evangelists, himself received of the Lord. Here, again, we have the same use of terms which no literal interpretation will enable even the Romanist to explain. The apostle, like the evangelist, tells us that the words of our Lord were expressly, ``This cup is the New Testament in my blood,''---language which, interpreted upon the principle of strict literality, would identify the vessel containing the wine with the Divine covenant, and which requires, therefore, even in the opinion of the Romanist, to be understood figuratively. And, further still, the apostle, after the giving of thanks by our Lord, still speaks of the elements, not in language which denotes their transubstantiation, but in terms which plainly declare that they were bread and wine still. ``For as often as ye eat this bread, and drink this cup, ye do shew the Lord's death till He come.'' In this case no literal rendering of these words will be sufficient to reconcile them with the dogma of transubstantiation; and even in supporting that dogma, the Romanist is compelled in this passage to fall back upon an interpretation not literal. We are warranted, then, by the practice of Romanists themselves, in the very case of the sacramental language employed in Scripture, to say that it is not possible to adhere to, or consistently to carry out, a strictly literal interpretation.3

\hypertarget{a-figurative-interpretation-of-the-sacramental-language-is-perfectly-competent-and-possible.}{%
\subsection{A figurative interpretation of the sacramental language is perfectly competent and possible.}\label{a-figurative-interpretation-of-the-sacramental-language-is-perfectly-competent-and-possible.}}

It cannot be denied---and we have no occasion or wish to deny it---that, as a general canon of interpretation, it is true that the literal rendering of any statement made by a writer ought, in the first instance, to be tried and to be adopted, if it be in accordance with the use of words and the import and object of the statement. But the necessities and use of language justify and demand a figurative interpretation of terms, rather than a literal, in manifold instances; and those instances in which words are to be rendered not literally, but figuratively, must plainly be determined by the nature, connection, and object of the words. Now, in reference to the use of the sacramental language found in the Bible, it has often been argued, and has never yet been fairly met by the advocates of a literal meaning, that many similar passages are to be found in Scripture in which the same words admit of, and indeed require, not a literal, but a figurative interpretation, by the confession of all parties; and the conclusion is drawn from this, and fairly drawn, that the terms used in regard to the ordinance of the Supper may be figurative too. The occurrence of such texts, demanding, as all parties allow, a figurative or non-literal rendering, is valid and relevant evidence in regard to the nature of Scripture language, and proves at least this, that the words employed in reference to the Supper may admit of a figurative rendering also. This citation of parallel language does not in itself, indeed, demonstrate that the sacramental terms must be figurative; but it unquestionably proves that they may be figurative. Cardinal Wiseman, in his discussion of the doctrine of transubstantiation, gives a list of some texts bearing on the question, which have been referred to by Protestants as evidence in their favour, to the effect that the language, ``This is my body,'' ``this is my blood,'' may be understood, not literally, but figuratively. They are to the following effect:

``The seven good kine are seven years.''\\
``The ten horns are ten kings.''\\
``The field is the world.''\\
``And that rock was Christ.''\\
``For these are the two covenants.''\\
``The seven stars are the angels of the seven churches.''\\
``I am the door.''\\
``I am the true vine.''\\
``This is my covenant between me and you.''\\
``It is the Lord's passover.''

In these instances, and many similar ones, it is admitted by all parties, Romanists as well as Protestants, that the verb to be must be understood in its non-literal signification, and cannot by any possibility be understood literally. From the nature of the assertion made, from the context, and from the manner in which the terms are made use of, there is no possibility of denying that these texts are to be understood not literally, but figuratively; and they seem, therefore, by this parallelism to the words employed in connection with the Supper, to prove all that they were ever quoted to prove, namely, that the expressions, ``This is my body,'' ``this is my blood,'' may be understood in a figurative sense too. Such texts are not quoted to demonstrate that the sacramental phraseology of Scripture must be figurative; they are only quoted to prove that there is nothing in the nature of Scripture language, judging by its use in similar cases, to prevent us, if the nature of the statement and the context should require it, from interpreting the language concerning the Supper in a non-literal or figurative sense also. The multitude of texts closely analogous in form to the phrases, ``This is my body,'' ``this is my blood,'' and which, as all parties allow, must be understood figuratively, may not indeed, taken singly, necessitate a non-literal rendering in the latter case also; but they, at the very least, authorize it, should the import and connection of the passage make the demand, if they do not go a step further, and of themselves recommend a figurative interpretation.

Now, how is it that Cardinal Wiseman in his Lectures deals with these passages, and disposes of the argument drawn from them? He bestows a vast deal of minute criticism upon them, in order to show that these passages must, either from the meaning of the statement made in each, or the sense of the context, or the express assertion of the sacred writer, be accounted figurative and symbolical; and that, therefore, the verb to be in each of these cases must be reckoned equivalent to the verb to signify. And having done this, he considers he has done enough to prove that the cases referred to are not parallel to the sacramental language, ``This is my body,'' ``this is my blood.'' Now, it is enough, in reference to such an argument, to say that we willingly adopt his explanation of these passages, accounting them, as he does, to be figurative, and reckoning, as he does, the verb to be, when employed in such texts, as equivalent to the verb to signify. And it is for this very reason that we quote them as a justification of our assertion, that the same verb, when employed in reference to the Lord's Supper, may be equivalent there also to the verb to signify. If these texts did not admit of a figurative interpretation, and if the verb to be did not in them appear equivalent to the verb to signify, we should not have quoted them, because they would not have served our purpose. The reasoning of the Cardinal is certainly a singular specimen of an attempt at logical argument. I shall give it in his own words: ``Suppose,'' says he in his Lectures, ``suppose I wish to illustrate one of these passages by another, I should say this text, `The seven kine are seven years,' is parallel with `The field is the world,' and both of them with the phrase, `These are the two covenants;' and I can illustrate them by one another. And why? Because in every one of them the same thing exists; that is to say, in every one of these passages there is the interpretation of an allegorical teaching,---a vision in the one, a parable in the second, and an allegory in the third. I do not put them into one class because they all contain the verb to be, but because they all contain the same thing. They speak of something mystical and typical,---the interpretation of a dream, an allegory, and a parable. Therefore, having ascertained that in one of these the verb to be means to represent, I conclude that it has the same sense in the others; and I frame a general rule, that wherever such symbolical teaching occurs, these verbs are synonymous. When, therefore, you tell me that `this is my body' may mean `this represents my body,' because in those passages the same word occurs with this sense, I must, in like manner, ascertain not only that the word to be is common to the text, but that the same thing is to be found in it as in them; in other words, that in the forms of institution there was given the explanation of some symbol, such as the interpretation of a vision, a parable, or a prophecy \ldots{} Until you have done this, you have no right to consider them all as parallel, or to interpret it by them.''

The objection here urged by Cardinal Wiseman seems to amount to this, that we have quoted passages which, by the nature of the statement they contain, or by the context, or by the direct assertion of the writer, are plainly demonstrated to be figurative, while the sacramental expressions, ``This is my body,'' ``this is my blood,'' are not so demonstrated to be figurative. The answer is obvious. We do not quote such texts to prove that the terms of the sacramental institution must be understood figuratively, but to prove that they may be understood figuratively; to demonstrate that there is no bar in the shape of Scripture usage in the way to prevent us from interpreting them figuratively, if it is necessary. We are prepared to prove, by the very same means as the Cardinal employs,---by the nature of the statement itself, by the context, and such like considerations,---that the sacramental terms are figurative, just as Cardinal Wiseman proves that the words, ``This cup is the New Testament,'' are to be understood figuratively, or as these other terms, ``The seven kine are seven years,'' must be interpreted figuratively. The very nature of the statement itself proves it to be a statement to be understood, not in a literal, but a figurative sense. We interpret the expression, ``The seven kine are seven years,'' in a figurative sense, not because these words occur in the interpretation of a dream,---for both the dream and the interpretation may be embodied in words, literal, and not figurative,---but because the very nature of the proposition and the sense of the context necessitate it, it being impossible that the seven kine can be literally seven years. Again, we interpret, and so does Cardinal Wiseman, the expression, ``This cup is the New Testament,'' not literally, but figuratively, for a similar reason,---that the very nature of the proposition, and the sense of the context, demand a non-literal rendering; and in like manner we interpret the expression, ``This is my body,'' ``this is my blood,'' not literally, but figuratively, for the very same reason, because the very nature of the proposition, and the sense of the context, necessitate such an interpretation. The citation of other passages of Scripture in which the verb to be is used for the verb to represent or signify, is had recourse to in the argument simply to prove that the usage of Scripture language does not forbid, but countenances such a kind of interpretation. And the numerous texts already referred to are both relevant and sufficient to accomplish that object.

\hypertarget{a-figurative-interpretation-of-the-sacramental-language-this-is-my-body-this-is-my-blood-is-not-only-possible-and-competent-but-necessary.}{%
\subsection{A figurative interpretation of the sacramental language, ``This is my body,'' ``this is my blood,'' is not only possible and competent, but necessary.}\label{a-figurative-interpretation-of-the-sacramental-language-this-is-my-body-this-is-my-blood-is-not-only-possible-and-competent-but-necessary.}}

In no other way can we ever discriminate between figurative and literal terms, whether scriptural or non-scriptural, whether used by inspired or uninspired men, than by a reference to the nature of the proposition which the language embodies, to the sense of the context, and to the object of the speaker or writer; unless in those exceptional cases in which he directly tells us that he is to be understood in the one way or in the other. Very seldom indeed, in regard to language not meant to deceive, is it difficult to understand, from a consideration of these points, whether it is to be interpreted figuratively or not. In the case of the Lord's Supper, the words employed in reference to the elements could have presented to the disciples who heard them no difficulty at all. The ordinance was grafted upon the passover, with the figurative language and actions of which the Apostles, as Jews, were abundantly familiar; and this circumstance alone must have familiarized their minds with, and prepared them for the figurative meaning of the words and elements in the Supper. Above all, the nature of the proposition, ``This is my body,'' ``this is my blood,'' interpreted by the commentary of our Lord, ``This do in remembrance of me,'' and understood in the light of His accompanying actions and words, renders it nearly impossible that they could believe that a miracle had been wrought on the bread and wine, and that the body and blood, soul and Divinity of the Lord Jesus Christ, then present to their eyes, could be at the same instant contained under the appearance of the morsel of bread and the mouthful of wine that they ate and drank. Nothing but the ``strong delusion that believes a lie'' can lead any man who reads and understands the simple narrative of Scripture, to deny that the interpretation of the sacramental phraseology employed must be figurative, and not literal.

There are two attempts commonly made by Romanists to explain away the impossibility of the Apostles,---or indeed any other man not wholly blinded by spiritual delusion,---believing in the literal interpretation of the sacramental words that refer to the Supper.

1st, The power of Christ to work a miracle, like that which is alleged to have been wrought in the case of the bread and wine, is asserted; and it is averred that the Apostles could not doubt the supernatural ability of their Lord and Master, so often in other days exerted before their eyes. ``What,'' asks Dr.~Wiseman, ``is possible or impossible to God? What is contradictory to His power? Who shall venture to define it further than what may be the obvious, the first, and simplest principle of contradiction,---the existence and simultaneous non-existence of a thing? But who will pretend to say that any ordinary mind would be able to measure this perplexed subject, and to reason thus: `The Almighty may indeed, for instance, change water into wine, but He cannot change bread into a body?' Who that looks on these two propositions with the eye of an uneducated man, could say that in his mind there was a broad distinction between them, that while he saw one effected by the power of a Being believed by him to be omnipotent, he still held the other to be of a class so widely different as to venture to pronounce it absolutely impossible?.\ldots{} Now, such as I have described were the minds of the Apostles,---those of illiterate, uncultivated men. They had been accustomed to see Christ perform the most extraordinary works. They had seen Him walking on the water, His body consequently deprived for a time of the usual properties of matter,---of that gravity which, according to the laws of nature, should have caused it to sink. They had seen Him, by His simple word, command the elements and raise the dead to life, etc. Can we, then, believe that with such minds as these, and with such evidences, the Apostles were likely to have words addressed to them by our Saviour, which they were to interpret rightly, only by the reasoning of our opponents,---that is, on the ground of what He asserted being philosophically impossible?''

It is hardly necessary to reply to such an argument as this. In the first place, the miracles with which the Apostles were familiar had no analogy whatsoever to the stupendous wonder of transubstantiation. Those miracles were appeals to the senses in proof of truths not seen; and they were tested by the senses, as things to be judged of by them all. The so-called miracle of transubstantiation is no appeal to the senses, but the reverse,---a thing not to be tested by the exercise of any one of them, if it were possible, and a thing denied by any one of them, because impossible. If it were a possible thing, it would subvert the very principle on which our perceptions are made to us by God the primary source of our beliefs, and the foundation of truth to us; and it would cause the very instincts which His hand has laid deep within our inmost being to be to us a lie. The conversion of water into wine at that marriage supper in Cana of Galilee of old was a wonder seen by the eye, and in agreement with the evidence of the senses, because the properties, first of the water, and afterwards of the wine, were seen and judged of by all. The conversion of the bread into the body of the Lord, while yet the properties of bread remain, is a wonder that contradicts the evidence of our senses, and involves an impossibility.

In the second place, even Cardinal Wiseman himself admits that there are impossibilities in the nature of things, not competent even for Almighty power to accomplish. Such an impossibility, according to his own statement, is the ``existence and simultaneous non-existence of a thing;'' and side by side with this one limitation, which, upon the authority of Dr.~Wiseman, is to be put even upon the power of God, we may put another limitation, and that upon higher authority than his: ``God cannot deny Himself.'' In that revelation which He has given to us in our instinctive and primary perceptions of sensible things, and in that other revelation which He has given to us in His Word, God, who is the Truth, cannot contradict Himself.3

2d, An attempt is made by Romanists to identify, as one and the same in principle, the dogma of transubstantiation and what are called the mysteries of revelation. ``What,'' says Cardinal Wiseman, ``becomes of the Trinity? What becomes of the incarnation of our Saviour? What of His birth from a virgin? And, in short, what of every mystery of the Christian religion?'' It will be time enough to answer such questions as these when it is proved that such mysteries contradict our rational nature, in the same manner as the dogma of transubstantiation contradicts our perceptive nature. Such mysteries as those referred to are above our reason, but not against it. They are beyond the powers of our rational nature fully to understand, but not contradictory to our rational nature so as to be inconsistent with it. The argument in defence of transubstantiation, drawn from such a source, is but one example out of many that could be quoted, of the common tactics of Romish controversialists, who are but too often prepared to hand over to the unbeliever the most sacred truths which the Scripture has recorded, rather than not make out a plea for their own superstitions.

\hypertarget{the-doctrine-of-the-real-presence-and-the-priestly-theory}{%
\section{The Doctrine Of The ``Real Presence'' And The Priestly Theory}\label{the-doctrine-of-the-real-presence-and-the-priestly-theory}}

With the dogma of transubstantiation, as held by the Church of Rome, stands very closely connected the question as to the manner in which Christ is present in the ordinance of the Supper. The doctrine of the ``real presence'' of Christ in the Sacrament has, more almost than any other in theology, been made the subject of prolonged and bitter controversy. By the Church of Rome, as we have seen, the real presence of Christ is explained to be the true and actual existence of the body and blood, the soul and Divinity of the Saviour, under the sensible appearances of bread and wine; so that in the elements Christ is as much present after a bodily sort, in consequence of their transubstantiation, as He ever was present to His disciples of old in the days of His flesh. By the Lutheran Church, the real presence of Christ in the ordinance is maintained, not upon the principle of such a change in the substance of the elements into Christ's body and blood as contradicts the testimony of our senses, but, rather upon the supposition that the bread and wine remaining the same, the real body and blood of Christ are nevertheless united to them in some mysterious manner, so as to be actually present with them, and actually received along with them, when they are partaken of by the communicant. By our own Church, as well as by many other Protestant communions, the real presence of Christ in the Sacrament is asserted on the ground that He is not in a bodily manner present in the substance of the elements, nor yet in a bodily manner mysteriously present with the elements, but only spiritually present to the faith of him who receives the ordinance in faith.

The influence of the fierce and frequent controversies waged in connection with the nature and efficacy of the Lord's Supper shortly after the date of the Reformation, and the disposition on the part of Luther, and the Churches affected by his influence, to depart as little as possible from the established phraseology of the ancient Church on the subject of the Sacrament, served to introduce, or to continue in theological discussions, a language somewhat exaggerated, and occasionally almost unintelligible, in regard to this question. Such, undoubtedly, was the phrase ``consubstantiation,'' used by some of the Lutherans to express the mysterious corporeal presence of Christ, not in, but with, or under, or somehow in connection with the elements; and such also was the phrase ``impanation,'' employed by others to elucidate, or rather to obscure, the doctrine of the manner in which Christ's bodily presence is connected with the sacramental bread. And I cannot help thinking that, under the power of very much the same influences, the term ``real presence'' has not unfrequently been employed and explained, even by orthodox divines, in such a way as to give a somewhat exaggerated and mysterious aspect to the connection subsisting between Christ and the Sacrament. That phrase has occasionally been employed in association with such language as to leave the impression that Christ was present in the Supper, not spiritually to the faith of the believer, and not corporeally to the senses of the communicant, but in some indefinite manner between the two, and after a sort mysterious and peculiar to the Sacrament of the Supper. Such language seems to have no warrant in the Word of God.

The Scriptures give us no ground to assert that Christ is present in the Sacrament of the Supper in a manner different from that in which He is present in the Sacrament of Baptism. I do not speak at present of the extent of the blessing or of the grace which He may impart in the one or the other Sacrament by His presence; I speak only of the manner of His presence. There is nothing, I think, in Scripture to warrant us in affirming that the manner of Christ's presence in the Supper is in itself unique or peculiar, or indeed in any respect different from the manner of His presence in Baptism, or any other of His own ordinances. In all of these He is present, after a spiritual manner, to the faith of the participator in the ordinance, and in no other way. The blessings which that presence may impart may be different in different ordinances, and may be more or less in one than in another. But there is nothing in the Word of God which would lead us to say that the real presence of Christ in any of His ordinances, whether sacramental or not, is anything else than Christ present, through his Spirit and power, to the faith of the believer. Such promises as these---``Lo, I am with you alway, even unto the end of the world;'' ``Where two or three are met together in my name, there am I in the midst of you;'' ``Behold, I stand at the door and knock: if any man hear my voice, and open the door, I will come in to him, and will sup with him, and he with me;'' and such like---plainly give us ground to affirm that Christ, through His Spirit, is present in His ordinances to the faith of the believer, imparting spiritual blessing and grace. But there is nothing that would lead us to make a difference or distinction between the presence of Christ in the Supper and the presence of Christ in His other ordinances, in so far as the manner of that presence is concerned. The efficacy of the Saviour's presence may be different in the way of imparting more or less of saving grace, according to the nature of the ordinance, and the degree of the believer's faith. But the manner of that presence is the same, being realized through the Spirit of Christ, and to the faith of the believer. The Sacramentarian controversy has tended in no small measure to introduce into the language of theology, in connection with the ``real presence,'' an ambiguity of thought and statement, not confined to Romanist, or even semi-Romanist divines.

But, passing from that part of the Popish theory of the Supper which refers to the alleged change produced on the elements by transubstantiation, and to the manner of Christ's presence in the ordinance, I go on to consider the other part of the Popish theory of the Supper which refers to the office of the ministering priest in the Sacrament, or his power to offer the body and blood of Christ, actually present, as a true sacrifice for sin. The first grand error in the Popish doctrine of the Lord's Supper is the monstrous figment of the transubstantiation of the elements; the second, intimately connected with the first, and perhaps yet more extensive and mischievous as an error in its practical bearings, is the doctrine of the power of the Church, in the ordinance of bread and wine, to offer a true and efficacious propitiation to God, both for the living and the dead. The sacrifice of the mass is founded upon, and very closely connected with, the dogma of transubstantiation,---in some sort following as an inference from the assumption that the priest stands in Christ's stead at the Communion Table, and, by a supernatural power not inferior to Christ's, changes, by the utterance of the words of institution, the elements of bread and wine into the actual body and blood, soul and Divinity, which were once the sacrifice offered up for this world upon the Cross. In the performance of this supernatural and mysterious office, which, according to its own theory, it is given to the Church of Rome to discharge, we see both the priest and the sacrifice,---the priest, acting as mediator between God and the people, offering a true satisfaction to God for sin, and promising remission and reconciliation; and the sacrifice presented to God, real and efficacious, because in fact the very same sacrifice, in its substance, of the flesh and blood of Christ, as He Himself once made and presented, and not less availing in its mighty virtue to propitiate God, and procure salvation for the sinner. A real office of priesthood, and a real offering of sacrifice, are the two features that characterize this second portion of the Popish theory of the Sacraments. Both are asserted, and both are essential in the sacrifice of the mass, which has been grafted on the dogma of transubstantiation, and both form integral parts of that monstrous system of sacerdotal usurpation by which the Church of Rome seeks to build up her spiritual tyranny. The position, then, laid down by the Church of Rome in connection with the subject of the mass, may be conveniently discussed under these two heads: first, the claim which she makes to possess and exercise the office of a true priesthood; and second, the power that she arrogates to make and offer a true sacrifice to God. Reserving the second of these points for future consideration, we shall now proceed to deal with the claim put forth by the Church of Rome to hold and exercise the office of a real priesthood.

This claim runs through the doctrine and practice of the Popish Church in all its departments, and is not restricted to the case of its views in connection with the Supper. The priestly office and sacerdotal pretensions are recognised in almost every branch of its administration as a Church, and, indeed, are fundamental to the system. But the priesthood which it pretends to exercise towards God and on behalf of man is perhaps developed most prominently and conspicuously in connection with its doctrine of the Lord's Supper. The question is one that lies at the very root of the difference between the Popish and Protestant systems, and on that account is of more than ordinary interest and importance.

The doctrine of a real priesthood residing in the Christian ministry, more especially in connection with its chief function of offering the sacrifice of the mass, is thus stated by the Council of Trent: ``Sacrifice and priesthood are so joined together by the ordinance of God that they existed under every dispensation. Since, therefore, under the New Testament the Catholic Church has received the holy visible sacrifice of the Eucharist by the institution of the Lord, it is necessary also to confess that there is in it a new, visible, and outward priesthood into which the old has been transferred. Now the sacred writings show, and the tradition of the Catholic Church has always taught, that this was instituted by the same Lord our Saviour, and that a power was given to the Apostles, and their successors in the priesthood, of consecrating, offering, and administering His body and blood, and also of remitting and retaining sins.'' ``If any shall say that by these words, `Do this in remembrance of me,' Christ did not appoint the Apostles to be priests, or did not ordain that they and other priests should offer His body and blood, let him be accursed.'' ``If any shall say that the sacrifice of the mass is only one of praise and thanksgiving, or a bare commemoration of the sacrifice accomplished upon the Cross, but not propitiatory; or that it only profits him who receives it, and ought not to be offered for the living and dead, for sins, pains, satisfactions, and other necessities,---let him be accursed.''

Amid the other errors contained in these statements by the Council of Trent, what we have chiefly to do with at present is the claim which is put forth on behalf of the Church of Rome and her ministers to hold and exercise the office of priesthood in the same sense as, ceremonially, the priests of a former dispensation did so; with power now, not ceremonially, but really, to act as priests in the absence of Christ in heaven, and truly to offer sacrifice to God for sin. The question in regard to such a claim is this: Have we any warrant to believe that a visible and external priesthood has been established in the New Testament Church, with powers to act as mediators between God and man, and offer the propitiatory sacrifice for the living and the dead; or has the office of priesthood which existed under a former economy no longer an existence now in the Gospel Church, there being none on earth authorized or qualified to undertake it,---the one Priesthood, in the end of the world for sin, having completed its work on earth, and the Priest who held the office having returned to heaven to continue it there? This is a vital and fundamental question, not only in order to enable us to form an estimate of the real character of the system of Romanists, but also because it enters so essentially into the principles held by High Churchmen of other denominations.

\hypertarget{the-existence-of-a-priesthood-as-a-standing-ordinance-in-the-christian-church-is-inconsistent-with-the-fact-that-such-an-office-was-abrogated-with-the-jewish-economy-and-necessarily-came-to-an-end-when-that-dispensation-gave-place-to-the-gospel-economy.}{%
\subsection{The existence of a priesthood as a standing ordinance in the Christian Church is inconsistent with the fact that such an office was abrogated with the Jewish economy, and necessarily came to an end when that dispensation gave place to the Gospel economy.}\label{the-existence-of-a-priesthood-as-a-standing-ordinance-in-the-christian-church-is-inconsistent-with-the-fact-that-such-an-office-was-abrogated-with-the-jewish-economy-and-necessarily-came-to-an-end-when-that-dispensation-gave-place-to-the-gospel-economy.}}

An earthly priesthood was an ordinance appointed for a special purpose and a special time; and the purpose having been served, and the time past, it is necessarily at an end. The priestly office, and the institution of sacrifice with which it stands inseparably connected, formed part of that instrumentality by which, for thousands of years, God prepared this world for the coming and the death of His own Son as its Saviour. First of all, it was the father of the family who was ordained the priest to offer the sacrifice for the rest, and to approach unto God on behalf of his household; the members of which drew near to God, and worshipped, and were accepted only through him. Such seems to have been the practice in patriarchal times, and apparently not without the appointment, or at least the sanction, of God. The father of the family, as well as the divinely appointed sacrifice he offered, thus in a general and distant way represented Christ as the medium whereby sinners might approach to God in worship. But the patriarchal institute was too general and vague a type of the One Mediator through whom alone, when fully revealed, men were to find access to God. Accordingly it was done away with, and another institute was ordained in its place, with priests specially set apart to the office of mediators between God and the people, and with more special authority given, and more distinct provision made for them to be the media through whom the rest were to present their worship and sacrifices, and themselves to make their approach to God and find acceptance. Under the Mosaic ritual, it was no longer lawful for the sinner himself directly to approach to God with his own offering of worship or sacrifice; it was no longer lawful for the sinner even to draw near with his sacrifice unto God through the head of the family, as under the patriarchal institute. The avenue of approach to God was, step by step, narrowed and restricted. First, the father of the family was marked out and selected as the recognised priest and mediator for the rest. Next, a further limitation took place, and the priest of Aaron's line was specially appointed to stand in the stead of the whole families of the nation in their approach to God; and strict provision was made---and guarded by the most solemn penalties---that no man should venture to present the sacrifice himself, or to worship except through the media of this one commissioned priesthood. The thousands of Israel were restricted in their legal worship to the one avenue, and forbidden to draw near to the Holy One of Israel except through the one mediation of the earthly priest of Aaron's lineage.

And why was it that this earthly priesthood was thus marked off from all the rest, and the other worshippers made dependent on the one appointed priest of Aaron's house? And why were men forbidden to approach to God directly and immediately themselves, or even indirectly through any other but this one mediator? The answer is obvious. The priesthood was so restricted, and so fenced about with solemn limitations, in order that it might be a type of Christ, ``the one Mediator between God and man.'' From age to age, and from step to step, the worshippers of God under the old economies were more and more shut up to the idea and the practice of approaching the Most High God only through the channel of one Priesthood and the person of one High Priest. The typical priests and priesthoods of former dispensations led men's hearts and habits to fix upon the one Mediator through whom alone we now draw near to God. They taught the worshippers to anticipate and to hope in that one Man, who is now the Priest, not of one family, as in patriarchal times, nor of one nation, as in Jewish times, but the Priest through whom all the families and all the nations of the world draw night to God. The earthly priesthoods of the former days of the Church all converged upon and pointed to and centred in Christ. With Christ, therefore, those priesthoods came to an end. The type was merged in the Antitype, and then was done away. The priests of other days, together with the sacrifices which they offered, have served the object designed by them, and are abolished. They can, from the very nature of their office, have no use, and no meaning, and no place in a Church to which another and a higher priesthood has been given, and when the sign has given place to the thing that was signified. The office of the priesthood on earth ceased with the former dispensation; and not only is there no re-appointment under the Gospel of such an order of men in the Church, but they would, from the very place and office that they occupied, be inconsistent with the Gospel economy. They formed part and parcel of a typical system which has been abolished.

\hypertarget{the-existence-of-a-priesthood-as-a-standing-ordinance-in-the-christian-church-is-inconsistent-with-the-privileges-of-believers-under-the-gospel.}{%
\subsection{The existence of a priesthood as a standing ordinance in the Christian Church is inconsistent with the privileges of believers under the Gospel.}\label{the-existence-of-a-priesthood-as-a-standing-ordinance-in-the-christian-church-is-inconsistent-with-the-privileges-of-believers-under-the-gospel.}}

It is not unfrequently argued by the advocates of Romanist or semi-Romanist principles on this subject, that the privilege of a human priesthood and a human mediatorship is one so great and precious that it cannot be conceived to exist, as we know it did, under the earlier and far inferior dispensation, and yet to be awanting under the later and far better dispensation of the Gospel. The presence of an earthly priesthood, it is urged, must be enjoyed by the Church now, inasmuch as it cannot be supposed to be deprived of one of the highest privileges which belonged to the former and less richly endowed Church of the Old Testament.

A comparison between the superior advantages of the Gospel Church, as measured by those of the Jewish, is the very consideration which, instead of proving that a human priesthood is continued to us now, most emphatically demonstrates that it is abrogated. The presence and office of a human priesthood, enjoyed by worshippers under the law, are far surpassed by the higher and more glorious privileges enjoyed by believers under the Gospel. No doubt it was an act of grace and condescension on the part of God, to permit sinners to approach His presence through the avenue of a visible priesthood and a visible sacrifice in former times, even although that boon was granted to them under solemn and jealous restrictions; and it was a great and precious privilege for the worshipper to be allowed to draw near to the mercy-seat through means of a human mediator, and by the intervention of a material offering. But the privilege of Christians in the New Testament Church is better and more glorious still. Through Christ a new and living way has been opened up for all to draw nigh to God, not indirectly through a human mediator, but directly, each man for himself. The whole brotherhood of believers are no longer dependent upon one of themselves for the liberty or opportunity of access to the common Father; and without distinction of special office, it is the freedom purchased for all, without earthly priest or earthly intercessor interposed, to go with boldness into the very holiest. The presence of an earthly and external priesthood is no evidence of superior privilege, but the reverse. It is the mark of an imperfect and carnal dispensation.

That it was necessary for the worshipper to employ the intervention of another than himself in order that he might approach to his Creator,---that a sinner should be dependent on another sinner for pardon or access to heaven,---that he should not dare to engage his heart to draw near to God except through the medium of a human priesthood,---were strong arguments to prove the essential imperfection of that dispensation which witnessed such things, and constituted a yoke of bondage which it was hard to bear. And what it was when the sons of Aaron by God's own appointment were the human priests and mediators, that it is now in the case of those Churches who bind upon their own necks the institute of a human priesthood, and then boast of it as their exclusive distinction and privilege. It is a spiritual yoke that is too heavy to bear; it is a retrogression from the freedom wherewith under the Gospel Christ has made His people free; it is a badge of the voluntary thraldom and debasement of a Church that has itself gone into bondage to men, instead of maintaining the liberty of Christ the Lord. The restriction of approaching God only through the earthly priest in the local temple at Jerusalem, and by the blood of bulls and goats,---the prohibition forbidding the sinner to draw near to the mercy-seat directly himself, or through any other medium,---those were evidences of essential imperfection in the Church state of the worshippers under a former economy. And the human priesthood of the Church of Rome,---the material sacrifice made and offered for the worshippers,---the priest standing between the sinner and God, and barring or opening the way of approach,---the mediator acting as the medium of communication between the Most High and His creatures, and retaining or remitting their sin,---these, too, are restrictions, and, because human and unauthorized, daring and impious restrictions, upon the freeness of God's grace and the liberties of His redeemed people.

It is a fact of much significance, and indeed of decisive force in this argument, that throughout the whole of the New Testament Scriptures there is no instance in which either the name of priest, or the functions belonging to the office of priesthood, are ascribed to the ministers of the Christian Church; that the only examples of the use of the term are those in which it is given, not to the minister, but to the people; and that the ascription of the privileges of the office is uniformly made to the members at large. On the one hand, the term ἱερευς, or ``priest,'' is never in any single instance in the New Testament applied to a minister of the Christian Church, although always made use of to designate the priest of the Aaronic dispensation. The usual name given to the minister of the New Testament Church is πρεσβυτερος,---the change of designation marking very decisively the change in the nature of the office. On the other hand, on the only occasions on which the word ἱερευς is used in the New Testament in reference to any except a Jewish priest, it is given to the members of the Christian Church at large, and not to the ministers of that Church. In the Book of Revelation, believers are spoken of as ``kings and priests to God;'' and in the first Epistle of Peter they are described as a ``royal priesthood.'' The name formerly appropriated to the sons of Aaron, selected and anointed from among the rest of the congregation to be priests to God, is not inherited by the ministers of the Christian Church in the same exclusive manner, but, on the contrary, is applied in an enlarged and extended sense to the whole body of believers. More than this: the privilege enjoyed by the priests of old, of alone of all the worshipping assembly drawing near to God without the intervention of any other, is a privilege uniformly represented in the New Testament as not peculiar to the ministers of the Church, but extended now to all its members, and common to all believers. The office peculiar to the minister of the Christian Church is described at large in the New Testament Scriptures, and is a ``ministry'' or ``service'' unto others (διακονια, λειτουργια), not a mediatorship on behalf of others. It is spoken of as an office of ``ministering,'' ``preaching,'' ``exhorting,'' ``ruling,'' amid the flock of Christ, not an office of sacrificing, and making reconciliation, and approaching to God as the mediator on behalf of the rest, and becoming the avenue for the access of their persons or worship to the Divine presence. On the contrary, this privilege of approaching directly to God without the intervention of any substitute or proxy on earth, is a privilege which is expressly attributed to all believers as their personal right: so that, if in any sense there are priests now on earth, those priests are the believing people of God at large; and if in any sense there are priestly sacrifices now offered up, they are the spiritual sacrifices of the prayer and praise of Christians, without distinction of office or place in the Church. The sacerdotal theory on which the Church system of Rome is built, and the priestly office which is so conspicuously developed in her practice as regards the Lord's Supper, are utterly repugnant to the spirit of the New Testament Church, and to the privileges which it has secured to believers. The privilege of a human priesthood, which existed under the law, is abolished under the Gospel; or rather, in its spirit and substance, the privilege is enlarged and extended to all believers under the New Testament Church. It was the peculiar and distinctive prerogative of the priests under the law, that they alone of all the worshippers drew near to God without a human mediator. That prerogative is common to all the royal priesthood of believers under the Gospel.

\hypertarget{the-existence-of-an-earthly-priesthood-as-a-standing-ordinance-of-the-christian-church-is-inconsistent-with-the-one-office-of-christ-as-the-priest-and-mediator-of-his-people.}{%
\subsection{The existence of an earthly priesthood as a standing ordinance of the Christian Church is inconsistent with the one office of Christ as the Priest and Mediator of His people.}\label{the-existence-of-an-earthly-priesthood-as-a-standing-ordinance-of-the-christian-church-is-inconsistent-with-the-one-office-of-christ-as-the-priest-and-mediator-of-his-people.}}

Earthly priest the New Testament Church has none. The very name is blotted out from the inspired history of the Church under the Gospel in its application to any office-bearer within its pale; and it is found, in so far as it can now be found on earth, only in connection with that spiritual and universal priesthood which belongs alike to all true believers, who have equally the privilege of free approach to God, equally the anointing which makes them His people, and equally the consecration that sets them apart for His service. In any other sense than this, there is no priest in the Christian Church on earth. The material sacrifice made by men has ceased, the incense kindled by men no longer burns, the atonement presented by men is no more offered up. The Gospel is a religion without a priest on earth, without a sacrifice, and without an altar. And yet there is a priesthood that belongs to the Christian Church still; and there is a Priest who yet discharges that office on behalf of His people. ``We have a great High Priest that hath passed into the heavens for us,''---not a mortal and dying man, but one ``of whom it is witnessed that He liveth for ever,''---not a priest who offers, as did the sons of Aaron of old, the typical sacrifices of blood, or, as the ministers of Rome do now, the pretended sacrifices of an unbloody offering of bread and wine,---but one who, once for all, offered up a Divine yet human sacrifice for men,---not an intercessor, who, like the high priest under the law, entered into God's presence with the blood of bulls and goats, nor yet like the priest of the Papacy with a consecrated wafer,---but an Intercessor, who, with His own precious and more than mortal blood, has passed into the presence of God,---an Intercessor, the Son of God, presenting the offering of Himself without spot or blemish, and pleading for us on the ground of His meritorious sacrifice. And this office which the Son of God now discharges in heaven for His Church passes not from Him to any other (ἀπαραβατον ἐχει την ἱερωσυνην.) His is an unchangeable and undying Priesthood; and He ever liveth to make intercession for His people. The office which He sustains and discharges in heaven is His own incommunicable office, which none save Himself has either the right or the power to discharge. The one Priest that has made the sacrifice and offered it to God for the sins of many,---there was none that could share with Him in that mighty and mysterious work. The one Priest to stand between God and a sinful world,---there was none but the Son that could undertake so to approach unto the Most High. The one Priest to intercede with an offended God for the guilty,---there was none but the equal of the Father that could so plead. The one Priest to dispense unto men throughout all ages the blessings of redemption and grace,---there is none equal to the task but He ``in whom dwelleth all the fulness of the Godhead bodily.'' Alone in His office as in His nature, unapproachable in His work as in His greatness, ``He abideth a Priest for ever,''---the ever-present and ever-living Mediator, who has no fellow to share in His priestly functions, and whose glory as Mediator He will not give unto another.

And what shall we say of those Church systems, Romanist and semi-Romanist, that give to mortal men that office of Priest which none can bear but the Son of God, and constitute sinners mediators on earth between their fellow-sinners and the Almighty? Such an encroachment upon His incommunicable office touches very nearly the honour of Christ. The assumption by men of His personal and inalienable prerogatives, inseparable from Himself as Mediator, is a dishonour done to Him in that very character in which He stands forth supreme and alone before the eyes of the universe. The very title of Mediator belongs in the Christian Church to none but One, and He the only-begotten Son of the Father. Our lips are now forbidden to name another Priest but Jesus. Even in the Old Testament Church, the name and the office of the Priest had something in them of awful and mysterious import, typical as they were of the fulness of the Gospel day, and of the greatness of the Gospel Mediator, and fenced about, as we know them to have been, with the solemn and irrevocable sentence of death upon those who should unwarrantably assume or encroach upon them. And still more awful are that name and office of Priest, now that in these latter days they have been sustained by the Son of God Incarnate, and mysteriously sanctified by the shedding of that more than mortal blood which was poured out on Calvary, and which He still day by day presents in heaven, as He continually pleads with the Father there. To stand between God and man, as Christ once stood amid the darkness of Calvary, was a work which none but He could do. To stand between God and man, as Christ now stands, a Priest in heaven no less than on earth, is a work which none but He can accomplish. To bear the burden of such an office now is as little competent to mortal man as it was to bear the burden of it in the Garden, or at the Cross. The name of Priest between God and man is Christ's inalienable and incommunicable name,---whether He bears the anger of an offended Judge, or pleads with the compassion of a reconciled Father,---whether He makes, as He once did, atonement by sacrifice, or makes, as He now does, intercession by prayer. It is the sin above others of the Church of Rome, that it has assumed to itself that name of Priest, which none in heaven or in earth is worthy to bear but the Son of God, and that its ministers pretend to stand between the creature and the Creator in the exercise of His priestly office among men.

\hypertarget{the-sacrifice-of-the-mass-and-other-forms-of-the-sacrificial-theory}{%
\section{The Sacrifice Of The Mass, And Other Forms Of The Sacrificial Theory}\label{the-sacrifice-of-the-mass-and-other-forms-of-the-sacrificial-theory}}

The claim to the possession of a real priesthood, and to the power of making and presenting to God a real propitiatory sacrifice, is fundamental to the theory of the Church of Rome, and is one of the great pillars on which its spiritual strength leans. The right to stand between God and man in the character of mediator, to exercise the priest's office in place of Christ on the earth, to negotiate as man's intercessor with God, and to arrange the terms of his acceptance or condemnation, to make and offer the sacrifice which alone can avail unto justification of life, to retain or remit sin, to give or withhold saving grace,---in short, the claim to the sacerdotal office lies at the very foundation of the Popish system. This one principle of a priestly power existing in her ministry, accompanying all their administrations, and sanctifying all their acts, runs through the whole details of the Church system of Rome, and is the grand secret of very much of its success. We see it fully and conspicuously developed in connection with the Romish doctrine of the Supper, and as the foundation of the sacrifice of the mass. But it is not confined to that one department of the Popish Church system. The sacerdotal principle pervades it, more or less, throughout its entire range; and the Church of Rome has thus added to its many sins the one emphatic sin of usurping the place of Him who has an unchangeable priesthood in heaven and on earth, and of seizing out of His hands the powers that He wields as ``Priest for ever.'' But great and awful though the sin be of arrogating the place and prerogatives of the one High Priest of His people, it is yet a sin which pays its price to the Church that commits it, in the spiritual prestige that it confers, and the spiritual authority that it brings along with it. A sense of the need of some mediator between the sinner and an offended God, a feeling of the absolute necessity of a priest and intercessor for a fallen creature, to negotiate the terms of his pardon and acceptance, can hardly ever be rooted out from the guilty conscience. And the Church of Rome, when it ventures to arrogate to itself on earth that very office which guilty nature needs, and succeeds in its perilous claim to be regarded as the only priest and intercessor between sinners and God, establishes for itself a spiritual dominion over the souls of its victims, greater and more absolute than any other dominion in this world. And hence the tenacity with which the Romish Church clings to the claim of a priestly or sacerdotal office, inseparably connected as it is with some of the most monstrous and incredible pretensions, with the dogma of transubstantiation, with the claim to forgive sin, which none but God can do, with the pretence of making and presenting a Divine and propitiatory sacrifice to the Almighty.

In spite of the explicit abrogation of the office with the abrogation of the Old Testament dispensation; in spite of the palpable inconsistency of the office with the spirit of the Gospel, and the privileges of believers; and, worse still, in spite of the inconsistency of the office with the sole priesthood of Christ, the Church of Rome ordains each one of her ministers to be a priest, and invests him with the power and authority of an earthly priesthood. It needs must be that a priest have a sacrifice to present unto God. ``This man must of necessity have somewhat to offer.'' And having ordained, as she alleges, a real priest, the Church of Rome proceeds to put into his hands a real sacrifice, and gives him warrant to offer it to God for the sins of the living and the dead.

The doctrine of the Church of Rome on this vital point is laid down in such a manner in her authorized formularies that it is impossible to explain it away. The Council of Trent has defined it in such terms, that the attempts made by more modern Romanists to soften down the atrocious dogma of the real offering-up of the sacrifice of the Lord, body and blood, soul and Divinity, in the Sacrament by the priest, are in vain. Speaking of ``the institution of the most holy sacrifice of the mass,'' the Council declares that it is ``a visible sacrifice, as the nature of man requires, by which that bloody one, once to be accomplished on the Cross, might be represented, and the memory of it remain even unto the end of the world.'' And with this statement, expressive of the representative or commemorative character of the ordinance, the apologists of the Church of Rome, whose desire is to conceal the real doctrine held by her on this subject, very often terminate their quotation, as if the Council of Trent held it to be no more than a symbolical sacrifice in memory of Christ's. But that this is not the case, the words of the Council's definition leave us no room to doubt. It proceeds: ``For after the celebration of the old passover, which the multitude of the children of Israel sacrificed in memory of their departure from Egypt, Christ instituted a new passover, even Himself, to be sacrificed by the Church through the priests under visible signs (Seipsum ab Ecclesiâ per sacerdotes sub signis visibilibus immolandum), in memory of His departure out of this world unto the Father, when by the shedding of His blood He redeemed us and snatched us from the power of darkness, and translated us into His kingdom.'' ``And since in this Divine sacrifice, which is performed in the mass, that same Christ is contained and immolated in an unbloody manner, who on the altar of the Cross once offered Himself with blood, the holy Synod teaches that that sacrifice is, and becomes of itself, truly propitiatory; so that if with a true heart and right faith, with fear and reverence, we approach to God, contrite and penitent, we may obtain mercy and find grace to help in time of need. Wherefore the Lord, being appeased by the offering of this, and granting grace and the gift of repentance, remits crimes and sins, even great ones. For it is one and the same victim,---He who then offered Himself on the Cross being the same Person who now offers through the ministry of the priests, the only difference being in the manner of offering (Una enim eademque est hostia, idem nunc offerens sacerdotum ministerio, qui Seipsum tunc in cruce obtulit, sola offerendi ratione diversa).'' And, once more: ``If any shall say that the sacrifice of the mass is only one of praise and thanksgiving, or a bare commemoration of the sacrifice which was made upon the Cross, but not propitiatory; or that it only profits him who receives it, and ought not to be offered for the living and the dead, for sins, pains, satisfactions, and other necessities,---let him be accursed.''

There are two things in regard to the doctrine of the Church of Rome put beyond all dispute or cavil by these statements. First, it is Christ Himself transubstantiated into the elements, and corporeally present in the Sacrament, that is offered up by the priest as a real sacrifice. It is utterly impossible for Romanists to escape from this dogma so long as the language of Trent remains uncancelled. No attempt can succeed to give it a mystical or symbolical meaning, and soften down the authoritative assertion of the Council, that in the Supper there is a real sacrifice of Christ Himself by the priest. Romish controversialists may indeed adopt different modes of explaining how the sacrifice of the mass stands related to the sacrifice of the Cross. Some of them, like Harding the Jesuit, in his reply to Bishop Jewel, may plainly and unhesitatingly assert ``that Christ offered and sacrificed His body and blood twice,---first in that holy Supper, unbloodily, when He took bread in His hands and brake it, and afterwards on the Cross with shedding of His blood.'' Others of them, like Möhler, in his Symbolism, with a view to make the doctrine less palpably inconsistent with Scripture, may assert another form of it, and maintain that there are not two sacrifices, but one, and that the sacrifice of the Supper constitutes a part of that sacrifice which Christ offered on the Cross; or, to use Möhler's own language, ``Christ's ministry and sufferings, as well as His perpetual condescension to our infirmity in the Eucharist, constitute one great sacrificial act, one mighty action undertaken out of love for us, and expiatory of our sins, consisting, indeed, of various individual parts, yet so that none by itself is, strictly speaking, the sacrifice.'' ``The will of Christ to manifest His gracious condescension to us in the Eucharist, forms no less an integral part of His great work than all besides, and in a way so necessary, indeed, that whilst we here find the whole scheme of redemption reflected, without it the other parts would not have sufficed for our complete atonement.'' But however Romanists may choose to explain it,---whether as a repetition of the sacrifice of the Cross, or a continuation of it,---the Supper is unquestionably, according to the doctrine of the Church of Rome, a real sacrifice, made up of Christ's body and blood. And second, this real sacrifice is truly propitiatory in its nature, having virtue in it to satisfy Divine justice, and to constitute a proper atonement for sin. These two doctrinal positions are clearly and undeniably laid down by the Council of Trent, and in such a manner that Romanists cannot evade them. And it is certainly one cause of thankfulness, and no small one, that the Council of Trent was overruled by Divine Providence to put this and other of the monstrous tenets of Romanism into such a dogmatic and articulate form, that it is now utterly impossible for the Church of Rome to deny or escape from them.

What, then, are we to say to the real sacrifice asserted by the Church of Rome, a true propitiation to God for sin, repeated day after day by countless priests who have authority and power to make and offer it?

\hypertarget{the-doctrine-of-the-church-of-rome-is-in-direct-contradiction-to-the-doctrine-of-scripture-which-declares-that-there-is-one-priest-and-no-more-than-one-under-the-gospel.}{%
\subsection{The doctrine of the Church of Rome is in direct contradiction to the doctrine of Scripture, which declares that there is one Priest, and no more than one under the Gospel.}\label{the-doctrine-of-the-church-of-rome-is-in-direct-contradiction-to-the-doctrine-of-scripture-which-declares-that-there-is-one-priest-and-no-more-than-one-under-the-gospel.}}

``Sacrifice and priesthood,'' say the Fathers of the Council of Trent, ``are so joined together by the ordinance of God, that they existed under every dispensation.'' There can be no doubt that the statement is correct in this sense, that wherever there is a sacrifice, there must be a priest to offer it, and wherever there is a priest, he must of necessity have a sacrifice to offer.2 And hence, as part of the sacrificial theory of the Supper and essential to it, the ordination by which the Church of Rome sets apart persons for the work of the ministry includes, as its main and characteristic feature, a commission not to preach the Gospel and to dispense its ordinances, but to make and offer sacrifices to God for the souls of men. Hers is mainly and distinctively an order of priests, and not an order of ministers,---a succession from age to age of sacrificers and intercessors, and not of preachers. And thus her system is distinctively opposed to the system of Scripture, which points to one Priest, and forbids our lips to name a second in the Gospel Church. The argument of the last section might be sufficient, without further illustration, to establish this. But the point is so vital, and it is brought out with such power and effect by the Apostle Paul, that I cannot help adverting to his statements on this subject.

The grand design of that magnificent exposition of the doctrine of Christ's office and nature and work in the Epistle to the Hebrews, is to prove that, far above and beyond the mediators and priests under the law, Christ was the one Son and the one Priest of God, in a way and manner altogether exclusive and peculiar, and such as to contrast Him with all others who ever, in any secondary sense, bore these names. In regard to the priesthood more especially, there were under former dispensations two orders of priests, with one of which the apostle compares our Lord, with the other of which the apostle contrasts Him; and both the comparison and the contrast serve to bring out more distinctly the singular and exclusive character that He bears as the Priest of God, who has neither partner nor successor in the office. There was, according to the apostle, a priesthood after the order of Melchisedec, and there was a priesthood after the order of Aaron. With the priesthood after the order of Melchisedec our Lord is compared. There was room in that order for but one Priest, and no more than one; and for this reason, as stated by the apostle, ``He abideth a Priest continually.'' In the office that he held He had no predecessor, and He had no successor. Melchisedec stood alone in the typical order that bears his name; and the more surely and distinctly to mark out this singularity of his position, we are told, with respect to his office, that he was ``fatherless, motherless, ungenealogied, having neither beginning of days nor end of life'' (ἀπατωρ, ἀμητωρ, ἀγενεαλογητος, μητε ἀρχην ἡμερων μητε ζωης τελος ἐχων). And such as the type was, so is the Antitype. The Lord Jesus Christ was ``made a Priest after the order of Melchisedec;'' and, like that of His type, His office is singular and exclusive; He knows neither predecessor nor successor in it; having not only in His Divine nature, but in His mediatorial character, ``neither beginning of days nor end of life.'' None went before, and none shall come after this Priest; or, as the apostle expresses it, His office is one ``that passeth not from Him to any other.''2 The comparison instituted between our Lord's priesthood and that of Melchisedec demonstrates that He is the one Priest, with none to go before or succeed Him in that character.

But again, with the priesthood of Aaron that of our Lord is contrasted by the apostle; and the contrast serves to bring out in like manner the very same grand doctrine. In that priesthood there were not one, but many priests, following each other in rapid succession. The mortal and dying men who inherited the blood and the office of Aaron ``were not,'' as the apostle tells us, ``suffered to continue by reason of death.'' One after another passed away in swift succession, so that in the not lengthened period of the Aaronic Church there were truly ``many priests,'' following each other rapidly in office, as ever and anon death removed them from beside the altar where they sacrificed and interceded. With them our Lord is contrasted, and not compared in this respect. ``This man, because He continueth ever, hath an unchangeable priesthood.'' ``He is consecrated for evermore.'' He is endued with ``the power of an endless life,'' and ``ever liveth to make intercession for His people.'' Compared with the order of Melchisedec, and contrasted with the order of Aaron, our Lord is emphatically marked out as the one Priest of God, who can have none to follow, even as He had none to go before Him in His office. And the many priests, anointed day by day continually, and succeeding each other in rapid succession in the Church of Rome, are most decisively declared to be inconsistent with His one glorious priesthood.

\hypertarget{the-popish-theory-of-the-lords-supper-is-in-direct-opposition-to-the-doctrine-of-scripture-which-declares-that-there-is-one-sacrifice-and-no-more-than-one-under-the-gospel.}{%
\subsection{The Popish theory of the Lord's Supper is in direct opposition to the doctrine of Scripture, which declares that there is one sacrifice, and no more than one, under the Gospel.}\label{the-popish-theory-of-the-lords-supper-is-in-direct-opposition-to-the-doctrine-of-scripture-which-declares-that-there-is-one-sacrifice-and-no-more-than-one-under-the-gospel.}}

This argument is likewise brought out with commanding force and effect---as if by way of anticipation of the very error of the Papacy---in Paul's Epistle to the Hebrews. He exhibits the contrast between the many priests under the law and the one Priest of God under the Gospel, immortal, and living ever to discharge that office of priesthood in which He had no predecessor and can have no follower, and in which, like Melchisedec, He stood alone. But in close relation with this, he exhibits the contrast also between the many sacrifices under the law with their ceaseless repetition, and the one sacrifice of the Lord Jesus Christ, which never was, and never could be, repeated. The argument by which the apostle demonstrates the unspeakable superiority of the sacrifice of Christ over the sacrifices offered by the sons of Aaron, is a brief and decisive one. The very fact of the repetition of the one, and the non-repetition of the other, was the conclusive evidence of that superiority. The sacrifices under the law were repeated day by day continually; the priest had never done with offering, and the altar never ceased to be wet with the blood of the victims. What was done to-day had to be repeated to-morrow; and the sacrifice was never so completely made and finished but that it had to be repeated afresh, and renewed times without number. And why? The reason was obvious. They were essentially imperfect. They could never so accomplish the great object of atoning for sin but that their renewal was necessary; and what was done on one day had to be supplemented by what was to be done on the next. ``The law,'' says the apostle, ``having a shadow of good things to come, and not the very image of the things, can never with those sacrifices which they offered year by year continually make the comers thereunto perfect. For then would they not have ceased to be offered? because that the worshippers once purged should have had no more conscience of sins. But in those sacrifices there is a remembrance again made of sins every year.'' The fact of their ceaseless repetition was the evidence of their essential imperfection. But in contrast with this, and as an evidence of its sufficiency, the apostle urges the consideration that the sacrifice made by Christ was offered up once, and no more than once. It stood alone, as an offering made once for all, and never again to be repeated,---a sacrifice so complete in its single presentation that it admits of no repetition or renewal. Christ cannot die a second time upon the Cross, as if His first death were incomplete in its efficacy or its merits; for ``by one offering He has perfected for ever them that are sanctified'' or atoned for. Again and again the apostle renews his argument, and his assertion of the fact on which the argument is founded. ``Christ was once offered to bear the sins of many.'' ``Nor yet that He should offer Himself often as the high priest.'' ``For then must He often have suffered since the foundation of the world.'' ``He entered in once into the holy place;'' and ``we are sanctified through the offering of the body of Jesus Christ once for all.'' ``By one offering He hath perfected for ever them that are sanctified.'' The argument is decisive. The perfection of Christ's sacrifice, and the non-repetition of Christ's sacrifice, are inseparable. If that sacrifice needs to be repeated, then it cannot be perfect.

And the reasoning of the apostle is conclusive, as if by anticipation, against the many sacrifices of the Church of Rome in the Supper, whatever explanation may be adopted by its advocates to explain away the contradiction between their practice and the doctrine of Scripture. Let the sacrifice of the mass be a repetition of the sacrifice of Christ upon the Cross, as some Romanist controversialists hold it to be,---and their explanation plainly and undeniably means, that the sacrifice of the Cross needs to be repeated day by day, in order to accomplish the salvation of sinners. Or, let the sacrifice of the mass be a continuation of the sacrifice of Christ on the Cross, and a part of the same atonement, as other Romanists expound it,---and this explanation plainly and undeniably means, that the sacrifice of the Cross was not finished when Christ bowed His head and gave up the ghost. Explain the connection as you will between the sacrifice of the mass and the atonement made upon the Cross, it is utterly inconsistent with the argument of the apostle by which he proves the unapproachable perfection of Christ's work, from its being that one offering which never can be repeated or followed by another.

\hypertarget{what-is-essential-to-the-very-nature-of-a-true-propitiatory-sacrifice-is-awanting-in-the-pretended-sacrifice-of-the-mass.}{%
\subsection{What is essential to the very nature of a true propitiatory sacrifice is awanting in the pretended sacrifice of the mass.}\label{what-is-essential-to-the-very-nature-of-a-true-propitiatory-sacrifice-is-awanting-in-the-pretended-sacrifice-of-the-mass.}}

What was offered on the altar in former times could be no propitiatory sacrifice to God unless it was dedicated to Him by death. Believing sacrifice itself to be a positive institution of God, we must look for the nature and import of the observance only in His Word, and in the practice sanctioned by His appointment. And taking the case of the Old Testament sacrifices, we are warranted in saying that they were uniformly dedicated to God by death, and that ``without shedding of blood there could be no remission.'' There were, indeed, offerings under the law not connected with the shedding of blood, and not accompanied by the destruction of life; but these were not propitiatory. In every case of a propitiatory offering the victim was slain, and the atonement made through the shedding of blood. Expiation and the death of the offering---atonement and shedding of blood---were so inseparably connected, that there could be no real sacrifice of a propitiatory nature when the sacrifice was not dedicated to God by death. From the very earliest times blood was accounted a holy thing, not to be eaten or made use of for common purposes; and the very terms of the prohibition explain the reason of it: ``For the life of the flesh is in the blood, and I have given it to you upon the altar to make an atonement for your soul; for it is the blood that maketh atonement for the soul.''3 Without blood shed there could be no expiation. And here lies one difficulty of the Romish dogma of the sacrifice of the mass. It is a propitiation for the sins of the living and the dead; it is no bare commemoration of a sacrifice, but itself a sacrifice, with virtue to satisfy Divine justice and atone for sin; it is an offering of expiation offered wherever there is a priest to consecrate the ordinance and present it to God. It is a sacrifice of Christ, offered up in propitiation of His Father's righteous displeasure, and efficacious for the remission of sin. But yet we are assured by the apostle that ``Christ dieth no more; death hath no more dominion over Him. For in that He died, He died unto sin once: but in that He liveth, He liveth unto God.'' The Lord Jesus Christ, in His glorified human nature, has long since passed away from the scene of His suffering and humiliation; seated at the Father's right hand, He has rested Him from His work of sorrow and blood, and can repeat no more the agony of the Garden or of the Cross. He does bear with Him indeed in heaven, impressed for ever on His human flesh, the tokens of suffering and crucifixion; ``as a lamb that has been slain,'' He appears on high in the sight of His Father and His angels, marked with the visible evidence of sacrifice and death. But He repeats the sacrifice no more; His blood is not afresh poured out. The proofs of His once finished sacrifice which He carries about in His person are enough; and with these silent but eloquent witnesses to make good His cause, He pleads the virtue of that sacrifice, and never pleads in vain. His uninterrupted and continual advocacy, founded on the merits of His one sacrifice, all-sufficient and complete, supersedes the necessity of its repetition; He needs to die no more for the many sins of His people, which they daily renew, because He once died a death enough for them all, and now lives a life of everlasting intercession, based upon that death, for His people. Without shedding of blood, without atoning suffering, without life rendered as expiation for life, the pretended sacrifice of the mass is inconsistent with the scriptural idea of sacrifice dedicated to God by death.2

Upon such grounds as these we are warranted to say that the sacrificial theory of the Church of Rome, more fully developed in her dogma of the mass, but running throughout her whole spiritual system, is entirely opposed to the doctrine of the Word of God, which asserts, as fundamental to the Gospel, that as there is but one Priest, so there is but one sacrifice known in the New Testament Church. But there are various modifications of this sacrificial theory which, avoiding the extreme doctrine of the Papacy, are held by many semi-Romanists, and still assert that the Lord's Supper is a sacrifice. There are two of these held very commonly by High Churchmen in the English Establishment, to which I would very briefly advert.

1st, In a sense very different from the Romish, it was held by not a few of the Christian Fathers in the early centuries,---and the doctrine has been revived in more recent times in the Church of England,---that the elements of bread and wine were a true material sacrifice, not indeed propitiatory, but eucharistic; very much in the same way as the first fruits laid upon the altar by appointment of the Mosaic law, were a thank-offering to God for the overflowing of His bounties to His creatures. According to this view, the elements of bread and wine, offered to God in the Supper as a material sacrifice without blood, are the fulfilment of the prophecy of Malachi, in which he foretells, in regard to Gospel times, that ``a pure offering,'' as contradistinguished from the bloody sacrifice of the law, should then be offered in to God's name. ``From the rising of the sun to the going down of the same, Thy name shall be great among the Gentiles; and in every place incense shall be offered unto Thy name, and a pure offering.'' This sacrificial theory of the Supper is certainly free from the vital and most fundamental error of the Church of Rome, when it ascribes to the sacrifice in the ordinance a propitiatory character; but it is open to insurmountable objections.

First, a material sacrifice, in the sense of a thank-offering to God for the bounties of His providence, has not the slightest countenance in any of those passages of the New Testament which describe the nature and design of the Supper. It is hardly anything else than a conceit, gratuitously invented by those who saw that it was impossible to regard the Supper as a propitiation for sin, but who were anxious, in conformity with the unguarded language of the patristic writers on the subject, to devise some plausible excuse for applying the term ``sacrifice'' to the Supper. Second, the theory is entirely inconsistent with the first and primary characteristic of the Supper, as clearly laid down in Scripture, namely, that it is an ordinance commemorative of the propitiatory sacrifice of Christ. Third, the theory of a material sacrifice in the Supper, in the sense of a thank-offering of bread and wine for the bounties of Providence, is repugnant to the spiritual nature of the Gospel dispensation, which stands opposed to typical worship.

2d, There is another sacrificial theory of the Supper, much more common than the one now mentioned, and indeed, with various but unimportant modifications, the prevalent theory among those High Churchmen of the English Establishment who reject the extreme views of Popery, as asserted in the doctrine of the mass, but who hold that in the Supper there is a real propitiatory sacrifice, and a real sacrificing priest. According to this view, the elements of bread and wine, not transubstantiated, but remaining unchanged, become, by the words of institution and the consecration of the priest, the body and blood of Christ symbolically and mystically; in consequence of the sacramental union between the sign and the thing signified in the Sacrament, the elements are both to God and to us equivalent to and of the same value with Christ Himself; and the offering up to God of the elements, thus both representing a crucified Saviour, and not inferior in virtue or worth to the Saviour Himself, becomes a true propitiatory sacrifice made to the Almighty for sin. Upon this theory of the Supper, the office of priest in the Christian Church is similar to that of priest under the law: both offer to God real, although symbolical sacrifices, equally pointing to Christ,---there being this difference, that the Aaronic priesthood offered a sacrifice of blood in the prospect of the Saviour's sacrifice to come; while the Christian priesthood offers an unbloody sacrifice in memory of the Saviour's sacrifice now past; and also, that the sacrifices presented now in the Supper, in consequence of their sacramental union with Christ, are infinitely more precious than the sacrifices of the former economy. Such, briefly, and so far as I am able to understand it, is the prevalent doctrine among the majority of the High Church party in the Church of England at the present day, who are not yet prepared, as an extreme section of them appear to be, to accept the Tridentine definitions of the nature and efficacy of the Sacrament of the Lord's Supper. It is maintained and expounded at length in a work recently republished in the Anglo-Catholic Library, entitled, The Unbloody Sacrifice and Altar Unveiled and Supported, by Johnson.

This theory, while excluding the dogma of transubstantiation, which Romanists feel to be necessary to give consistency and foundation to their doctrine of the Supper, approaches in other essential respects very closely to that doctrine, asserting, as it does, a real sacrificing priest and a real propitiatory sacrifice in the Supper. The principles already laid down in opposition to the Popish theory of the Supper are almost all equally available against the now mentioned modification of it. It is subversive of the whole doctrine and character of the Gospel. Under the Christian dispensation there is no priest but One, and He is in heaven. It is His incommunicable name, which none in heaven or on earth may bear but Himself. There is no sacrifice or propitiation but one, and that was finished on the Cross erected upon Calvary, looking back, as it does, for thousands of years over the long array of bloody offerings, which were but the types that pointed towards it, not yet come; and looking forward, as it does, over the long array of ordinances in the Christian Church, commemorative of it, now that it is past. Neither type beforehand, nor commemoration afterhand, could share in its character as an expiatory sacrifice for sin. There is now no dedication of victims to God by death,---life given for life, and blood exchanged for blood,---in order to make a propitiation. The tragedy of the Cross cannot now be renewed, nor atoning blood be shed afresh; and yet ``without the shedding of blood there is no remission'' in Sacrament or in sacrifice. Under whatever form, or modification the sacramental theory be held, which asserts in the Supper a real sacrifice, and a true propitiation for sin, it is a dishonour done to the Lamb of God, who ``by the one offering of Himself has perfected for ever them that are sanctified,'' and who, in virtue of that one Divine offering, now ``liveth for ever to make intercession for His people.''

\bibliography{book.bib,packages.bib}

\end{document}
